% !TeX program = pdflatex
% !TeX encoding = UTF-8
\documentclass[11pt,aspectratio=169,t]{beamer}
% Все презентации курса компилируются только pdflatex.
\usepackage{iftex}
\RequirePDFTeX
\usepackage{cmap}
\usepackage[T2A]{fontenc}
\usepackage[utf8]{inputenc}
\usepackage[english,russian]{babel}
\usepackage{PTSerif}
\usepackage{amsmath,amsfonts,amssymb}
\usepackage{graphicx,microtype,anyfontsize,tikz}
\usepackage{ragged2e}
\usetikzlibrary{arrows.meta,calc}

\usetheme{default}
\usefonttheme{serif}
\DeclareMathOperator{\sen}{sen}
\bibliographystyle{apalike}

% Общие сведения; номер лекции, её название и дата задаются в N.tex.
\newcommand{\coursename}{Введение в маневрирование КА}
\author{Чиняев Владимир Николаевич}
\institute{Московский физико-технический институт}

% Книжное оформление с голубым акцентом.
\definecolor{coursepaper}{HTML}{FBF9F3}
\definecolor{courseink}{HTML}{302B28}
\definecolor{courseaccent}{HTML}{187AA5}
\definecolor{coursemuted}{HTML}{6D6560}
\definecolor{courserule}{HTML}{A9CADA}
\definecolor{coursetint}{HTML}{EDF4F6}
\tikzset{
  course diagram/.style={font=\fontsize{9.5}{11.5}\selectfont,
    line cap=round,line join=round,>=Latex},
  orbit line/.style={draw=courseaccent,line width=1pt},
  aux line/.style={draw=courserule,line width=.5pt},
  vector line/.style={draw=courseink,line width=.65pt,->},
  angle line/.style={draw=courseaccent,line width=.8pt,->},
  orbit point/.style={circle,fill=courseaccent,inner sep=1.7pt}
}
\setbeamercolor{normal text}{fg=courseink,bg=coursepaper}
\setbeamercolor{background canvas}{bg=coursepaper}
\setbeamercolor{structure}{fg=courseaccent}
\setbeamercolor{alerted text}{fg=courseaccent}
\setbeamercolor{frametitle}{fg=courseink,bg=coursepaper}
\setbeamercolor{footline}{fg=coursemuted,bg=coursepaper}
\setbeamercolor{caption name}{fg=courseaccent}

% Основной текст — 11/14 pt; подписи к рисункам — 9.5/11.5 pt.
\let\coursenormalsize\normalsize
\renewcommand{\normalsize}{\coursenormalsize\fontsize{11}{14}\selectfont}
\setbeamerfont{normal text}{size*={11pt}{14pt}}
\setbeamerfont{frametitle}{size*={20pt}{23pt},series=\mdseries}
\setbeamerfont{title}{size*={27pt}{31pt},series=\mdseries}
\setbeamerfont{subtitle}{size*={12.8pt}{15pt},series=\mdseries}
\setbeamerfont{author}{size*={10.8pt}{13pt}}
\setbeamerfont{institute}{size*={9.2pt}{11pt}}
\setbeamerfont{date}{size*={9pt}{11pt}}
\setbeamerfont{footline}{size*={6.8pt}{8pt}}
\setbeamerfont{caption}{size*={9.5pt}{11.5pt}}
\setbeamerfont{itemize/enumerate body}{size*={11pt}{14pt}}
\setbeamerfont{itemize/enumerate subbody}{size*={11pt}{14pt}}
\setbeamerfont{itemize/enumerate subsubbody}{size*={11pt}{14pt}}
\AtBeginDocument{\normalsize}

\setbeamersize{text margin left=6mm,text margin right=6mm}
\setbeamertemplate{navigation symbols}{}
\setbeamertemplate{headline}{}
\setbeamertemplate{caption}[numbered]
\setbeamertemplate{itemize items}[circle]

% Стандартные блоки Beamer позволяют группировать текст и формулы.
\setbeamertemplate{blocks}[default]
\setbeamerfont{block title}{size*={11pt}{14pt},series=\bfseries}
\setbeamerfont{block body}{size*={11pt}{14pt}}
\setbeamercolor{block title}{fg=courseaccent,bg=coursepaper}
\setbeamercolor{block body}{fg=courseink,bg=coursetint}
\setbeamercolor{block title alerted}{parent=block title}
\setbeamercolor{block body alerted}{parent=block body}
\setbeamercolor{block title example}{parent=block title}
\setbeamercolor{block body example}{parent=block body}

% Компактный заголовок; его высота учитывается при размещении содержания.
\setbeamertemplate{frametitle}{%
  \vskip4.5mm%
  \begin{beamercolorbox}[wd=\textwidth,sep=0pt]{frametitle}
    \usebeamerfont{frametitle}\strut\insertframetitle\strut\par
    \vskip2mm%
    {\color{courseaccent}\hrule height .55pt}
  \end{beamercolorbox}%
}

% Нижний колонтитул: курс слева, лекция и номер слайда справа.
\setbeamertemplate{footline}{%
  \begin{beamercolorbox}[wd=\paperwidth,ht=5.5mm,dp=0pt]{footline}
    \begin{tikzpicture}[remember picture,overlay,x=1mm,y=1mm]
      \begin{scope}[shift={(current page.south west)}]
        \draw[courserule,line width=.3pt] (6,5.5)--(154,5.5);
        \node[text=coursemuted,anchor=west,inner sep=0pt,
          font=\usebeamerfont{footline}] at (6,2.8) {\coursename};
        \node[text=coursemuted,anchor=east,inner sep=0pt,
          font=\usebeamerfont{footline}] at (154,2.8)
          {Лекция \arabic{lecture}\quad\textbullet\quad\insertframenumber};
      \end{scope}
    \end{tikzpicture}%
  \end{beamercolorbox}%
}

% Титульник рассчитан на 16:9 (160 × 90 мм), без нижнего колонтитула.
\setbeamertemplate{title page}{%
  \begin{tikzpicture}[remember picture,overlay,x=1mm,y=1mm]
    \begin{scope}[shift={(current page.south west)}]
      \node[text=coursemuted,anchor=north,font=\usebeamerfont{institute}]
        at (80,82) {\insertinstitute};
      \draw[courseaccent,line width=.55pt] (71,71)--(89,71);
      \node[text=courseink,align=center,text width=125mm,
        font=\usebeamerfont{title}] at (80,53)
        {\hyphenpenalty=10000\coursename\par};
      \node[text=courseaccent,font=\usebeamerfont{subtitle}]
        at (80,32) {Лекция \arabic{lecture}. \inserttitle};
      \node[text=courseink,font=\usebeamerfont{author}]
        at (80,19) {\insertauthor};
      \node[text=coursemuted,font=\usebeamerfont{date}]
        at (80,9) {\insertdate};
    \end{scope}
  \end{tikzpicture}%
}

\usepackage{booktabs,array,tabularx,pgfplots}
\pgfplotsset{compat=1.18,
  course plot/.style={
    font=\fontsize{9.5}{11.5}\selectfont,
    axis background/.style={fill=coursepaper},
    axis line style={draw=coursemuted,line width=.45pt},
    tick style={draw=coursemuted,line width=.4pt},
    grid=major,grid style={draw=courserule!65,line width=.3pt},
    tick align=outside,axis lines=left,
    scaled ticks=false,
    tick label style={/pgf/number format/use comma,/pgf/number format/1000 sep={\,}},
    legend style={draw=none,fill=coursepaper,font=\fontsize{9}{11}\selectfont},
    legend cell align=left,
    label style={font=\fontsize{9.5}{11.5}\selectfont},
    clip=true}}
\setlength{\abovedisplayskip}{5pt}
\setlength{\belowdisplayskip}{5pt}
\setlength{\abovedisplayshortskip}{2pt}
\setlength{\belowdisplayshortskip}{4pt}
\newcommand{\figurelegend}{%
 {\usebeamerfont{caption}\color{coursemuted}
 \textcolor{courseink}{\rule[.6ex]{5mm}{.5pt}} исходная\quad
 \textcolor{courseaccent}{\rule[.6ex]{5mm}{.6pt}} A\quad
 \textcolor{coursemuted}{\textendash\,\textendash} B\par}}

\setcounter{lecture}{3}
\title{Одноимпульсные манёвры}
\date{22 сентября 2026 г.}

\begin{document}

% 0–15 мин: модель импульса и геометрия пересечения.
% Слайд 1
\begin{frame}[plain]
\titlepage
\end{frame}

% Слайд 2
\begin{frame}{Импульсная модель манёвра}
\begin{columns}[T,onlytextwidth]
\begin{column}{.49\textwidth}
\begin{block}{Состояние в точке манёвра}
\centering
$\mathbf r^+=\mathbf r^-=\mathbf r$,\par\smallskip
$\mathbf v^+=\mathbf v^-+\Delta\mathbf v$.
\end{block}
\justifying
Продолжительностью работы ДУ пренебрегаем. До и после импульса КА движется в центральном поле с параметром $\mu$.
\par\smallskip
Компоненты в орбитальном базисе:
\[
\Delta\mathbf v=\Delta v_r\mathbf e_r+
\Delta v_t\mathbf e_t+\Delta v_z\mathbf e_z.
\]
\end{column}
\begin{column}{.49\textwidth}
\centering% Velocity vectors meet at one fixed position. Local radial/transverse basis.
\begin{tikzpicture}[course diagram,x=1mm,y=1mm]
\path[use as bounding box] (-6,-12) rectangle (63,39);
\coordinate (F) at (3,0);
\coordinate (P) at (31,9);
\filldraw[fill=coursetint,draw=courserule] (F) circle[radius=4];
\node at (F) {$\oplus$};
\draw[vector line] (F)--node[below=1mm] {$\mathbf r$}(P);
\draw[aux line,dashed] (17,-10)..controls(27,-10)and(29.5,-2)..(P)
 ..controls(32.5,20)and(33,28)..(30,35);
\draw[aux line,dashed] (10,-9)..controls(18,-4)and(25,2)..(P)
 ..controls(37,16)and(40,21)..(41,27);
\draw[vector line] (P)--(55,17) node[below right] {$\mathbf e_r$};
\draw[vector line] (P)--(24,30) node[above left] {$\mathbf e_t$};
\draw[courseink,line width=1pt,->] (P)--(34,31) node[above] {$\mathbf v^-$};
\draw[courseaccent,line width=1.1pt,->] (P)--(49,30) node[above] {$\mathbf v^+$};
\draw[courseaccent,line width=.9pt,->] (34,31)--(49,30)
 node[midway,above=2mm] {$\Delta\mathbf v$};
\fill[courseink] (P) circle[radius=.9];
\node[anchor=north west,inner sep=1mm] at (P) {$P$};
\end{tikzpicture}
\par
{\usebeamerfont{caption}\color{coursemuted}
Обе орбиты проходят через точку $P$.\par}
\end{column}
\end{columns}
\end{frame}

% Слайд 3
\begin{frame}{Изменение энергии и орбитального момента}
\begin{columns}[T,onlytextwidth]
\begin{column}{.48\textwidth}
\begin{block}{Удельная энергия}
\centering
$\displaystyle\varepsilon=\frac{v^2}{2}-\frac\mu r=-\frac\mu{2a}.$
\par\medskip
$\displaystyle\varepsilon^+-\varepsilon^-
=\mathbf v^-\!\cdot\Delta\mathbf v+
\frac{\|\Delta\mathbf v\|^2}{2}.$
\end{block}
\justifying
Изменение большой полуоси связано с изменением энергии. Вклад зависит от направления импульса относительно скорости.
\end{column}
\begin{column}{.48\textwidth}
\begin{block}{Удельный орбитальный момент}
\centering
$\mathbf h=\mathbf r\times\mathbf v$,\par\smallskip
$\mathbf h^+=\mathbf h^-+\mathbf r\times\Delta\mathbf v$.
\par\smallskip
$h^2=\mu p=\mu a(1-e^2)$.
\end{block}
\justifying
Радиальный импульс сохраняет $\mathbf h$. Нормальная компонента меняет ориентацию плоскости.
\end{column}
\end{columns}
\medskip
\justifying
На эллипсе скорость обычно имеет обе компоненты: $v_r$ и $v_t$.
Трансверсальное направление совпадает с направлением скорости только в апсидах.
\end{frame}

% Слайд 4
\begin{frame}{Общая точка начальной и конечной орбит}
\centering% All conics share their focus, shown by the Earth symbol.
\begin{tikzpicture}[course diagram,x=1mm,y=1mm,scale=.78]
\path[use as bounding box] (-19,-23) rectangle (132,25);
\foreach \shift/\title in {0/{Нет общих точек},38/{Касание},76/{Две точки},114/{Совпадение}} {
 \node[align=center,text width=29mm,text=courseaccent] at (\shift,23) {\title};
 \node[inner sep=0pt] at (\shift,0) {$\oplus$};
}
\draw[courseink,line width=.75pt] (0,0) circle[radius=8];
\draw[orbit line] (0,0) circle[radius=14];
\begin{scope}[xshift=38mm]
 \draw[courseink,line width=.75pt] (0,0) circle[radius=8];
 \draw[orbit line] (-5,0) ellipse[x radius=13,y radius=12];
 \fill[courseaccent] (8,0) circle[radius=.8];
\end{scope}
\begin{scope}[xshift=76mm]
 \draw[courseink,line width=.75pt] (0,0) circle[radius=11];
 \draw[orbit line] (-5.2,0) ellipse[x radius=13,y radius=11.9147];
 \pgfmathsetmacro{\xx}{(10.92/11-1)/.4*11}
 \pgfmathsetmacro{\yy}{sqrt(121-\xx*\xx)}
 \fill[courseaccent] (\xx,\yy) circle[radius=.8];
 \fill[courseaccent] (\xx,-\yy) circle[radius=.8];
\end{scope}
\draw[courseink,line width=1.25pt] (114,0) circle[radius=12];
\draw[courseaccent,line width=.8pt,dashed] (114,0) circle[radius=12];
\foreach \shift/\label in {0/0,38/1,76/2,114/\infty}
 \node at (\shift,-21) {$\label$};
\end{tikzpicture}
\par
\begin{block}{Условие одноимпульсного перехода}
\justifying
Орбиты должны иметь общую точку. В ней требуемый импульс равен
$\Delta\mathbf v=\mathbf v_2-\mathbf v_1$.
\end{block}
\end{frame}

% Слайд 5
\begin{frame}{Радиус-векторы компланарных орбит}
\justifying
Выберем неподвижные орты $\mathbf i,\mathbf j$ в общей плоскости орбит.
Углы отсчитываем от $\mathbf i$ в одном направлении; $p_j=a_j(1-e_j^2)$.
\begin{align*}
\mathbf r_1&=\frac{p_1}{1+e_1\cos\nu_1}
 \bigl[\cos(\nu_1+\omega_1)\mathbf i+\sin(\nu_1+\omega_1)\mathbf j\bigr],\\[3pt]
\mathbf r_2&=\frac{p_2}{1+e_2\cos\nu_2}
 \bigl[\cos(\nu_2+\omega_2)\mathbf i+\sin(\nu_2+\omega_2)\mathbf j\bigr].
\end{align*}
\justifying
В общей точке $\mathbf r_1=\mathbf r_2$: совпадают направление и длина радиус-вектора.
Поскольку $r_1,r_2>0$, получаем два скалярных условия:
\[
\omega_1+\nu_1=\omega_2+\nu_2,\qquad
\frac{p_1}{1+e_1\cos\nu_1}=\frac{p_2}{1+e_2\cos\nu_2}.
\]
\end{frame}

% Слайд 6
\begin{frame}{Уравнение для точки пересечения}
Из равенства направлений: $\Delta\omega=\omega_2-\omega_1$,
$\nu_2=\nu_1-\Delta\omega$. Подставим во второе условие:
\[
p_1\bigl[1+e_2\cos(\nu_1-\Delta\omega)\bigr]
=p_2(1+e_1\cos\nu_1).
\]
Раскроем косинус разности:
\[
\cos(\nu_1-\Delta\omega)
=\cos\nu_1\cos\Delta\omega+\sin\nu_1\sin\Delta\omega.
\]
Соберём коэффициенты при $\cos\nu_1$, $\sin\nu_1$ и свободный член:
\[
\underbrace{(e_2p_1\cos\Delta\omega-e_1p_2)}_{A}\cos\nu_1+
\underbrace{e_2p_1\sin\Delta\omega}_{B}\sin\nu_1+
\underbrace{(p_1-p_2)}_{C}=0.
\]
\justifying
Теперь неизвестна только $\nu_1$. Для каждого корня найдём
$\nu_2=\nu_1-\Delta\omega$, общий радиус и обе скорости.
\end{frame}

% Слайд 7
\begin{frame}{Решение тригонометрического уравнения}
Пусть $A\cos\nu_1+B\sin\nu_1=-C$ и $\rho=\sqrt{A^2+B^2}>0$.
Введём угол $\alpha$:
\[
\cos\alpha=\frac A\rho,\qquad \sin\alpha=\frac B\rho,
\qquad \alpha=\operatorname{atan2}(B,A).
\]
Тогда
\[
A\cos\nu_1+B\sin\nu_1
=\rho\bigl(\cos\alpha\cos\nu_1+\sin\alpha\sin\nu_1\bigr)
=\rho\cos(\nu_1-\alpha).
\]
\[
\cos(\nu_1-\alpha)=-\frac C\rho
\quad\Longrightarrow\quad
\boxed{\nu_1=\alpha\pm\arccos(-C/\rho)\pmod{2\pi}.}
\]
\justifying
Функция $\operatorname{atan2}(y,x)$ определяет угол по обоим знакам $x,y$.
Обычная $\arctan(B/A)$ может выбрать неверную четверть.
\end{frame}

% Слайд 8
\begin{frame}{Число точек пересечения}
При $\rho>0$ число решений следует из диапазона значений косинуса:
\[
\cos(\nu_1-\alpha)=-C/\rho.
\]
\renewcommand{\arraystretch}{1.2}
\begin{tabularx}{.99\textwidth}{@{}lX@{}}
$|C|>\rho$ & Правая часть вне $[-1,1]$: общих точек нет.\\
$|C|=\rho$ & Корни со знаками $+$ и $-$ совпадают по модулю $2\pi$: касание.\\
$|C|<\rho$ & Два различных корня: две точки пересечения.\\
\end{tabularx}
\begin{block}{Отдельно проверяем $A=B=0$}
\justifying
Уравнение сводится к $C=0$. Если $C\ne0$, пересечений нет.
Если $C=0$, орбиты совпадают и подходит любая точка.
\end{block}
\justifying
Например, для кругов $e_1=e_2=0$: $A=B=0$, $C=a_1-a_2$.
Круги совпадают при $a_1=a_2$ и не пересекаются при $a_1\ne a_2$.
\end{frame}

% 15–25 мин: компланарный пример.
% Слайд 9
\begin{frame}{Компланарный переход: исходные данные}
\begin{columns}[T,onlytextwidth]
\begin{column}{.47\textwidth}
\begin{block}{Две орбиты в общей плоскости}
\centering
\renewcommand{\arraystretch}{1.18}
\begin{tabular}{@{}lrr@{}}
 & Начальная & Конечная\\\midrule
$a$, км & $10\,000$ & $12\,000$\\
$e$ & $0{,}10$ & $0{,}35$\\
$\omega$, град & $20$ & $90$
\end{tabular}
\end{block}
\justifying
Углы $\omega$ здесь отсчитываются от общей оси в плоскости орбит.
\par\medskip
\textbf{Найти:} точки пересечения и импульс перехода в каждой из них.
\end{column}
\begin{column}{.51\textwidth}
\centering\begin{tikzpicture}[course diagram,x=1mm,y=1mm]
\path[use as bounding box] (-34.0, -24.5) rectangle (34.0, 24.5);
\filldraw[fill=coursetint,draw=courserule,line width=.4pt] (0.0000,5.4859) circle[radius=10.6247];
\node[inner sep=0pt,font=\fontsize{8.5}{10}\selectfont] at (0.0000,5.4859) {$\oplus$};
\draw[draw=courseink,line width=.75pt] plot coordinates {(14.08801,10.61348) (13.94939,10.98068) (13.80207,11.34450) (13.64610,11.70475) (13.48157,12.06127) (13.30855,12.41386) (13.12712,12.76235) (12.93737,13.10656) (12.73937,13.44631) (12.53322,13.78143) (12.31901,14.11173) (12.09683,14.43705) (11.86677,14.75721) (11.62895,15.07203) (11.38345,15.38135) (11.13039,15.68499) (10.86988,15.98279) (10.60203,16.27457) (10.32695,16.56016) (10.04476,16.83940) (9.75559,17.11212) (9.45956,17.37816) (9.15680,17.63735) (8.84744,17.88952) (8.53163,18.13453) (8.20950,18.37220) (7.88120,18.60238) (7.54687,18.82491) (7.20667,19.03964) (6.86075,19.24640) (6.50928,19.44505) (6.15241,19.63545) (5.79033,19.81742) (5.42319,19.99085) (5.05118,20.15556) (4.67448,20.31144) (4.29328,20.45832) (3.90777,20.59609) (3.51814,20.72460) (3.12459,20.84372) (2.72733,20.95332) (2.32656,21.05327) (1.92251,21.14347) (1.51539,21.22377) (1.10542,21.29408) (0.69283,21.35428) (0.27786,21.40426) (-0.13926,21.44391) (-0.55829,21.47315) (-0.97897,21.49188) (-1.40106,21.50000) (-1.82429,21.49744) (-2.24841,21.48412) (-2.67315,21.45996) (-3.09823,21.42490) (-3.52339,21.37887) (-3.94835,21.32184) (-4.37282,21.25373) (-4.79651,21.17453) (-5.21912,21.08419) (-5.64038,20.98268) (-6.05997,20.86999) (-6.47759,20.74611) (-6.89294,20.61104) (-7.30570,20.46478) (-7.71557,20.30734) (-8.12223,20.13876) (-8.52536,19.95906) (-8.92464,19.76828) (-9.31976,19.56649) (-9.71038,19.35373) (-10.09619,19.13008) (-10.47685,18.89562) (-10.85205,18.65045) (-11.22145,18.39467) (-11.58473,18.12838) (-11.94157,17.85173) (-12.29164,17.56483) (-12.63462,17.26784) (-12.97018,16.96092) (-13.29802,16.64423) (-13.61780,16.31796) (-13.92923,15.98229) (-14.23198,15.63744) (-14.52576,15.28362) (-14.81027,14.92105) (-15.08520,14.54997) (-15.35027,14.17063) (-15.60520,13.78330) (-15.84971,13.38824) (-16.08353,12.98574) (-16.30640,12.57609) (-16.51806,12.15960) (-16.71827,11.73657) (-16.90680,11.30734) (-17.08341,10.87225) (-17.24790,10.43162) (-17.40006,9.98583) (-17.53969,9.53522) (-17.66662,9.08018) (-17.78067,8.62108) (-17.88169,8.15830) (-17.96954,7.69225) (-18.04407,7.22331) (-18.10518,6.75191) (-18.15276,6.27843) (-18.18673,5.80332) (-18.20700,5.32698) (-18.21353,4.84984) (-18.20626,4.37233) (-18.18516,3.89487) (-18.15023,3.41791) (-18.10146,2.94187) (-18.03888,2.46719) (-17.96251,1.99431) (-17.87240,1.52365) (-17.76863,1.05565) (-17.65127,0.59074) (-17.52041,0.12934) (-17.37617,-0.32812) (-17.21868,-0.78122) (-17.04808,-1.22955) (-16.86452,-1.67270) (-16.66818,-2.11026) (-16.45925,-2.54184) (-16.23792,-2.96706) (-16.00441,-3.38552) (-15.75895,-3.79686) (-15.50177,-4.20071) (-15.23314,-4.59673) (-14.95331,-4.98456) (-14.66257,-5.36387) (-14.36120,-5.73433) (-14.04951,-6.09565) (-13.72779,-6.44752) (-13.39637,-6.78965) (-13.05558,-7.12176) (-12.70575,-7.44360) (-12.34723,-7.75491) (-11.98036,-8.05547) (-11.60551,-8.34504) (-11.22304,-8.62342) (-10.83331,-8.89042) (-10.43670,-9.14585) (-10.03359,-9.38954) (-9.62435,-9.62134) (-9.20938,-9.84112) (-8.78906,-10.04875) (-8.36378,-10.24411) (-7.93392,-10.42711) (-7.49987,-10.59766) (-7.06203,-10.75570) (-6.62079,-10.90116) (-6.17653,-11.03401) (-5.72964,-11.15422) (-5.28051,-11.26176) (-4.82951,-11.35663) (-4.37703,-11.43883) (-3.92344,-11.50840) (-3.46911,-11.56535) (-3.01442,-11.60974) (-2.55972,-11.64160) (-2.10537,-11.66102) (-1.65174,-11.66807) (-1.19915,-11.66282) (-0.74797,-11.64538) (-0.29851,-11.61586) (0.14888,-11.57435) (0.59389,-11.52099) (1.03620,-11.45591) (1.47551,-11.37924) (1.91150,-11.29114) (2.34389,-11.19175) (2.77238,-11.08124) (3.19671,-10.95977) (3.61659,-10.82752) (4.03177,-10.68466) (4.44198,-10.53138) (4.84698,-10.36787) (5.24652,-10.19433) (5.64038,-10.01095) (6.02832,-9.81793) (6.41014,-9.61548) (6.78562,-9.40381) (7.15455,-9.18314) (7.51675,-8.95368) (7.87203,-8.71564) (8.22020,-8.46926) (8.56110,-8.21475) (8.89455,-7.95234) (9.22042,-7.68225) (9.53853,-7.40472) (9.84875,-7.11996) (10.15095,-6.82822) (10.44499,-6.52972) (10.73076,-6.22470) (11.00812,-5.91338) (11.27699,-5.59600) (11.53724,-5.27278) (11.78879,-4.94397) (12.03154,-4.60979) (12.26541,-4.27048) (12.49032,-3.92626) (12.70619,-3.57737) (12.91296,-3.22403) (13.11055,-2.86648) (13.29893,-2.50494) (13.47802,-2.13963) (13.64779,-1.77079) (13.80819,-1.39864) (13.95919,-1.02341) (14.10074,-0.64531) (14.23282,-0.26457) (14.35541,0.11860) (14.46848,0.50397) (14.57202,0.89133) (14.66601,1.28046) (14.75045,1.67114) (14.82533,2.06317) (14.89066,2.45633) (14.94642,2.85041) (14.99264,3.24520) (15.02932,3.64050) (15.05647,4.03609) (15.07410,4.43178) (15.08225,4.82736) (15.08093,5.22263) (15.07016,5.61738) (15.04998,6.01142) (15.02042,6.40456) (14.98151,6.79658) (14.93329,7.18730) (14.87580,7.57652) (14.80908,7.96406) (14.73319,8.34971) (14.64817,8.73329) (14.55407,9.11460) (14.45094,9.49347) (14.33885,9.86969) (14.21785,10.24309) (14.08801,10.61348)};
\draw[draw=courseaccent,line width=1pt] plot coordinates {(0.00000,18.47906) (-0.34015,18.47576) (-0.68025,18.46587) (-1.02025,18.44937) (-1.36009,18.42626) (-1.69972,18.39654) (-2.03909,18.36019) (-2.37815,18.31720) (-2.71683,18.26754) (-3.05509,18.21122) (-3.39286,18.14819) (-3.73009,18.07844) (-4.06672,18.00195) (-4.40269,17.91868) (-4.73793,17.82860) (-5.07238,17.73169) (-5.40598,17.62790) (-5.73865,17.51719) (-6.07032,17.39954) (-6.40092,17.27490) (-6.73037,17.14321) (-7.05860,17.00445) (-7.38551,16.85855) (-7.71103,16.70548) (-8.03506,16.54518) (-8.35751,16.37759) (-8.67829,16.20266) (-8.99729,16.02034) (-9.31441,15.83057) (-9.62954,15.63329) (-9.94257,15.42844) (-10.25337,15.21595) (-10.56183,14.99578) (-10.86781,14.76786) (-11.17118,14.53211) (-11.47181,14.28849) (-11.76953,14.03693) (-12.06421,13.77737) (-12.35568,13.50974) (-12.64378,13.23398) (-12.92833,12.95004) (-13.20916,12.65786) (-13.48609,12.35737) (-13.75890,12.04853) (-14.02742,11.73128) (-14.29142,11.40557) (-14.55069,11.07135) (-14.80500,10.72859) (-15.05411,10.37725) (-15.29780,10.01728) (-15.53579,9.64867) (-15.76783,9.27139) (-15.99366,8.88542) (-16.21299,8.49077) (-16.42553,8.08742) (-16.63098,7.67538) (-16.82903,7.25467) (-17.01938,6.82532) (-17.20168,6.38737) (-17.37561,5.94086) (-17.54081,5.48587) (-17.69694,5.02246) (-17.84363,4.55072) (-17.98050,4.07077) (-18.10717,3.58273) (-18.22326,3.08673) (-18.32837,2.58294) (-18.42210,2.07153) (-18.50403,1.55272) (-18.57374,1.02671) (-18.63083,0.49375) (-18.67486,-0.04588) (-18.70541,-0.59189) (-18.72204,-1.14395) (-18.72433,-1.70173) (-18.71185,-2.26483) (-18.68417,-2.83286) (-18.64087,-3.40537) (-18.58152,-3.98189) (-18.50573,-4.56192) (-18.41308,-5.14493) (-18.30319,-5.73034) (-18.17568,-6.31756) (-18.03019,-6.90594) (-17.86637,-7.49481) (-17.68391,-8.08348) (-17.48251,-8.67119) (-17.26188,-9.25717) (-17.02180,-9.84063) (-16.76204,-10.42071) (-16.48242,-10.99655) (-16.18281,-11.56725) (-15.86310,-12.13189) (-15.52323,-12.68950) (-15.16320,-13.23912) (-14.78303,-13.77975) (-14.38281,-14.31037) (-13.96268,-14.82996) (-13.52283,-15.33747) (-13.06352,-15.83186) (-12.58505,-16.31207) (-12.08779,-16.77707) (-11.57217,-17.22580) (-11.03869,-17.65723) (-10.48790,-18.07035) (-9.92043,-18.46416) (-9.33694,-18.83769) (-8.73819,-19.19001) (-8.12496,-19.52020) (-7.49814,-19.82742) (-6.85862,-20.11084) (-6.20737,-20.36971) (-5.54543,-20.60331) (-4.87384,-20.81101) (-4.19372,-20.99223) (-3.50621,-21.14645) (-2.81250,-21.27324) (-2.11378,-21.37223) (-1.41129,-21.44314) (-0.70628,-21.48578) (-0.00000,-21.50000) (0.70628,-21.48578) (1.41129,-21.44314) (2.11378,-21.37223) (2.81250,-21.27324) (3.50621,-21.14645) (4.19372,-20.99223) (4.87384,-20.81101) (5.54543,-20.60331) (6.20737,-20.36971) (6.85862,-20.11084) (7.49814,-19.82742) (8.12496,-19.52020) (8.73819,-19.19001) (9.33694,-18.83769) (9.92043,-18.46416) (10.48790,-18.07035) (11.03869,-17.65723) (11.57217,-17.22580) (12.08779,-16.77707) (12.58505,-16.31207) (13.06352,-15.83186) (13.52283,-15.33747) (13.96268,-14.82996) (14.38281,-14.31037) (14.78303,-13.77975) (15.16320,-13.23912) (15.52323,-12.68950) (15.86310,-12.13189) (16.18281,-11.56725) (16.48242,-10.99655) (16.76204,-10.42071) (17.02180,-9.84063) (17.26188,-9.25717) (17.48251,-8.67119) (17.68391,-8.08348) (17.86637,-7.49481) (18.03019,-6.90594) (18.17568,-6.31756) (18.30319,-5.73034) (18.41308,-5.14493) (18.50573,-4.56192) (18.58152,-3.98189) (18.64087,-3.40537) (18.68417,-2.83286) (18.71185,-2.26483) (18.72433,-1.70173) (18.72204,-1.14395) (18.70541,-0.59189) (18.67486,-0.04588) (18.63083,0.49375) (18.57374,1.02671) (18.50403,1.55272) (18.42210,2.07153) (18.32837,2.58294) (18.22326,3.08673) (18.10717,3.58273) (17.98050,4.07077) (17.84363,4.55072) (17.69694,5.02246) (17.54081,5.48587) (17.37561,5.94086) (17.20168,6.38737) (17.01938,6.82532) (16.82903,7.25467) (16.63098,7.67538) (16.42553,8.08742) (16.21299,8.49077) (15.99366,8.88542) (15.76783,9.27139) (15.53579,9.64867) (15.29780,10.01728) (15.05411,10.37725) (14.80500,10.72859) (14.55069,11.07135) (14.29142,11.40557) (14.02742,11.73128) (13.75890,12.04853) (13.48609,12.35737) (13.20916,12.65786) (12.92833,12.95004) (12.64378,13.23398) (12.35568,13.50974) (12.06421,13.77737) (11.76953,14.03693) (11.47181,14.28849) (11.17118,14.53211) (10.86781,14.76786) (10.56183,14.99578) (10.25337,15.21595) (9.94257,15.42844) (9.62954,15.63329) (9.31441,15.83057) (8.99729,16.02034) (8.67829,16.20266) (8.35751,16.37759) (8.03506,16.54518) (7.71103,16.70548) (7.38551,16.85855) (7.05860,17.00445) (6.73037,17.14321) (6.40092,17.27490) (6.07032,17.39954) (5.73865,17.51719) (5.40598,17.62790) (5.07238,17.73169) (4.73793,17.82860) (4.40269,17.91868) (4.06672,18.00195) (3.73009,18.07844) (3.39286,18.14819) (3.05509,18.21122) (2.71683,18.26754) (2.37815,18.31720) (2.03909,18.36019) (1.69972,18.39654) (1.36009,18.42626) (1.02025,18.44937) (0.68025,18.46587) (0.34015,18.47576) (0.00000,18.47906)};
\fill[courseaccent] (13.1489,12.7216) circle[radius=.8];
\node[anchor=south west,inner sep=1mm] at (13.1489,12.7216) {1};
\fill[courseaccent] (-18.1494,3.4085) circle[radius=.8];
\node[anchor=north east,inner sep=1mm] at (-18.1494,3.4085) {2};
\end{tikzpicture}
\par
{\usebeamerfont{caption}\color{coursemuted}
Тёмная линия --- начальная орбита; голубая --- конечная.\par}
\end{column}
\end{columns}
\end{frame}

% Слайд 10
\begin{frame}{Компланарный переход: точки пересечения}
Подставим элементы в уравнение пересечения:
\[
p_1=9900\text{ км},\qquad p_2=10\,530\text{ км},\qquad
\Delta\omega=70^\circ,
\]
\[
132{,}0998\cos\nu_1+3256{,}0349\sin\nu_1-630=0.
\]
Приводим к одному косинусу:
\[
\cos(\nu_1-87{,}6767^\circ)=0{,}1933278.
\]
\centering
\renewcommand{\arraystretch}{1.05}
\begin{tabular}{@{}lcc@{}}
\toprule
 & Точка 1 & Точка 2\\\midrule
$\nu_1$ & $8{,}8238^\circ$ & $166{,}5297^\circ$\\
$\nu_2=\nu_1-\Delta\omega$ & $298{,}8238^\circ$ & $96{,}5297^\circ$\\\bottomrule
\end{tabular}\par\smallskip
\justifying
Все значения $\nu$ приводим к диапазону $0\le\nu<360^\circ$, добавляя или вычитая полный оборот.
\end{frame}

% Слайд 11
\begin{frame}{Расчёт импульса в общей точке}
\begin{columns}[T,onlytextwidth]
\begin{column}{.58\textwidth}
Компоненты скорости в плоскости орбиты:
\[
v_r=\sqrt{\frac\mu p}\,e\sin\nu,\qquad
v_t=\sqrt{\frac\mu p}\,(1+e\cos\nu).
\]
При $u=\omega+\nu$ переходим к общим осям:
\[
v_x=v_r\cos u-v_t\sin u,
\qquad
v_y=v_r\sin u+v_t\cos u.
\]
\[
\boxed{\Delta\mathbf v=\mathbf v_2-\mathbf v_1},\qquad \Delta v=\|\Delta\mathbf v\|.
\]
\end{column}
\begin{column}{.38\textwidth}
\centering% Example 2.1, first intersection. Components in the common local RT basis.
\begin{tikzpicture}[course diagram,x=1mm,y=1mm,scale=.88]
\path[use as bounding box] (-22,-8) rectangle (32,46);
\coordinate (O) at (0,0);
\coordinate (A) at (.584,41.834);
\coordinate (B) at (-11.320,43.144);
\draw[vector line] (O)--(27,0) node[right] {$v_r$};
\draw[vector line] (O)--(0,46) node[above right,inner sep=.5mm] {$v_t$};
\draw[courseink,line width=1.05pt,->] (O)--(A);
\node[anchor=west,inner sep=1mm] at (3,31) {$\mathbf v_1$};
\draw[courseaccent,line width=1.05pt,->] (O)--(B);
\node[anchor=east,inner sep=1mm] at (-9,30) {$\mathbf v_2$};
\draw[courseaccent,line width=1.05pt,->] (A)--(B);
\node[anchor=south,inner sep=1mm] at (-6,44) {$\Delta\mathbf v$};
\draw[aux line,dashed] (B)--(-11.320,0);
\node[below] at (-11.320,0) {$v_{r,2}<0$};
\fill[courseink] (O) circle[radius=.7];
\end{tikzpicture}

\end{column}
\end{columns}
\smallskip
\justifying
В примерах $\mu=398\,600{,}44158$ км$^3$/с$^2$.
При $p$ в км получаем км/с; векторы скоростей далее приводим в м/с.
\end{frame}

% Слайд 12
\begin{frame}{Компланарный переход: первая точка}
При $\nu_1=8{,}823797^\circ$, $\nu_2=298{,}823797^\circ$:
\[
r=\frac{9900}{1+0{,}1\cos8{,}823797^\circ}
=9009{,}6938\text{ км},\quad u=28{,}823797^\circ.
\]
\begin{align*}
\mathbf r_1=\mathbf r_2&=(7893{,}4514,\;4343{,}7320,\;0)^{\mathsf T}\text{ км},\\
\mathbf v_1&=(-3276{,}1963,\;6155{,}4095,\;0)^{\mathsf T}\text{ м/с},\\
\mathbf v_2&=(-5119{,}6398,\;5390{,}2835,\;0)^{\mathsf T}\text{ м/с}.
\end{align*}
Вычитаем компоненты, затем находим длину вектора:
\[
\Delta\mathbf v=(-1843{,}4435,\;-765{,}1259,\;0)^{\mathsf T}\text{ м/с},
\]
\[
\Delta v=\sqrt{1843{,}4435^2+765{,}1259^2}=1995{,}9213\text{ м/с}.
\]
\end{frame}

% Слайд 13
\begin{frame}{Компланарный переход: вторая точка}
При $\nu_1=166{,}529685^\circ$, $\nu_2=96{,}529685^\circ$:
\[
r=\frac{9900}{1+0{,}1\cos166{,}529685^\circ}
=10\,966{,}4800\text{ км},\quad u=186{,}529685^\circ.
\]
\begin{align*}
\mathbf r_1=\mathbf r_2&=(-10\,895{,}3413,\;-1247{,}0858,\;0)^{\mathsf T}\text{ км},\\
\mathbf v_1&=(504{,}5515,\;-5707{,}8639,\;0)^{\mathsf T}\text{ м/с},\\
\mathbf v_2&=(-1453{,}7355,\;-6112{,}6325,\;0)^{\mathsf T}\text{ м/с}.
\end{align*}
\[
\Delta\mathbf v=(-1958{,}2871,\;-404{,}7686,\;0)^{\mathsf T}\text{ м/с},
\]
\[
\Delta v=\sqrt{1958{,}2871^2+404{,}7686^2}=1999{,}6815\text{ м/с}.
\]
\justifying
Первая точка требует на $3{,}7602$ м/с меньше. Результат сравниваем
после вычисления обоих векторных импульсов.
\end{frame}

% 25–40 мин: пространственный переход.
% Слайд 14
\begin{frame}{Пересечение некомпланарных орбит}
\begin{columns}[T,onlytextwidth]
\begin{column}{.49\textwidth}
\begin{block}{Линия пересечения плоскостей}
\centering
$\displaystyle\mathbf L=
\frac{\widehat{\mathbf h}_1\times\widehat{\mathbf h}_2}
{\|\widehat{\mathbf h}_1\times\widehat{\mathbf h}_2\|},\quad
\widehat{\mathbf h}_j=\frac{\mathbf h_j}{\|\mathbf h_j\|}.$
\end{block}
\justifying
Проверяются оба направления:
$+\mathbf L$ и $-\mathbf L$.
\par\smallskip
Аномалию $\nu_j$ отсчитываем от перицентра до выбранного луча по движению КА.
\par\smallskip
$\displaystyle r_j=\frac{p_j}{1+e_j\cos\nu_j}$.
\end{column}
\begin{column}{.49\textwidth}
\centering% Schematic, orthographic view. Generated by intersection_schematic().
\begin{tikzpicture}[course diagram,x=1mm,y=1mm]
\path[use as bounding box] (-35,-21) rectangle (35,21);
\draw[coursemuted,densely dashed,line width=.55pt] (14.2167,-9.3915)--(-4.4930,17.3288);
\draw[courseink,line width=.85pt] plot coordinates {(13.01481,17.40319) (12.91746,17.46374) (12.81806,17.52110) (12.71663,17.57527) (12.61319,17.62623) (12.50774,17.67398) (12.40029,17.71851) (12.29086,17.75981) (12.17946,17.79787) (12.06610,17.83268) (11.95079,17.86424) (11.83354,17.89253) (11.71435,17.91754) (11.59326,17.93926) (11.47025,17.95769) (11.34535,17.97281) (11.21856,17.98461) (11.08990,17.99308) (10.95937,17.99822) (10.82700,18.00000) (10.69278,17.99842) (10.55672,17.99346) (10.41885,17.98512) (10.27917,17.97338) (10.13770,17.95823) (9.99444,17.93966) (9.84940,17.91765) (9.70260,17.89219) (9.55405,17.86327) (9.40376,17.83087) (9.25174,17.79499) (9.09801,17.75561) (8.94258,17.71271) (8.78545,17.66628) (8.62666,17.61631) (8.46620,17.56278) (8.30409,17.50568) (8.14034,17.44500) (7.97498,17.38072) (7.80800,17.31282) (7.63944,17.24130) (7.46930,17.16614) (7.29760,17.08733) (7.12436,17.00484) (6.94959,16.91867) (6.77332,16.82881) (6.59554,16.73523) (6.41630,16.63793) (6.23560,16.53689) (6.05346,16.43210) (5.86991,16.32354) (5.68496,16.21120) (5.49864,16.09508) (5.31096,15.97514) (5.12195,15.85140) (4.93164,15.72382) (4.74004,15.59241) (4.54718,15.45714) (4.35309,15.31802) (4.15779,15.17502) (3.96132,15.02814) (3.76369,14.87737) (3.56493,14.72269) (3.36508,14.56412) (3.16417,14.40162) (2.96223,14.23521) (2.75929,14.06486) (2.55539,13.89059) (2.35055,13.71237) (2.14483,13.53022) (1.93824,13.34412) (1.73084,13.15408) (1.52266,12.96010) (1.31374,12.76216) (1.10412,12.56029) (0.89385,12.35448) (0.68298,12.14473) (0.47154,11.93106) (0.25959,11.71346) (0.04717,11.49195) (-0.16567,11.26653) (-0.37887,11.03722) (-0.59239,10.80404) (-0.80616,10.56699) (-1.02013,10.32610) (-1.23424,10.08137) (-1.44844,9.83284) (-1.66266,9.58053) (-1.87685,9.32446) (-2.09093,9.06466) (-2.30484,8.80115) (-2.51852,8.53397) (-2.73189,8.26316) (-2.94488,7.98874) (-3.15743,7.71076) (-3.36946,7.42927) (-3.58088,7.14429) (-3.79163,6.85589) (-4.00163,6.56411) (-4.21079,6.26900) (-4.41902,5.97063) (-4.62626,5.66904) (-4.83240,5.36431) (-5.03737,5.05649) (-5.24107,4.74565) (-5.44341,4.43188) (-5.64429,4.11523) (-5.84363,3.79580) (-6.04133,3.47367) (-6.23729,3.14891) (-6.43140,2.82162) (-6.62358,2.49190) (-6.81371,2.15984) (-7.00169,1.82554) (-7.18743,1.48910) (-7.37081,1.15065) (-7.55172,0.81028) (-7.73006,0.46811) (-7.90572,0.12428) (-8.07858,-0.22111) (-8.24854,-0.56792) (-8.41549,-0.91602) (-8.57930,-1.26527) (-8.73988,-1.61553) (-8.89709,-1.96666) (-9.05084,-2.31851) (-9.20100,-2.67092) (-9.34747,-3.02375) (-9.49012,-3.37682) (-9.62885,-3.72999) (-9.76353,-4.08307) (-9.89406,-4.43591) (-10.02032,-4.78832) (-10.14220,-5.14014) (-10.25959,-5.49117) (-10.37238,-5.84123) (-10.48046,-6.19015) (-10.58372,-6.53772) (-10.68206,-6.88376) (-10.77537,-7.22807) (-10.86355,-7.57046) (-10.94650,-7.91073) (-11.02412,-8.24867) (-11.09631,-8.58408) (-11.16300,-8.91677) (-11.22407,-9.24651) (-11.27945,-9.57312) (-11.32906,-9.89637) (-11.37281,-10.21607) (-11.41063,-10.53200) (-11.44244,-10.84396) (-11.46818,-11.15174) (-11.48778,-11.45512) (-11.50118,-11.75391) (-11.50832,-12.04791) (-11.50916,-12.33689) (-11.50364,-12.62067) (-11.49173,-12.89905) (-11.47339,-13.17182) (-11.44858,-13.43879) (-11.41729,-13.69977) (-11.37948,-13.95457) (-11.33515,-14.20300) (-11.28429,-14.44489) (-11.22689,-14.68005) (-11.16295,-14.90832) (-11.09249,-15.12952) (-11.01551,-15.34349) (-10.93204,-15.55008) (-10.84210,-15.74914) (-10.74572,-15.94051) (-10.64295,-16.12405) (-10.53382,-16.29965) (-10.41838,-16.46716) (-10.29670,-16.62646) (-10.16882,-16.77746) (-10.03483,-16.92003) (-9.89479,-17.05408) (-9.74877,-17.17952) (-9.59688,-17.29627) (-9.43919,-17.40425) (-9.27580,-17.50340) (-9.10681,-17.59365) (-8.93233,-17.67495) (-8.75246,-17.74727) (-8.56733,-17.81057) (-8.37704,-17.86482) (-8.18174,-17.91000) (-7.98153,-17.94611) (-7.77657,-17.97314) (-7.56697,-17.99109) (-7.35289,-18.00000) (-7.13446,-17.99987) (-6.91183,-17.99075) (-6.68515,-17.97266) (-6.45457,-17.94566) (-6.22025,-17.90980) (-5.98234,-17.86514) (-5.74100,-17.81175) (-5.49638,-17.74971) (-5.24866,-17.67910) (-4.99799,-17.60001) (-4.74454,-17.51254) (-4.48847,-17.41679) (-4.22995,-17.31288) (-3.96913,-17.20090) (-3.70619,-17.08100) (-3.44130,-16.95328) (-3.17460,-16.81788) (-2.90628,-16.67495) (-2.63649,-16.52461) (-2.36540,-16.36701) (-2.09316,-16.20230) (-1.81993,-16.03064) (-1.54588,-15.85217) (-1.27116,-15.66706) (-0.99593,-15.47546) (-0.72032,-15.27755) (-0.44451,-15.07349) (-0.16863,-14.86344) (0.10717,-14.64759) (0.38275,-14.42609) (0.65797,-14.19912) (0.93269,-13.96686) (1.20677,-13.72949) (1.48009,-13.48717) (1.75252,-13.24008) (2.02393,-12.98841) (2.29420,-12.73232) (2.56321,-12.47200) (2.83085,-12.20761) (3.09700,-11.93933) (3.36156,-11.66734) (3.62442,-11.39181) (3.88548,-11.11291) (4.14464,-10.83080) (4.40180,-10.54567) (4.65688,-10.25767) (4.90978,-9.96697) (5.16043,-9.67374) (5.40874,-9.37812) (5.65463,-9.08030) (5.89803,-8.78041) (6.13886,-8.47862) (6.37706,-8.17507) (6.61257,-7.86992) (6.84532,-7.56331) (7.07525,-7.25539) (7.30231,-6.94630) (7.52645,-6.63617) (7.74762,-6.32515) (7.96576,-6.01336) (8.18084,-5.70094) (8.39282,-5.38802) (8.60165,-5.07471) (8.80731,-4.76115) (9.00976,-4.44744) (9.20896,-4.13371) (9.40489,-3.82006) (9.59754,-3.50662) (9.78686,-3.19347) (9.97284,-2.88074) (10.15547,-2.56852) (10.33472,-2.25690) (10.51058,-1.94600) (10.68303,-1.63588) (10.85208,-1.32666) (11.01770,-1.01841) (11.17989,-0.71122) (11.33865,-0.40518) (11.49397,-0.10035) (11.64584,0.20317) (11.79428,0.50532) (11.93927,0.80603) (12.08082,1.10523) (12.21893,1.40284) (12.35361,1.69880) (12.48486,1.99306) (12.61269,2.28555) (12.73710,2.57621) (12.85811,2.86499) (12.97573,3.15183) (13.08996,3.43668) (13.20082,3.71950) (13.30832,4.00023) (13.41247,4.27882) (13.51328,4.55524) (13.61078,4.82944) (13.70497,5.10138) (13.79587,5.37102) (13.88350,5.63833) (13.96787,5.90326) (14.04901,6.16579) (14.12692,6.42588) (14.20163,6.68350) (14.27315,6.93862) (14.34151,7.19121) (14.40671,7.44125) (14.46879,7.68870) (14.52776,7.93354) (14.58364,8.17576) (14.63646,8.41532) (14.68622,8.65220) (14.73295,8.88639) (14.77667,9.11786) (14.81740,9.34659) (14.85517,9.57258) (14.88998,9.79579) (14.92187,10.01622) (14.95085,10.23385) (14.97694,10.44866) (15.00016,10.66065) (15.02054,10.86979) (15.03810,11.07607) (15.05284,11.27949) (15.06480,11.48003) (15.07400,11.67768) (15.08045,11.87243) (15.08417,12.06427) (15.08519,12.25319) (15.08352,12.43919) (15.07918,12.62224) (15.07220,12.80236) (15.06258,12.97951) (15.05036,13.15371) (15.03554,13.32494) (15.01815,13.49320) (14.99821,13.65848) (14.97573,13.82077) (14.95073,13.98006) (14.92323,14.13636) (14.89325,14.28965) (14.86080,14.43993) (14.82590,14.58719) (14.78857,14.73143) (14.74883,14.87264) (14.70669,15.01081) (14.66216,15.14595) (14.61527,15.27805) (14.56602,15.40710) (14.51445,15.53309) (14.46055,15.65602) (14.40435,15.77589) (14.34586,15.89269) (14.28509,16.00642) (14.22207,16.11706) (14.15680,16.22462) (14.08930,16.32909) (14.01958,16.43046) (13.94765,16.52872) (13.87354,16.62389) (13.79725,16.71593) (13.71880,16.80486) (13.63819,16.89066) (13.55545,16.97333) (13.47059,17.05286) (13.38361,17.12924) (13.29453,17.20247) (13.20336,17.27255) (13.11012,17.33945) (13.01481,17.40319)};
\draw[courseaccent,line width=.85pt] plot coordinates {(18.37820,9.74347) (18.26406,9.82800) (18.14657,9.91120) (18.02575,9.99306) (17.90161,10.07358) (17.77416,10.15275) (17.64342,10.23054) (17.50940,10.30696) (17.37212,10.38199) (17.23158,10.45562) (17.08779,10.52784) (16.94078,10.59863) (16.79054,10.66799) (16.63710,10.73590) (16.48047,10.80236) (16.32065,10.86734) (16.15766,10.93084) (15.99152,10.99285) (15.82222,11.05335) (15.64979,11.11233) (15.47424,11.16978) (15.29557,11.22569) (15.11381,11.28004) (14.92895,11.33281) (14.74103,11.38400) (14.55003,11.43359) (14.35599,11.48157) (14.15891,11.52792) (13.95880,11.57263) (13.75568,11.61569) (13.54956,11.65708) (13.34045,11.69678) (13.12836,11.73478) (12.91332,11.77107) (12.69533,11.80563) (12.47441,11.83845) (12.25058,11.86950) (12.02384,11.89878) (11.79422,11.92627) (11.56173,11.95194) (11.32639,11.97580) (11.08821,11.99781) (10.84721,12.01796) (10.60342,12.03624) (10.35684,12.05262) (10.10750,12.06710) (9.85542,12.07965) (9.60062,12.09025) (9.34311,12.09890) (9.08294,12.10556) (8.82010,12.11023) (8.55464,12.11288) (8.28657,12.11349) (8.01592,12.11206) (7.74272,12.10855) (7.46699,12.10296) (7.18877,12.09526) (6.90808,12.08544) (6.62495,12.07347) (6.33942,12.05933) (6.05152,12.04302) (5.76128,12.02451) (5.46874,12.00378) (5.17394,11.98081) (4.87691,11.95558) (4.57770,11.92808) (4.27634,11.89829) (3.97289,11.86619) (3.66737,11.83176) (3.35985,11.79498) (3.05037,11.75584) (2.73898,11.71431) (2.42572,11.67038) (2.11066,11.62404) (1.79385,11.57526) (1.47534,11.52403) (1.15520,11.47033) (0.83348,11.41415) (0.51025,11.35547) (0.18558,11.29427) (-0.14048,11.23055) (-0.46784,11.16428) (-0.79643,11.09545) (-1.12619,11.02405) (-1.45704,10.95007) (-1.78889,10.87349) (-2.12167,10.79431) (-2.45528,10.71250) (-2.78965,10.62807) (-3.12468,10.54101) (-3.46028,10.45129) (-3.79636,10.35893) (-4.13283,10.26390) (-4.46957,10.16621) (-4.80648,10.06585) (-5.14347,9.96282) (-5.48042,9.85711) (-5.81722,9.74873) (-6.15375,9.63767) (-6.48991,9.52393) (-6.82556,9.40752) (-7.16059,9.28845) (-7.49488,9.16671) (-7.82828,9.04232) (-8.16067,8.91528) (-8.49192,8.78560) (-8.82188,8.65330) (-9.15042,8.51838) (-9.47739,8.38087) (-9.80266,8.24079) (-10.12606,8.09814) (-10.44745,7.95295) (-10.76667,7.80525) (-11.08357,7.65506) (-11.39798,7.50240) (-11.70976,7.34731) (-12.01872,7.18981) (-12.32471,7.02995) (-12.62755,6.86776) (-12.92708,6.70327) (-13.22313,6.53652) (-13.51551,6.36757) (-13.80405,6.19645) (-14.08858,6.02321) (-14.36891,5.84791) (-14.64487,5.67059) (-14.91627,5.49132) (-15.18292,5.31014) (-15.44466,5.12713) (-15.70128,4.94235) (-15.95262,4.75586) (-16.19847,4.56773) (-16.43867,4.37803) (-16.67302,4.18684) (-16.90134,3.99423) (-17.12344,3.80029) (-17.33915,3.60509) (-17.54829,3.40872) (-17.75067,3.21126) (-17.94612,3.01282) (-18.13445,2.81347) (-18.31551,2.61332) (-18.48911,2.41245) (-18.65510,2.21098) (-18.81330,2.00899) (-18.96355,1.80660) (-19.10570,1.60391) (-19.23959,1.40103) (-19.36507,1.19806) (-19.48201,0.99512) (-19.59026,0.79233) (-19.68968,0.58978) (-19.78015,0.38761) (-19.86154,0.18592) (-19.93374,-0.01516) (-19.99664,-0.21551) (-20.05013,-0.41503) (-20.09411,-0.61357) (-20.12850,-0.81103) (-20.15322,-1.00728) (-20.16818,-1.20220) (-20.17332,-1.39566) (-20.16859,-1.58755) (-20.15392,-1.77774) (-20.12928,-1.96612) (-20.09464,-2.15255) (-20.04996,-2.33693) (-19.99523,-2.51913) (-19.93045,-2.69903) (-19.85561,-2.87652) (-19.77072,-3.05147) (-19.67581,-3.22379) (-19.57090,-3.39336) (-19.45603,-3.56006) (-19.33125,-3.72379) (-19.19661,-3.88444) (-19.05218,-4.04191) (-18.89803,-4.19611) (-18.73425,-4.34693) (-18.56093,-4.49427) (-18.37817,-4.63806) (-18.18608,-4.77819) (-17.98478,-4.91459) (-17.77439,-5.04718) (-17.55505,-5.17586) (-17.32691,-5.30058) (-17.09011,-5.42127) (-16.84482,-5.53784) (-16.59118,-5.65025) (-16.32939,-5.75843) (-16.05962,-5.86233) (-15.78206,-5.96190) (-15.49688,-6.05709) (-15.20430,-6.14787) (-14.90452,-6.23418) (-14.59774,-6.31601) (-14.28418,-6.39331) (-13.96405,-6.46607) (-13.63758,-6.53426) (-13.30498,-6.59787) (-12.96650,-6.65689) (-12.62235,-6.71130) (-12.27278,-6.76110) (-11.91802,-6.80630) (-11.55831,-6.84689) (-11.19389,-6.88289) (-10.82499,-6.91430) (-10.45187,-6.94114) (-10.07477,-6.96343) (-9.69392,-6.98119) (-9.30958,-6.99446) (-8.92197,-7.00325) (-8.53136,-7.00760) (-8.13797,-7.00755) (-7.74205,-7.00313) (-7.34383,-6.99439) (-6.94355,-6.98138) (-6.54144,-6.96413) (-6.13774,-6.94271) (-5.73268,-6.91716) (-5.32647,-6.88754) (-4.91935,-6.85390) (-4.51153,-6.81631) (-4.10323,-6.77483) (-3.69466,-6.72951) (-3.28604,-6.68043) (-2.87756,-6.62765) (-2.46943,-6.57123) (-2.06184,-6.51125) (-1.65499,-6.44779) (-1.24906,-6.38090) (-0.84424,-6.31066) (-0.44071,-6.23714) (-0.03863,-6.16043) (0.36181,-6.08059) (0.76045,-5.99769) (1.15715,-5.91182) (1.55173,-5.82306) (1.94405,-5.73146) (2.33397,-5.63713) (2.72134,-5.54012) (3.10603,-5.44051) (3.48790,-5.33839) (3.86684,-5.23383) (4.24271,-5.12690) (4.61541,-5.01769) (4.98483,-4.90625) (5.35084,-4.79268) (5.71336,-4.67704) (6.07228,-4.55941) (6.42752,-4.43985) (6.77897,-4.31845) (7.12656,-4.19528) (7.47021,-4.07039) (7.80983,-3.94388) (8.14537,-3.81579) (8.47674,-3.68621) (8.80388,-3.55520) (9.12674,-3.42282) (9.44526,-3.28914) (9.75937,-3.15422) (10.06905,-3.01814) (10.37422,-2.88094) (10.67487,-2.74269) (10.97093,-2.60346) (11.26239,-2.46329) (11.54920,-2.32225) (11.83133,-2.18040) (12.10876,-2.03779) (12.38146,-1.89447) (12.64941,-1.75050) (12.91259,-1.60593) (13.17098,-1.46081) (13.42456,-1.31519) (13.67334,-1.16913) (13.91728,-1.02266) (14.15639,-0.87584) (14.39066,-0.72871) (14.62009,-0.58132) (14.84467,-0.43371) (15.06440,-0.28592) (15.27929,-0.13799) (15.48933,0.01003) (15.69453,0.15811) (15.89490,0.30620) (16.09045,0.45427) (16.28117,0.60229) (16.46709,0.75021) (16.64821,0.89800) (16.82454,1.04564) (16.99610,1.19307) (17.16291,1.34029) (17.32497,1.48724) (17.48230,1.63390) (17.63492,1.78025) (17.78285,1.92625) (17.92611,2.07187) (18.06471,2.21709) (18.19867,2.36188) (18.32802,2.50622) (18.45277,2.65008) (18.57294,2.79343) (18.68856,2.93626) (18.79965,3.07853) (18.90623,3.22023) (19.00832,3.36134) (19.10595,3.50183) (19.19914,3.64168) (19.28791,3.78087) (19.37228,3.91939) (19.45229,4.05721) (19.52794,4.19431) (19.59927,4.33069) (19.66630,4.46631) (19.72906,4.60116) (19.78756,4.73523) (19.84184,4.86850) (19.89191,5.00096) (19.93780,5.13258) (19.97953,5.26335) (20.01713,5.39327) (20.05062,5.52230) (20.08002,5.65045) (20.10535,5.77769) (20.12665,5.90402) (20.14393,6.02941) (20.15722,6.15386) (20.16653,6.27736) (20.17189,6.39988) (20.17332,6.52143) (20.17085,6.64198) (20.16449,6.76153) (20.15427,6.88006) (20.14021,6.99757) (20.12233,7.11404) (20.10064,7.22945) (20.07518,7.34381) (20.04596,7.45709) (20.01299,7.56930) (19.97631,7.68041) (19.93592,7.79042) (19.89186,7.89931) (19.84413,8.00709) (19.79275,8.11372) (19.73775,8.21922) (19.67913,8.32356) (19.61693,8.42674) (19.55115,8.52874) (19.48182,8.62956) (19.40895,8.72919) (19.33255,8.82762) (19.25264,8.92483) (19.16924,9.02081) (19.08237,9.11557) (18.99204,9.20908) (18.89826,9.30134) (18.80105,9.39233) (18.70043,9.48205) (18.59640,9.57049) (18.48899,9.65763) (18.37820,9.74347)};
\fill[courseink] (-0.2128,11.2161) circle[radius=.55];
\fill[courseink] (9.8283,-3.1242) circle[radius=.55];
\draw[vector line,line width=.9pt] (4.5946,4.3504)--(-3.4239,15.8019) node[above left,inner sep=.8mm] {$\mathbf L$};
\draw[vector line,line width=.9pt] (4.5946,4.3504)--(13.0148,17.4032) node[pos=.32,right,inner sep=1mm] {$\mathbf P$};
\draw[vector line,line width=.9pt] (4.5946,4.3504)--(-0.9250,7.9110) node[below left,inner sep=.7mm] {$\mathbf Q$};
\draw[courseink,line width=.55pt] (6.0541,6.6129)--(4.4350,7.6573)--(2.9755,5.3948);
\fill[courseink] (13.0148,17.4032) circle[radius=.65];
\node[anchor=west,inner sep=1mm] at (13.0148,17.4032) {Перицентр};
\fill[courseink] (4.5946,4.3504) circle[radius=.65];
\node[anchor=west,xshift=1mm,yshift=.5mm,inner sep=.5mm] at (4.5946,4.3504) {$O$};
\node[courseink,anchor=west,inner sep=.5mm] at (6,-15) {Орбита 1};
\node[courseaccent,anchor=west,inner sep=.5mm] at (19,3) {Орбита 2};
\end{tikzpicture}
\par
{\usebeamerfont{caption}\color{coursemuted}
$\mathbf P$ --- к перицентру орбиты 1;\par
$\mathbf Q$ --- на $90^\circ$ вперёд в её плоскости.\par}
\end{column}
\end{columns}
\smallskip
\justifying
Равенство радиусов на одном луче означает пересечение орбит.
Параллельные нормали — компланарный случай.
\end{frame}

% Слайд 15
\begin{frame}{Переход от элементов к векторам состояния}
Единичный вектор перицентра в неподвижных пространственных осях:
\[
\mathbf P=\begin{pmatrix}
\cos\Omega\cos\omega-\sin\Omega\sin\omega\cos i\\
\sin\Omega\cos\omega+\cos\Omega\sin\omega\cos i\\
\sin\omega\sin i
\end{pmatrix},\qquad
\mathbf Q=\widehat{\mathbf h}\times\mathbf P.
\]
$\mathbf Q$ лежит в плоскости орбиты, на $90^\circ$ впереди $\mathbf P$ по движению.
\[
\boxed{\mathbf r=\frac{p}{1+e\cos\nu}(\cos\nu\,\mathbf P+\sin\nu\,\mathbf Q),}
\]
\[
\boxed{\mathbf v=\sqrt{\frac\mu p}\bigl[-\sin\nu\,\mathbf P+(e+\cos\nu)\mathbf Q\bigr].}
\]
\justifying
На луче $\mathbf s$ находим $\cos\nu=\mathbf P\cdot\mathbf s$, $\sin\nu=\mathbf Q\cdot\mathbf s$. По ним выбираем угол $\nu$ от $0$ до $360^\circ$.
\end{frame}

% Слайд 16
\begin{frame}{Некомпланарный переход: исходные данные}
\centering
\renewcommand{\arraystretch}{1.05}
\begin{tabular}{@{}lrr@{}}
\toprule
 & Начальная орбита & Конечная орбита\\\midrule
$a$, км & $18\,654{,}3640$ & $20\,679{,}0085$\\
$e$ & $0{,}28969592$ & $0{,}16932267$\\
$i$, град & $35{,}62718$ & $38{,}46693$\\
$\Omega$, град & $89{,}44357$ & $107{,}25881$\\
$\omega$, град & $265{,}55428$ & $234{,}60634$\\\bottomrule
\end{tabular}\par\smallskip
\raggedright
Находим нормали и направление пересечения плоскостей:
\[
\widehat{\mathbf h}_j=(\sin i_j\sin\Omega_j,\;-\sin i_j\cos\Omega_j,\;\cos i_j)^{\mathsf T},
\]
\[
\mathbf L=\frac{\widehat{\mathbf h}_1\times\widehat{\mathbf h}_2}
{\|\widehat{\mathbf h}_1\times\widehat{\mathbf h}_2\|}
=(-0{,}804421,\;0{,}139578,\;0{,}577430)^{\mathsf T}.
\]

\end{frame}

% Слайд 17
\begin{frame}{Проверка радиусов на двух лучах}
Параметры орбит:
$p_1=17\,088{,}8203$ км, $p_2=20\,086{,}1379$ км.
Например, на положительном луче
\[
r_1=\frac{17\,088{,}8203}{1+0{,}28969592\cos176{,}87446^\circ}
=24\,043{,}870\text{ км}.
\]
\centering
\renewcommand{\arraystretch}{1.0}
\begin{tabular}{@{}lrr@{}}
\toprule
 & $+\mathbf L$ & $-\mathbf L$\\\midrule
$\nu_1$, град & $176{,}87446$ & $356{,}87446$\\
$\nu_2$, град & $193{,}55739$ & $13{,}55739$\\
$r_1$, км & $24\,043{,}870$ & $13\,254{,}700$\\
$r_2$, км & $24\,043{,}869$ & $17\,247{,}174$\\\bottomrule
\end{tabular}\par\smallskip
\justifying
На $+\mathbf L$ радиусы совпадают с точностью округления исходных данных
(расхождение $0{,}8$ м). На $-\mathbf L$ разность $3992$ км: переход невозможен.
\end{frame}

% Слайд 18
\begin{frame}{Некомпланарный переход: вектор импульса}
Для найденных аномалий вычислим состояния в общих неподвижных осях:
\begin{align*}
\mathbf r_1&=(-19\,341{,}3833,\;3355{,}9909,\;13\,883{,}6559)^{\mathsf T}\text{ км},\\
\mathbf v_1&=(-462{,}0121,\;-3388{,}2747,\;307{,}5030)^{\mathsf T}\text{ м/с},\\
\mathbf v_2&=(132{,}1330,\;-3645{,}1336,\;758{,}9684)^{\mathsf T}\text{ м/с}.
\end{align*}
Радиус-вектор целевой орбиты совпадает с $\mathbf r_1$ в пределах $0{,}8$ м.
\[
\Delta\mathbf v=\mathbf v_2-\mathbf v_1
=(594{,}1451,\;-256{,}8588,\;451{,}4654)^{\mathsf T}\text{ м/с},
\]
\[
\Delta v=\sqrt{594{,}1451^2+256{,}8588^2+451{,}4654^2}
=789{,}1805\text{ м/с}.
\]
\justifying
Нормальная к исходной плоскости компонента импульса обеспечивает переход
в плоскость целевой орбиты.
\end{frame}

% Коррекция отдельных элементов: число независимых требований.
% Слайд 19
\begin{frame}{Коррекция отдельных элементов}
\justifying
Теперь задаём не всю конечную орбиту, а только часть её элементов.
Место манёвра сначала считаем известным.
\begin{columns}[T,onlytextwidth]
\begin{column}{.48\textwidth}
\begin{block}{Импульс в плоскости орбиты}
Две компоненты: $\Delta v_r$, $\Delta v_t$.
\par\medskip
Зададим $a_2,e_2$ или $a_2,\omega_2$.
\end{block}
\justifying
При сохранении направления движения:
\[
i_2=i_1,\qquad\Omega_2=\Omega_1,
\]
\[
u_2=u_1,\qquad u=\omega+\nu.
\]
\end{column}
\begin{column}{.48\textwidth}
\begin{block}{Импульс вне плоскости орбиты}
Три компоненты: $\Delta v_r$, $\Delta v_t$, $\Delta v_z$.
\par\medskip
Зададим, например, $a_2,e_2,\Omega_2$.
\end{block}
\justifying
Остальные элементы определятся из условия, что конечная орбита проходит через точку манёвра.
\end{column}
\end{columns}
\end{frame}

% Линеаризация в фиксированной точке манёвра.
% Слайд 20
\begin{frame}{Линеаризация изменения элементов}
\setlength{\abovedisplayskip}{3pt}
\setlength{\belowdisplayskip}{3pt}
Точка манёвра фиксирована: $\mathbf r_2=\mathbf r_1=\mathbf r$,
$\mathbf v_2=\mathbf v_1+\Delta\mathbf v$.
\begin{gather*}
\mathbf y=(a,e,\omega,i,\Omega)^{\mathsf T}=\mathbf F(\mathbf r,\mathbf v),\\[3pt]
\Delta\mathbf y=
\mathbf F(\mathbf r,\mathbf v_1+\Delta\mathbf v)
-\mathbf F(\mathbf r,\mathbf v_1).
\end{gather*}
\begin{block}{Разложение по малому импульсу}
\centering
$\displaystyle
\Delta\mathbf y=B\,\Delta\mathbf v
+O\!\left(\|\Delta\mathbf v\|^2\right),\qquad
B_{kj}=\left.\frac{\partial y_k}{\partial v_j}\right|_{\mathbf r,\mathbf v_1}.
$
\par\medskip
В первом приближении: $\boxed{B\,\Delta\mathbf v\approx\Delta\mathbf y}$.
\end{block}
\justifying
$j=r,t,z$ --- компоненты в исходном орбитальном базисе.
Матрицу $B$ вычисляем до манёвра. Линеаризация требует малости импульса.
\end{frame}

% Слайд 21
\begin{frame}{Две компоненты импульса в плоскости}
\setlength{\abovedisplayskip}{3pt}
\setlength{\belowdisplayskip}{3pt}
При $\Delta v_z=0$ получаем пять уравнений для двух неизвестных:
\begin{columns}[T,onlytextwidth]
\begin{column}{.56\textwidth}
\[
\begin{pmatrix}
\Delta a\\ \Delta e\\ \Delta\omega\\ \Delta i\\ \Delta\Omega
\end{pmatrix}
\approx
\underbrace{\begin{pmatrix}
B_{ar}&B_{at}\\
B_{er}&B_{et}\\
B_{\omega r}&B_{\omega t}\\
0&0\\
0&0
\end{pmatrix}}_{B_{\parallel}\; (5\times2)}
\begin{pmatrix}\Delta v_r\\ \Delta v_t\end{pmatrix}.
\]
\end{column}
\begin{column}{.41\textwidth}
\smallskip\raggedright
Зададим $\Delta a,\Delta e$. Если первые два уравнения независимы,
из них найдём $\Delta v_r,\Delta v_t$.
\par\smallskip
Тогда третье уравнение определит $\Delta\omega$.
Его значение уже нельзя выбрать независимо.
\end{column}
\end{columns}
\begin{block}{Не более двух независимых требований}
\justifying
$\operatorname{rank}B_{\parallel}\leq2$. Могут измениться сразу $a,e,\omega$,
но все три уравнения должны выполняться при одних и тех же $\Delta v_r,\Delta v_t$.
\end{block}
\end{frame}

% Слайд 22
\begin{frame}{Три компоненты импульса}
\setlength{\abovedisplayskip}{3pt}
\setlength{\belowdisplayskip}{3pt}
В первом приближении пять изменений определяются тремя компонентами:
\begin{columns}[T,onlytextwidth]
\begin{column}{.59\textwidth}
\[
\begin{pmatrix}
\Delta a\\ \Delta e\\ \Delta\omega\\ \Delta i\\ \Delta\Omega
\end{pmatrix}
\approx
\underbrace{\begin{pmatrix}
B_{ar}&B_{at}&0\\
B_{er}&B_{et}&0\\
B_{\omega r}&B_{\omega t}&B_{\omega z}\\
0&0&B_{iz}\\
0&0&B_{\Omega z}
\end{pmatrix}}_{B\; (5\times3)}
\begin{pmatrix}\Delta v_r\\ \Delta v_t\\ \Delta v_z\end{pmatrix}.
\]
\end{column}
\begin{column}{.38\textwidth}
\smallskip\raggedright
Можно задать, например, $\Delta a,\Delta e,\Delta\Omega$,
если эти три уравнения независимы.
\par\smallskip
Найдём импульс, а затем $\Delta\omega,\Delta i$.
Изменения $i,\Omega$ связаны: оба зависят от одной $\Delta v_z$.
\end{column}
\end{columns}
\begin{block}{Число компонент задаёт верхнюю границу}
\justifying
$\operatorname{rank}B\leq3$. Выбранные уравнения должны быть независимы,
а остальные требования --- совместны с ними.
\end{block}
\end{frame}

% Слайд 23
\begin{frame}{Коррекция $a$ и $e$}
Заданы $a_2,e_2$. Из уравнения конечного эллипса получаем
\[
r=\frac{a_2(1-e_2^2)}{1+e_2\cos\nu_2}
\quad\Longrightarrow\quad
\cos\nu_2=\frac{a_2(1-e_2^2)/r-1}{e_2}=c.
\]
При $|c|<1$ в точке манёвра возможны два варианта:
\begin{columns}[T,onlytextwidth]
\begin{column}{.48\textwidth}
\begin{block}{Вариант A}
$\nu_{2A}=\arccos c$,\par\smallskip
$\omega_{2A}=u-\nu_{2A}$.
\end{block}
\end{column}
\begin{column}{.48\textwidth}
\begin{block}{Вариант B}
$\nu_{2B}=360^\circ-\arccos c$,\par\smallskip
$\omega_{2B}=u-\nu_{2B}$.
\end{block}
\end{column}
\end{columns}
\medskip
\justifying
Обе орбиты имеют заданные размер и форму, но по-разному ориентированы.
Для каждой находим скорость и импульс $\Delta\mathbf v=\mathbf v_2-\mathbf v_1$.
Углы приводим к диапазону $[0,360^\circ)$.
\end{frame}

% Слайд 24
\begin{frame}{Коррекция $a$ и $e$: расчёт элементов}
\begin{columns}[T,onlytextwidth]
\begin{column}{.48\textwidth}
\begin{block}{Исходное состояние}
$a_1=26\,000$ км,\quad $e_1=0{,}6$,\par
$i_1=64^\circ$,\quad $\Omega_1=0^\circ$,\par
$\omega_1=250^\circ$,\quad $\nu_1=60^\circ$.
\end{block}
\end{column}
\begin{column}{.48\textwidth}
\begin{block}{Требуется получить}
$a_2=26\,554$ км,\quad $e_2=0{,}7$.
\par\medskip
Найти обе конечные орбиты и импульсы.
\end{block}
\end{column}
\end{columns}
\[
r=\frac{26\,000(1-0{,}6^2)}{1+0{,}6\cos60^\circ}=12\,800\text{ км},
\qquad u=250^\circ+60^\circ=310^\circ.
\]
\[
\cos\nu_2=\frac{26\,554(1-0{,}7^2)/12\,800-1}{0{,}7}=0{,}08287277.
\]
Отсюда $\nu_{2A}=85{,}2463^\circ$, $\nu_{2B}=274{,}7537^\circ$.
Для каждого варианта $\omega_2=310^\circ-\nu_2$.
\end{frame}

% Слайд 25
\begin{frame}{Коррекция $a$ и $e$: конечные орбиты}
\begin{columns}[T,onlytextwidth]
\begin{column}{.48\textwidth}
\centering
\renewcommand{\arraystretch}{1.14}
\begin{tabular}{@{}lrr@{}}
\toprule
Элемент & A & B\\\midrule
$a_2$, км & $26\,554$ & $26\,554$\\
$e_2$ & $0{,}7$ & $0{,}7$\\
$i_2$, град & $64$ & $64$\\
$\Omega_2$, град & $0$ & $0$\\
$\omega_2$, град & $224{,}7537$ & $35{,}2463$\\
$\nu_2$, град & $85{,}2463$ & $274{,}7537$\\\bottomrule
\end{tabular}\par\medskip
\justifying
В варианте A КА удаляется от перицентра, в варианте B приближается к нему.
\end{column}
\begin{column}{.50\textwidth}
\centering\begin{tikzpicture}[course diagram,x=1mm,y=1mm]
\path[use as bounding box] (-34.0, -24.5) rectangle (34.0, 24.5);
\filldraw[fill=coursetint,draw=courserule,line width=.4pt] (0.9358,-2.2467) circle[radius=3.7771];
\node[inner sep=0pt,font=\fontsize{8.5}{10}\selectfont] at (0.9358,-2.2467) {$\oplus$};
\draw[draw=courseink,line width=.75pt] plot coordinates {(-1.17060,-8.03409) (-1.01864,-8.08800) (-0.86575,-8.13943) (-0.71194,-8.18839) (-0.55718,-8.23488) (-0.40146,-8.27888) (-0.24477,-8.32040) (-0.08708,-8.35944) (0.07162,-8.39597) (0.23134,-8.42999) (0.39212,-8.46149) (0.55396,-8.49045) (0.71688,-8.51685) (0.88092,-8.54066) (1.04608,-8.56187) (1.21238,-8.58045) (1.37986,-8.59638) (1.54852,-8.60960) (1.71840,-8.62011) (1.88952,-8.62785) (2.06189,-8.63279) (2.23554,-8.63489) (2.41050,-8.63410) (2.58678,-8.63037) (2.76441,-8.62364) (2.94341,-8.61388) (3.12382,-8.60101) (3.30564,-8.58497) (3.48890,-8.56571) (3.67363,-8.54314) (3.85984,-8.51721) (4.04756,-8.48782) (4.23682,-8.45490) (4.42762,-8.41837) (4.62000,-8.37814) (4.81396,-8.33410) (5.00954,-8.28618) (5.20674,-8.23425) (5.40558,-8.17822) (5.60608,-8.11797) (5.80825,-8.05338) (6.01209,-7.98433) (6.21763,-7.91070) (6.42486,-7.83234) (6.63379,-7.74912) (6.84442,-7.66088) (7.05674,-7.56749) (7.27076,-7.46878) (7.48647,-7.36458) (7.70384,-7.25473) (7.92286,-7.13904) (8.14352,-7.01733) (8.36577,-6.88941) (8.58959,-6.75508) (8.81494,-6.61413) (9.04176,-6.46635) (9.27000,-6.31151) (9.49959,-6.14940) (9.73047,-5.97977) (9.96254,-5.80239) (10.19571,-5.61699) (10.42988,-5.42334) (10.66492,-5.22115) (10.90070,-5.01018) (11.13707,-4.79013) (11.37387,-4.56073) (11.61091,-4.32170) (11.84800,-4.07274) (12.08490,-3.81357) (12.32139,-3.54389) (12.55719,-3.26341) (12.79201,-2.97183) (13.02554,-2.66886) (13.25743,-2.35420) (13.48732,-2.02759) (13.71478,-1.68874) (13.93940,-1.33738) (14.16069,-0.97327) (14.37815,-0.59617) (14.59122,-0.20587) (14.79933,0.19783) (15.00184,0.61508) (15.19810,1.04602) (15.38738,1.49074) (15.56892,1.94929) (15.74194,2.42167) (15.90558,2.90782) (16.05895,3.40762) (16.20113,3.92090) (16.33114,4.44739) (16.44797,4.98675) (16.55057,5.53854) (16.63786,6.10224) (16.70874,6.67720) (16.76207,7.26268) (16.79671,7.85781) (16.81151,8.46160) (16.80534,9.07292) (16.77704,9.69053) (16.72553,10.31301) (16.64974,10.93885) (16.54864,11.56636) (16.42129,12.19374) (16.26683,12.81905) (16.08450,13.44022) (15.87365,14.05508) (15.63375,14.66134) (15.36446,15.25663) (15.06555,15.83853) (14.73702,16.40454) (14.37901,16.95216) (13.99190,17.47888) (13.57626,17.98222) (13.13288,18.45978) (12.66274,18.90921) (12.16706,19.32830) (11.64726,19.71498) (11.10494,20.06737) (10.54190,20.38378) (9.96010,20.66273) (9.36164,20.90300) (8.74874,21.10362) (8.12375,21.26391) (7.48906,21.38345) (6.84710,21.46210) (6.20036,21.50000) (5.55126,21.49757) (4.90223,21.45549) (4.25561,21.37467) (3.61366,21.25625) (2.97855,21.10159) (2.35230,20.91221) (1.73680,20.68980) (1.13379,20.43617) (0.54485,20.15325) (-0.02861,19.84303) (-0.58536,19.50755) (-1.12432,19.14891) (-1.64458,18.76918) (-2.14540,18.37043) (-2.62621,17.95471) (-3.08655,17.52402) (-3.52614,17.08028) (-3.94479,16.62535) (-4.34248,16.16102) (-4.71925,15.68898) (-5.07525,15.21081) (-5.41074,14.72803) (-5.72602,14.24203) (-6.02149,13.75410) (-6.29758,13.26545) (-6.55478,12.77718) (-6.79360,12.29030) (-7.01462,11.80572) (-7.21839,11.32425) (-7.40553,10.84666) (-7.57663,10.37358) (-7.73230,9.90562) (-7.87317,9.44328) (-7.99983,8.98701) (-8.11290,8.53720) (-8.21297,8.09418) (-8.30063,7.65823) (-8.37645,7.22958) (-8.44097,6.80841) (-8.49476,6.39488) (-8.53832,5.98909) (-8.57216,5.59112) (-8.59677,5.20102) (-8.61263,4.81882) (-8.62017,4.44451) (-8.61983,4.07808) (-8.61201,3.71949) (-8.59712,3.36867) (-8.57553,3.02557) (-8.54759,2.69009) (-8.51365,2.36215) (-8.47402,2.04165) (-8.42902,1.72847) (-8.37893,1.42251) (-8.32403,1.12364) (-8.26457,0.83174) (-8.20082,0.54668) (-8.13299,0.26833) (-8.06132,-0.00343) (-7.98600,-0.26875) (-7.90724,-0.52776) (-7.82521,-0.78058) (-7.74010,-1.02735) (-7.65207,-1.26821) (-7.56128,-1.50327) (-7.46786,-1.73268) (-7.37195,-1.95656) (-7.27369,-2.17503) (-7.17319,-2.38822) (-7.07057,-2.59624) (-6.96593,-2.79922) (-6.85938,-2.99727) (-6.75100,-3.19050) (-6.64088,-3.37903) (-6.52911,-3.56295) (-6.41576,-3.74238) (-6.30090,-3.91741) (-6.18459,-4.08815) (-6.06690,-4.25468) (-5.94789,-4.41711) (-5.82761,-4.57552) (-5.70610,-4.73000) (-5.58342,-4.88063) (-5.45960,-5.02749) (-5.33469,-5.17067) (-5.20871,-5.31024) (-5.08171,-5.44626) (-4.95370,-5.57882) (-4.82473,-5.70798) (-4.69480,-5.83379) (-4.56396,-5.95634) (-4.43220,-6.07567) (-4.29956,-6.19184) (-4.16605,-6.30490) (-4.03167,-6.41492) (-3.89645,-6.52193) (-3.76040,-6.62599) (-3.62351,-6.72714) (-3.48580,-6.82542) (-3.34727,-6.92088) (-3.20793,-7.01354) (-3.06778,-7.10346) (-2.92681,-7.19065) (-2.78504,-7.27516) (-2.64245,-7.35700) (-2.49904,-7.43622) (-2.35481,-7.51282) (-2.20976,-7.58685) (-2.06387,-7.65831) (-1.91715,-7.72722) (-1.76957,-7.79360) (-1.62114,-7.85747) (-1.47184,-7.91883) (-1.32167,-7.97771) (-1.17060,-8.03409)};
\draw[draw=courseaccent,line width=1pt] plot coordinates {(-2.41428,-5.56812) (-2.32665,-5.65516) (-2.23765,-5.74087) (-2.14727,-5.82527) (-2.05549,-5.90837) (-1.96229,-5.99017) (-1.86766,-6.07069) (-1.77156,-6.14993) (-1.67399,-6.22789) (-1.57490,-6.30459) (-1.47429,-6.38001) (-1.37211,-6.45417) (-1.26835,-6.52707) (-1.16297,-6.59869) (-1.05594,-6.66904) (-0.94723,-6.73811) (-0.83680,-6.80589) (-0.72461,-6.87237) (-0.61064,-6.93755) (-0.49483,-7.00141) (-0.37716,-7.06394) (-0.25757,-7.12512) (-0.13602,-7.18493) (-0.01247,-7.24335) (0.11314,-7.30036) (0.24084,-7.35593) (0.37070,-7.41003) (0.50277,-7.46264) (0.63710,-7.51372) (0.77375,-7.56324) (0.91278,-7.61116) (1.05426,-7.65743) (1.19824,-7.70202) (1.34479,-7.74487) (1.49398,-7.78594) (1.64588,-7.82516) (1.80057,-7.86248) (1.95812,-7.89784) (2.11860,-7.93116) (2.28211,-7.96239) (2.44871,-7.99143) (2.61850,-8.01822) (2.79157,-8.04266) (2.96800,-8.06466) (3.14789,-8.08412) (3.33133,-8.10094) (3.51842,-8.11501) (3.70927,-8.12620) (3.90398,-8.13441) (4.10265,-8.13948) (4.30540,-8.14129) (4.51233,-8.13968) (4.72355,-8.13449) (4.93919,-8.12556) (5.15936,-8.11271) (5.38418,-8.09575) (5.61377,-8.07448) (5.84826,-8.04869) (6.08776,-8.01815) (6.33240,-7.98264) (6.58230,-7.94188) (6.83759,-7.89564) (7.09839,-7.84361) (7.36482,-7.78551) (7.63700,-7.72103) (7.91505,-7.64983) (8.19906,-7.57156) (8.48916,-7.48587) (8.78543,-7.39236) (9.08796,-7.29063) (9.39684,-7.18025) (9.71212,-7.06077) (10.03387,-6.93172) (10.36212,-6.79262) (10.69687,-6.64294) (11.03814,-6.48215) (11.38588,-6.30969) (11.74003,-6.12498) (12.10050,-5.92742) (12.46715,-5.71639) (12.83981,-5.49124) (13.21826,-5.25132) (13.60221,-4.99595) (13.99134,-4.72446) (14.38523,-4.43615) (14.78340,-4.13033) (15.18531,-3.80631) (15.59030,-3.46341) (15.99763,-3.10097) (16.40644,-2.71834) (16.81578,-2.31494) (17.22456,-1.89020) (17.63156,-1.44364) (18.03544,-0.97484) (18.43470,-0.48349) (18.82769,0.03063) (19.21263,0.56760) (19.58756,1.12735) (19.95038,1.70964) (20.29884,2.31403) (20.63055,2.93987) (20.94297,3.58628) (21.23347,4.25210) (21.49932,4.93591) (21.73772,5.63599) (21.94581,6.35032) (22.12077,7.07655) (22.25981,7.81205) (22.36021,8.55383) (22.41941,9.29865) (22.43505,10.04298) (22.40502,10.78304) (22.32750,11.51486) (22.20106,12.23431) (22.02466,12.93717) (21.79772,13.61920) (21.52015,14.27617) (21.19237,14.90399) (20.81533,15.49876) (20.39050,16.05682) (19.91987,16.57484) (19.40591,17.04991) (18.85152,17.47952) (18.26002,17.86166) (17.63504,18.19482) (16.98048,18.47803) (16.30043,18.71083) (15.59911,18.89326) (14.88078,19.02589) (14.14965,19.10969) (13.40988,19.14608) (12.66544,19.13684) (11.92014,19.08404) (11.17752,18.99002) (10.44086,18.85732) (9.71314,18.68861) (8.99705,18.48666) (8.29495,18.25430) (7.60887,17.99433) (6.94058,17.70957) (6.29151,17.40271) (5.66284,17.07640) (5.05548,16.73315) (4.47009,16.37534) (3.90714,16.00524) (3.36688,15.62493) (2.84940,15.23637) (2.35463,14.84135) (1.88238,14.44152) (1.43234,14.03837) (1.00410,13.63326) (0.59720,13.22740) (0.21107,12.82190) (-0.15486,12.41770) (-0.50123,12.01567) (-0.82870,11.61657) (-1.13793,11.22103) (-1.42961,10.82964) (-1.70444,10.44286) (-1.96310,10.06112) (-2.20626,9.68474) (-2.43461,9.31403) (-2.64879,8.94921) (-2.84944,8.59045) (-3.03718,8.23790) (-3.21262,7.89166) (-3.37634,7.55179) (-3.52889,7.21833) (-3.67081,6.89129) (-3.80262,6.57066) (-3.92481,6.25642) (-4.03784,5.94850) (-4.14217,5.64685) (-4.23822,5.35140) (-4.32641,5.06205) (-4.40711,4.77872) (-4.48070,4.50130) (-4.54752,4.22968) (-4.60791,3.96376) (-4.66217,3.70342) (-4.71061,3.44853) (-4.75351,3.19899) (-4.79113,2.95466) (-4.82372,2.71543) (-4.85153,2.48118) (-4.87477,2.25178) (-4.89367,2.02712) (-4.90841,1.80707) (-4.91919,1.59151) (-4.92619,1.38034) (-4.92958,1.17343) (-4.92952,0.97068) (-4.92615,0.77197) (-4.91962,0.57720) (-4.91007,0.38626) (-4.89761,0.19905) (-4.88237,0.01547) (-4.86445,-0.16458) (-4.84397,-0.34119) (-4.82102,-0.51446) (-4.79570,-0.68448) (-4.76808,-0.85133) (-4.73827,-1.01509) (-4.70632,-1.17586) (-4.67232,-1.33370) (-4.63633,-1.48870) (-4.59841,-1.64094) (-4.55863,-1.79048) (-4.51704,-1.93739) (-4.47370,-2.08175) (-4.42864,-2.22362) (-4.38192,-2.36305) (-4.33358,-2.50013) (-4.28365,-2.63489) (-4.23218,-2.76741) (-4.17919,-2.89773) (-4.12472,-3.02591) (-4.06880,-3.15200) (-4.01144,-3.27605) (-3.95268,-3.39810) (-3.89253,-3.51822) (-3.83102,-3.63642) (-3.76816,-3.75277) (-3.70396,-3.86730) (-3.63844,-3.98006) (-3.57161,-4.09106) (-3.50348,-4.20037) (-3.43405,-4.30800) (-3.36334,-4.41399) (-3.29134,-4.51838) (-3.21806,-4.62119) (-3.14350,-4.72245) (-3.06766,-4.82218) (-2.99054,-4.92043) (-2.91213,-5.01720) (-2.83243,-5.11252) (-2.75143,-5.20642) (-2.66913,-5.29891) (-2.58551,-5.39001) (-2.50056,-5.47974) (-2.41428,-5.56812)};
\draw[draw=coursemuted,line width=.8pt,dashed] plot coordinates {(4.78857,0.47579) (4.71652,0.57611) (4.64290,0.67534) (4.56770,0.77351) (4.49091,0.87063) (4.41250,0.96670) (4.33247,1.06174) (4.25078,1.15576) (4.16743,1.24878) (4.08237,1.34078) (3.99560,1.43179) (3.90707,1.52181) (3.81678,1.61084) (3.72467,1.69889) (3.63073,1.78595) (3.53492,1.87202) (3.43721,1.95711) (3.33754,2.04122) (3.23590,2.12433) (3.13224,2.20644) (3.02650,2.28755) (2.91866,2.36764) (2.80866,2.44670) (2.69646,2.52473) (2.58199,2.60170) (2.46522,2.67760) (2.34608,2.75241) (2.22452,2.82611) (2.10047,2.89868) (1.97387,2.97009) (1.84467,3.04031) (1.71278,3.10932) (1.57814,3.17708) (1.44068,3.24355) (1.30032,3.30869) (1.15698,3.37246) (1.01059,3.43482) (0.86104,3.49572) (0.70827,3.55509) (0.55217,3.61290) (0.39265,3.66906) (0.22962,3.72352) (0.06296,3.77621) (-0.10741,3.82705) (-0.28161,3.87596) (-0.45976,3.92285) (-0.64196,3.96763) (-0.82834,4.01019) (-1.01902,4.05045) (-1.21412,4.08827) (-1.41378,4.12354) (-1.61814,4.15613) (-1.82732,4.18590) (-2.04147,4.21272) (-2.26074,4.23641) (-2.48527,4.25681) (-2.71522,4.27376) (-2.95074,4.28705) (-3.19200,4.29649) (-3.43914,4.30187) (-3.69235,4.30296) (-3.95177,4.29952) (-4.21758,4.29128) (-4.48995,4.27799) (-4.76904,4.25935) (-5.05503,4.23505) (-5.34807,4.20478) (-5.64833,4.16818) (-5.95598,4.12489) (-6.27116,4.07452) (-6.59403,4.01668) (-6.92472,3.95092) (-7.26336,3.87679) (-7.61007,3.79381) (-7.96496,3.70148) (-8.32809,3.59927) (-8.69954,3.48662) (-9.07933,3.36295) (-9.46748,3.22764) (-9.86396,3.08007) (-10.26869,2.91957) (-10.68157,2.74545) (-11.10243,2.55701) (-11.53105,2.35352) (-11.96715,2.13423) (-12.41038,1.89838) (-12.86028,1.64520) (-13.31635,1.37390) (-13.77795,1.08372) (-14.24434,0.77387) (-14.71469,0.44362) (-15.18801,0.09224) (-15.66319,-0.28096) (-16.13895,-0.67661) (-16.61388,-1.09526) (-17.08640,-1.53741) (-17.55474,-2.00342) (-18.01698,-2.49355) (-18.47100,-3.00791) (-18.91450,-3.54644) (-19.34503,-4.10890) (-19.75993,-4.69482) (-20.15642,-5.30351) (-20.53157,-5.93402) (-20.88232,-6.58511) (-21.20555,-7.25525) (-21.49806,-7.94262) (-21.75667,-8.64504) (-21.97822,-9.36005) (-22.15963,-10.08487) (-22.29800,-10.81639) (-22.39062,-11.55124) (-22.43505,-12.28581) (-22.42918,-13.01627) (-22.37130,-13.73861) (-22.26013,-14.44875) (-22.09489,-15.14254) (-21.87531,-15.81589) (-21.60169,-16.46476) (-21.27488,-17.08532) (-20.89628,-17.67397) (-20.46784,-18.22741) (-19.99203,-18.74269) (-19.47177,-19.21728) (-18.91041,-19.64910) (-18.31162,-20.03653) (-17.67936,-20.37846) (-17.01781,-20.67423) (-16.33125,-20.92368) (-15.62401,-21.12710) (-14.90040,-21.28519) (-14.16467,-21.39903) (-13.42088,-21.47006) (-12.67293,-21.50000) (-11.92447,-21.49080) (-11.17888,-21.44460) (-10.43927,-21.36371) (-9.70843,-21.25051) (-8.98884,-21.10744) (-8.28269,-20.93697) (-7.59186,-20.74154) (-6.91792,-20.52355) (-6.26221,-20.28533) (-5.62576,-20.02914) (-5.00941,-19.75710) (-4.41375,-19.47126) (-3.83920,-19.17351) (-3.28598,-18.86564) (-2.75417,-18.54930) (-2.24372,-18.22602) (-1.75445,-17.89721) (-1.28610,-17.56414) (-0.83829,-17.22798) (-0.41063,-16.88978) (-0.00261,-16.55049) (0.38628,-16.21095) (0.75660,-15.87193) (1.10892,-15.53409) (1.44386,-15.19802) (1.76202,-14.86424) (2.06401,-14.53321) (2.35046,-14.20530) (2.62195,-13.88087) (2.87910,-13.56018) (3.12250,-13.24348) (3.35273,-12.93098) (3.57033,-12.62282) (3.77587,-12.31914) (3.96986,-12.02003) (4.15281,-11.72558) (4.32523,-11.43583) (4.48757,-11.15082) (4.64029,-10.87054) (4.78382,-10.59501) (4.91859,-10.32421) (5.04498,-10.05810) (5.16338,-9.79664) (5.27415,-9.53979) (5.37763,-9.28750) (5.47416,-9.03969) (5.56403,-8.79631) (5.64756,-8.55728) (5.72502,-8.32253) (5.79668,-8.09197) (5.86280,-7.86552) (5.92361,-7.64312) (5.97936,-7.42466) (6.03024,-7.21006) (6.07648,-6.99925) (6.11827,-6.79213) (6.15579,-6.58863) (6.18921,-6.38865) (6.21871,-6.19212) (6.24445,-5.99894) (6.26656,-5.80905) (6.28519,-5.62235) (6.30048,-5.43877) (6.31256,-5.25824) (6.32153,-5.08067) (6.32751,-4.90599) (6.33062,-4.73413) (6.33095,-4.56501) (6.32859,-4.39857) (6.32364,-4.23473) (6.31617,-4.07344) (6.30628,-3.91462) (6.29403,-3.75822) (6.27950,-3.60416) (6.26275,-3.45239) (6.24384,-3.30286) (6.22283,-3.15550) (6.19979,-3.01026) (6.17475,-2.86709) (6.14777,-2.72593) (6.11889,-2.58673) (6.08816,-2.44945) (6.05561,-2.31403) (6.02128,-2.18043) (5.98520,-2.04861) (5.94741,-1.91853) (5.90793,-1.79013) (5.86678,-1.66338) (5.82400,-1.53825) (5.77960,-1.41469) (5.73361,-1.29266) (5.68603,-1.17214) (5.63689,-1.05309) (5.58619,-0.93547) (5.53396,-0.81925) (5.48019,-0.70441) (5.42490,-0.59091) (5.36809,-0.47872) (5.30977,-0.36783) (5.24993,-0.25819) (5.18858,-0.14980) (5.12572,-0.04262) (5.06135,0.06337) (4.99545,0.16818) (4.92803,0.27184) (4.85907,0.37437) (4.78857,0.47579)};
\draw[aux line] (0.9358,-2.2467)--(5.8082,-8.0534);
\fill[courseink] (5.8082,-8.0534) circle[radius=.85];
\node[anchor=south west,inner sep=1mm] at (5.8082,-8.0534) {$P$};
\end{tikzpicture}
\par\figurelegend
\end{column}
\end{columns}
\end{frame}

% Слайд 26
\begin{frame}{Коррекция $a$ и $e$: скорости и импульсы}
Для обоих вариантов получаем одну точку:
\[
\mathbf r_1=\mathbf r_{2A}=\mathbf r_{2B}
=(8227{,}6814,\;-4298{,}3908,\;-8813{,}0072)^{\mathsf T}\text{ км}.
\]
\centering
\renewcommand{\arraystretch}{1.04}
\begin{tabular}{@{}lrrr@{}}
\toprule
Скорость, м/с & $v_x$ & $v_y$ & $v_z$\\\midrule
$\mathbf v_1$ & $6508{,}7585$ & $938{,}8305$ & $1924{,}8879$\\
$\mathbf v_{2A}$ & $6829{,}7618$ & $346{,}4893$ & $710{,}4084$\\
$\mathbf v_{2B}$ & $1964{,}3695$ & $2888{,}3181$ & $5921{,}9298$\\\bottomrule
\end{tabular}\par
\begin{align*}
\Delta\mathbf v_A&=\mathbf v_{2A}-\mathbf v_1
=(321{,}0033,\;-592{,}3412,\;-1214{,}4795)^{\mathsf T}\text{ м/с},\\
\Delta\mathbf v_B&=\mathbf v_{2B}-\mathbf v_1
=(-4544{,}3890,\;1949{,}4876,\;3997{,}0419)^{\mathsf T}\text{ м/с}.
\end{align*}
\[
\boxed{\Delta v_A=1388{,}8382\text{ м/с}},\qquad
\boxed{\Delta v_B=6358{,}3266\text{ м/с}}.
\]
\end{frame}

% Слайд 27
\begin{frame}{Выбор места коррекции $a$ и $e$}
\centering\begin{tikzpicture}
\begin{axis}[course plot,
width=128mm,height=48mm,ymin=0,ymax=7000,ytick={0,1000,2000,3000,4000,5000,6000,7000},
ylabel={$\Delta v$, м/с},
xlabel={$\nu_1$, град},
xmin=0,xmax=360,xtick={0,60,120,180,240,300,360},
unbounded coords=jump,filter discard warning=false,
legend style={at={(0.5,.98)},anchor=north,legend columns=2},
]
\addplot[courseaccent,solid,line width=.85pt,no marks] table[x=nu,y=dvA] {images/opt-ae.dat};
\addplot[coursemuted,densely dashed,line width=.85pt,no marks] table[x=nu,y=dvB] {images/opt-ae.dat};
\legend{A,B}
\end{axis}
\end{tikzpicture}
\par
\smallskip\raggedright
Те же исходные и целевые элементы; шаг перебора $\nu_1$ --- $0{,}5^\circ$.\newline
Для апсид: $a_2=(r_{\text{п},2}+r_{\text{а},2})/2$,\quad
$e_2=(r_{\text{а},2}-r_{\text{п},2})/(r_{\text{а},2}+r_{\text{п},2})$.
\end{frame}

% Коррекция a,omega и ограничения.
% Слайд 28
\begin{frame}{Коррекция $a$ и $\omega$}
Заданы $a_2,\omega_2$. Плоскость и положение точки не меняются, поэтому
\[
u=\omega_1+\nu_1,\qquad \nu_2=u-\omega_2\pmod{360^\circ}.
\]
В уравнении конечной орбиты неизвестен только $e_2$:
\[
r=\frac{a_2(1-e_2^2)}{1+e_2\cos\nu_2}.
\]
После умножения на знаменатель и переноса слагаемых получаем
\[
\boxed{a_2e_2^2+(r\cos\nu_2)e_2+r-a_2=0.}
\]
\justifying
Решаем квадратное уравнение и оставляем корни $0<e_2<1$.
Для каждого найденного эллипса вычисляем скорость в точке манёвра
и импульс перехода.
\end{frame}

% Слайд 29
\begin{frame}{Коррекция $a$ и $\omega$: расчёт эксцентриситета}
\begin{columns}[T,onlytextwidth]
\begin{column}{.48\textwidth}
\begin{block}{Исходное состояние}
$a_1=26\,222$ км,\quad $e_1=0{,}7$,\par
$i_1=63{,}434^\circ$,\quad $\Omega_1=0^\circ$,\par
$\omega_1=266^\circ$,\quad $\nu_1=222^\circ$.
\end{block}
\end{column}
\begin{column}{.48\textwidth}
\begin{block}{Требуется получить}
$a_2=26\,554$ км,\quad $\omega_2=270^\circ$.
\par\medskip
Найти конечные орбиты и импульсы.
\end{block}
\end{column}
\end{columns}
\[
r=\frac{26\,222(1-0{,}7^2)}{1+0{,}7\cos222^\circ}=27\,872{,}5686\text{ км},
\qquad \nu_2=218^\circ.
\]
Подставляем в квадратное уравнение:
\[
26\,554e_2^2-21\,963{,}8838e_2+1318{,}5686=0.
\]
Его корни: $e_{2A}=0{,}761972442$, $e_{2B}=0{,}065167867$.
\end{frame}

% Слайд 30
\begin{frame}{Коррекция $a$ и $\omega$: конечные орбиты}
\begin{columns}[T,onlytextwidth]
\begin{column}{.48\textwidth}
\centering
\renewcommand{\arraystretch}{1.07}
\begin{tabular}{@{}lrr@{}}
\toprule
Элемент & A & B\\\midrule
$a_2$, км & $26\,554$ & $26\,554$\\
$e_2$ & $0{,}761972$ & $0{,}065168$\\
$i_2$, град & $63{,}434$ & $63{,}434$\\
$\Omega_2$, град & $0$ & $0$\\
$\omega_2$, град & $270$ & $270$\\
$\nu_2$, град & $218$ & $218$\\
$r_{\text{п},2}$, км & $6320{,}6$ & $24\,823{,}5$\\\bottomrule
\end{tabular}\par\medskip
\justifying
Вариант A пересекает Землю: $r_{\text{п},2}<R_\oplus$.
\end{column}
\begin{column}{.50\textwidth}
\centering\begin{tikzpicture}[course diagram,x=1mm,y=1mm]
\path[use as bounding box] (-34.0, -24.5) rectangle (34.0, 24.5);
\filldraw[fill=coursetint,draw=courserule,line width=.4pt] (0.0000,-6.5943) circle[radius=3.8299];
\node[inner sep=0pt,font=\fontsize{8.5}{10}\selectfont] at (0.0000,-6.5943) {$\oplus$};
\draw[draw=courseink,line width=.75pt] plot coordinates {(-0.32950,-11.30642) (-0.20607,-11.31409) (-0.08249,-11.31987) (0.04127,-11.32375) (0.16523,-11.32572) (0.28939,-11.32578) (0.41379,-11.32393) (0.53844,-11.32015) (0.66337,-11.31443) (0.78860,-11.30676) (0.91414,-11.29712) (1.04002,-11.28550) (1.16625,-11.27188) (1.29287,-11.25623) (1.41989,-11.23854) (1.54733,-11.21877) (1.67521,-11.19690) (1.80356,-11.17290) (1.93240,-11.14673) (2.06174,-11.11836) (2.19161,-11.08775) (2.32203,-11.05486) (2.45302,-11.01965) (2.58460,-10.98207) (2.71679,-10.94207) (2.84962,-10.89960) (2.98310,-10.85461) (3.11726,-10.80703) (3.25211,-10.75681) (3.38768,-10.70388) (3.52399,-10.64817) (3.66104,-10.58962) (3.79887,-10.52814) (3.93749,-10.46365) (4.07692,-10.39608) (4.21717,-10.32533) (4.35826,-10.25131) (4.50022,-10.17392) (4.64304,-10.09307) (4.78674,-10.00864) (4.93135,-9.92053) (5.07685,-9.82861) (5.22328,-9.73276) (5.37063,-9.63285) (5.51891,-9.52875) (5.66811,-9.42031) (5.81826,-9.30739) (5.96933,-9.18983) (6.12133,-9.06747) (6.27425,-8.94014) (6.42808,-8.80766) (6.58281,-8.66984) (6.73840,-8.52650) (6.89485,-8.37742) (7.05213,-8.22240) (7.21018,-8.06122) (7.36899,-7.89364) (7.52850,-7.71943) (7.68864,-7.53834) (7.84937,-7.35010) (8.01061,-7.15445) (8.17228,-6.95110) (8.33428,-6.73977) (8.49651,-6.52014) (8.65885,-6.29192) (8.82117,-6.05477) (8.98333,-5.80835) (9.14516,-5.55233) (9.30647,-5.28635) (9.46708,-5.01005) (9.62675,-4.72304) (9.78524,-4.42496) (9.94228,-4.11540) (10.09757,-3.79399) (10.25078,-3.46031) (10.40156,-3.11398) (10.54951,-2.75459) (10.69419,-2.38174) (10.83516,-1.99504) (10.97188,-1.59412) (11.10381,-1.17860) (11.23034,-0.74815) (11.35082,-0.30243) (11.46455,0.15884) (11.57076,0.63592) (11.66864,1.12901) (11.75732,1.63827) (11.83585,2.16378) (11.90324,2.70554) (11.95843,3.26347) (12.00029,3.83740) (12.02765,4.42702) (12.03927,5.03189) (12.03385,5.65145) (12.01007,6.28493) (11.96654,6.93142) (11.90186,7.58980) (11.81462,8.25872) (11.70337,8.93661) (11.56672,9.62166) (11.40330,10.31180) (11.21179,11.00469) (10.99095,11.69773) (10.73968,12.38804) (10.45700,13.07247) (10.14210,13.74762) (9.79440,14.40987) (9.41354,15.05537) (8.99944,15.68010) (8.55232,16.27993) (8.07273,16.85062) (7.56157,17.38795) (7.02011,17.88773) (6.44997,18.34590) (5.85318,18.75860) (5.23209,19.12224) (4.58942,19.43361) (3.92820,19.68988) (3.25171,19.88875) (2.56348,20.02844) (1.86719,20.10775) (1.16664,20.12612) (0.46566,20.08357) (-0.23192,19.98079) (-0.92237,19.81903) (-1.60212,19.60014) (-2.26777,19.32647) (-2.91620,19.00084) (-3.54455,18.62648) (-4.15030,18.20692) (-4.73127,17.74596) (-5.28561,17.24757) (-5.81186,16.71581) (-6.30888,16.15479) (-6.77587,15.56858) (-7.21235,14.96116) (-7.61815,14.33640) (-7.99334,13.69798) (-8.33823,13.04936) (-8.65337,12.39380) (-8.93945,11.73430) (-9.19733,11.07362) (-9.42799,10.41425) (-9.63250,9.75844) (-9.81199,9.10818) (-9.96767,8.46521) (-10.10075,7.83107) (-10.21246,7.20706) (-10.30405,6.59429) (-10.37673,5.99369) (-10.43169,5.40600) (-10.47011,4.83183) (-10.49311,4.27164) (-10.50178,3.72578) (-10.49715,3.19445) (-10.48021,2.67781) (-10.45190,2.17590) (-10.41312,1.68868) (-10.36470,1.21607) (-10.30742,0.75792) (-10.24203,0.31404) (-10.16921,-0.11579) (-10.08962,-0.53184) (-10.00385,-0.93439) (-9.91245,-1.32375) (-9.81596,-1.70023) (-9.71486,-2.06418) (-9.60957,-2.41593) (-9.50053,-2.75583) (-9.38810,-3.08423) (-9.27263,-3.40147) (-9.15446,-3.70790) (-9.03387,-4.00387) (-8.91114,-4.28971) (-8.78652,-4.56576) (-8.66024,-4.83235) (-8.53250,-5.08978) (-8.40350,-5.33838) (-8.27342,-5.57844) (-8.14241,-5.81027) (-8.01061,-6.03414) (-7.87817,-6.25032) (-7.74521,-6.45910) (-7.61182,-6.66072) (-7.47811,-6.85543) (-7.34417,-7.04348) (-7.21008,-7.22509) (-7.07592,-7.40049) (-6.94174,-7.56989) (-6.80760,-7.73349) (-6.67357,-7.89150) (-6.53967,-8.04410) (-6.40596,-8.19148) (-6.27247,-8.33380) (-6.13922,-8.47124) (-6.00626,-8.60396) (-5.87359,-8.73210) (-5.74125,-8.85583) (-5.60924,-8.97527) (-5.47758,-9.09057) (-5.34628,-9.20185) (-5.21535,-9.30923) (-5.08479,-9.41284) (-4.95461,-9.51278) (-4.82481,-9.60917) (-4.69539,-9.70210) (-4.56635,-9.79168) (-4.43768,-9.87800) (-4.30939,-9.96115) (-4.18145,-10.04122) (-4.05388,-10.11828) (-3.92666,-10.19241) (-3.79978,-10.26369) (-3.67322,-10.33219) (-3.54699,-10.39798) (-3.42107,-10.46111) (-3.29545,-10.52165) (-3.17010,-10.57966) (-3.04503,-10.63519) (-2.92022,-10.68829) (-2.79565,-10.73901) (-2.67130,-10.78739) (-2.54717,-10.83349) (-2.42323,-10.87733) (-2.29947,-10.91897) (-2.17587,-10.95842) (-2.05243,-10.99573) (-1.92911,-11.03093) (-1.80590,-11.06405) (-1.68279,-11.09511) (-1.55975,-11.12413) (-1.43677,-11.15114) (-1.31383,-11.17617) (-1.19091,-11.19922) (-1.06800,-11.22031) (-0.94507,-11.23946) (-0.82210,-11.25668) (-0.69908,-11.27198) (-0.57599,-11.28536) (-0.45280,-11.29684) (-0.32950,-11.30642)};
\draw[draw=courseaccent,line width=1pt] plot coordinates {(0.00000,-10.38959) (0.09936,-10.38885) (0.19875,-10.38664) (0.29817,-10.38294) (0.39766,-10.37776) (0.49723,-10.37110) (0.59689,-10.36293) (0.69668,-10.35326) (0.79662,-10.34207) (0.89671,-10.32936) (0.99699,-10.31510) (1.09747,-10.29928) (1.19817,-10.28189) (1.29912,-10.26290) (1.40034,-10.24229) (1.50184,-10.22005) (1.60364,-10.19614) (1.70578,-10.17053) (1.80826,-10.14321) (1.91112,-10.11413) (2.01436,-10.08326) (2.11801,-10.05057) (2.22210,-10.01603) (2.32664,-9.97957) (2.43166,-9.94118) (2.53717,-9.90079) (2.64320,-9.85837) (2.74976,-9.81385) (2.85689,-9.76719) (2.96460,-9.71833) (3.07291,-9.66720) (3.18184,-9.61374) (3.29141,-9.55789) (3.40165,-9.49958) (3.51258,-9.43872) (3.62420,-9.37524) (3.73656,-9.30906) (3.84965,-9.24009) (3.96351,-9.16823) (4.07815,-9.09339) (4.19360,-9.01547) (4.30985,-8.93435) (4.42695,-8.84993) (4.54489,-8.76209) (4.66369,-8.67070) (4.78338,-8.57563) (4.90395,-8.47674) (5.02542,-8.37389) (5.14781,-8.26692) (5.27110,-8.15566) (5.39533,-8.03997) (5.52047,-7.91964) (5.64654,-7.79450) (5.77353,-7.66435) (5.90144,-7.52899) (6.03025,-7.38819) (6.15996,-7.24173) (6.29053,-7.08937) (6.42195,-6.93085) (6.55419,-6.76592) (6.68722,-6.59429) (6.82098,-6.41568) (6.95542,-6.22977) (7.09050,-6.03626) (7.22613,-5.83480) (7.36223,-5.62504) (7.49872,-5.40661) (7.63549,-5.17914) (7.77241,-4.94221) (7.90936,-4.69542) (8.04616,-4.43833) (8.18266,-4.17048) (8.31865,-3.89140) (8.45391,-3.60060) (8.58820,-3.29759) (8.72125,-2.98183) (8.85273,-2.65280) (8.98232,-2.30994) (9.10963,-1.95270) (9.23424,-1.58051) (9.35569,-1.19278) (9.47345,-0.78896) (9.58696,-0.36845) (9.69559,0.06930) (9.79866,0.52485) (9.89540,0.99871) (9.98500,1.49140) (10.06654,2.00335) (10.13906,2.53496) (10.20149,3.08656) (10.25267,3.65838) (10.29138,4.25056) (10.31627,4.86308) (10.32593,5.49582) (10.31885,6.14843) (10.29344,6.82038) (10.24803,7.51092) (10.18088,8.21898) (10.09020,8.94325) (9.97415,9.68204) (9.83089,10.43330) (9.65858,11.19461) (9.45545,11.96307) (9.21978,12.73537) (8.95000,13.50773) (8.64471,14.27588) (8.30275,15.03510) (7.92324,15.78024) (7.50565,16.50571) (7.04986,17.20562) (6.55620,17.87378) (6.02553,18.50384) (5.45924,19.08940) (4.85930,19.62414) (4.22827,20.10196) (3.56929,20.51716) (2.88604,20.86456) (2.18271,21.13971) (1.46392,21.33895) (0.73462,21.45960) (0.00000,21.50000) (-0.73462,21.45960) (-1.46392,21.33895) (-2.18271,21.13971) (-2.88604,20.86456) (-3.56929,20.51716) (-4.22827,20.10196) (-4.85930,19.62414) (-5.45924,19.08940) (-6.02553,18.50384) (-6.55620,17.87378) (-7.04986,17.20562) (-7.50565,16.50571) (-7.92324,15.78024) (-8.30275,15.03510) (-8.64471,14.27588) (-8.95000,13.50773) (-9.21978,12.73537) (-9.45545,11.96307) (-9.65858,11.19461) (-9.83089,10.43330) (-9.97415,9.68204) (-10.09020,8.94325) (-10.18088,8.21898) (-10.24803,7.51092) (-10.29344,6.82038) (-10.31885,6.14843) (-10.32593,5.49582) (-10.31627,4.86308) (-10.29138,4.25056) (-10.25267,3.65838) (-10.20149,3.08656) (-10.13906,2.53496) (-10.06654,2.00335) (-9.98500,1.49140) (-9.89540,0.99871) (-9.79866,0.52485) (-9.69559,0.06930) (-9.58696,-0.36845) (-9.47345,-0.78896) (-9.35569,-1.19278) (-9.23424,-1.58051) (-9.10963,-1.95270) (-8.98232,-2.30994) (-8.85273,-2.65280) (-8.72125,-2.98183) (-8.58820,-3.29759) (-8.45391,-3.60060) (-8.31865,-3.89140) (-8.18266,-4.17048) (-8.04616,-4.43833) (-7.90936,-4.69542) (-7.77241,-4.94221) (-7.63549,-5.17914) (-7.49872,-5.40661) (-7.36223,-5.62504) (-7.22613,-5.83480) (-7.09050,-6.03626) (-6.95542,-6.22977) (-6.82098,-6.41568) (-6.68722,-6.59429) (-6.55419,-6.76592) (-6.42195,-6.93085) (-6.29053,-7.08937) (-6.15996,-7.24173) (-6.03025,-7.38819) (-5.90144,-7.52899) (-5.77353,-7.66435) (-5.64654,-7.79450) (-5.52047,-7.91964) (-5.39533,-8.03997) (-5.27110,-8.15566) (-5.14781,-8.26692) (-5.02542,-8.37389) (-4.90395,-8.47674) (-4.78338,-8.57563) (-4.66369,-8.67070) (-4.54489,-8.76209) (-4.42695,-8.84993) (-4.30985,-8.93435) (-4.19360,-9.01547) (-4.07815,-9.09339) (-3.96351,-9.16823) (-3.84965,-9.24009) (-3.73656,-9.30906) (-3.62420,-9.37524) (-3.51258,-9.43872) (-3.40165,-9.49958) (-3.29141,-9.55789) (-3.18184,-9.61374) (-3.07291,-9.66720) (-2.96460,-9.71833) (-2.85689,-9.76719) (-2.74976,-9.81385) (-2.64320,-9.85837) (-2.53717,-9.90079) (-2.43166,-9.94118) (-2.32664,-9.97957) (-2.22210,-10.01603) (-2.11801,-10.05057) (-2.01436,-10.08326) (-1.91112,-10.11413) (-1.80826,-10.14321) (-1.70578,-10.17053) (-1.60364,-10.19614) (-1.50184,-10.22005) (-1.40034,-10.24229) (-1.29912,-10.26290) (-1.19817,-10.28189) (-1.09747,-10.29928) (-0.99699,-10.31510) (-0.89671,-10.32936) (-0.79662,-10.34207) (-0.69668,-10.35326) (-0.59689,-10.36293) (-0.49723,-10.37110) (-0.39766,-10.37776) (-0.29817,-10.38294) (-0.19875,-10.38664) (-0.09936,-10.38885) (-0.00000,-10.38959)};
\draw[draw=coursemuted,line width=.8pt,dashed] plot coordinates {(-0.00000,-21.50000) (0.39019,-21.49520) (0.78017,-21.48082) (1.16971,-21.45685) (1.55859,-21.42331) (1.94660,-21.38022) (2.33352,-21.32758) (2.71914,-21.26543) (3.10322,-21.19379) (3.48556,-21.11270) (3.86594,-21.02218) (4.24414,-20.92228) (4.61995,-20.81304) (4.99315,-20.69451) (5.36352,-20.56675) (5.73086,-20.42980) (6.09494,-20.28374) (6.45555,-20.12863) (6.81248,-19.96453) (7.16551,-19.79153) (7.51445,-19.60970) (7.85907,-19.41912) (8.19916,-19.21989) (8.53452,-19.01209) (8.86494,-18.79583) (9.19021,-18.57121) (9.51013,-18.33833) (9.82450,-18.09730) (10.13310,-17.84824) (10.43575,-17.59128) (10.73224,-17.32654) (11.02238,-17.05414) (11.30596,-16.77423) (11.58281,-16.48694) (11.85271,-16.19243) (12.11549,-15.89084) (12.37096,-15.58232) (12.61893,-15.26705) (12.85923,-14.94517) (13.09166,-14.61687) (13.31607,-14.28233) (13.53226,-13.94171) (13.74007,-13.59521) (13.93934,-13.24302) (14.12991,-12.88533) (14.31160,-12.52235) (14.48427,-12.15428) (14.64776,-11.78133) (14.80192,-11.40373) (14.94662,-11.02168) (15.08170,-10.63542) (15.20705,-10.24518) (15.32252,-9.85119) (15.42800,-9.45370) (15.52336,-9.05295) (15.60848,-8.64919) (15.68327,-8.24267) (15.74762,-7.83366) (15.80143,-7.42241) (15.84461,-7.00920) (15.87708,-6.59429) (15.89876,-6.17797) (15.90958,-5.76051) (15.90948,-5.34219) (15.89840,-4.92330) (15.87630,-4.50414) (15.84312,-4.08499) (15.79884,-3.66615) (15.74344,-3.24792) (15.67689,-2.83060) (15.59919,-2.41450) (15.51033,-1.99992) (15.41033,-1.58717) (15.29921,-1.17656) (15.17698,-0.76840) (15.04368,-0.36300) (14.89936,0.03933) (14.74407,0.43827) (14.57788,0.83351) (14.40085,1.22473) (14.21307,1.61163) (14.01464,1.99389) (13.80564,2.37120) (13.58621,2.74325) (13.35644,3.10973) (13.11649,3.47034) (12.86648,3.82478) (12.60658,4.17274) (12.33693,4.51393) (12.05772,4.84806) (11.76912,5.17483) (11.47132,5.49396) (11.16453,5.80517) (10.84895,6.10819) (10.52480,6.40275) (10.19231,6.68858) (9.85172,6.96543) (9.50327,7.23305) (9.14721,7.49118) (8.78382,7.73960) (8.41337,7.97809) (8.03612,8.20640) (7.65238,8.42435) (7.26243,8.63171) (6.86658,8.82831) (6.46514,9.01395) (6.05843,9.18845) (5.64676,9.35166) (5.23046,9.50342) (4.80988,9.64358) (4.38534,9.77201) (3.95719,9.88858) (3.52578,9.99319) (3.09146,10.08573) (2.65459,10.16611) (2.21552,10.23426) (1.77462,10.29011) (1.33225,10.33360) (0.88878,10.36470) (0.44458,10.38337) (0.00000,10.38959) (-0.44458,10.38337) (-0.88878,10.36470) (-1.33225,10.33360) (-1.77462,10.29011) (-2.21552,10.23426) (-2.65459,10.16611) (-3.09146,10.08573) (-3.52578,9.99319) (-3.95719,9.88858) (-4.38534,9.77201) (-4.80988,9.64358) (-5.23046,9.50342) (-5.64676,9.35166) (-6.05843,9.18845) (-6.46514,9.01395) (-6.86658,8.82831) (-7.26243,8.63171) (-7.65238,8.42435) (-8.03612,8.20640) (-8.41337,7.97809) (-8.78382,7.73960) (-9.14721,7.49118) (-9.50327,7.23305) (-9.85172,6.96543) (-10.19231,6.68858) (-10.52480,6.40275) (-10.84895,6.10819) (-11.16453,5.80517) (-11.47132,5.49396) (-11.76912,5.17483) (-12.05772,4.84806) (-12.33693,4.51393) (-12.60658,4.17274) (-12.86648,3.82478) (-13.11649,3.47034) (-13.35644,3.10973) (-13.58621,2.74325) (-13.80564,2.37120) (-14.01464,1.99389) (-14.21307,1.61163) (-14.40085,1.22473) (-14.57788,0.83351) (-14.74407,0.43827) (-14.89936,0.03933) (-15.04368,-0.36300) (-15.17698,-0.76840) (-15.29921,-1.17656) (-15.41033,-1.58717) (-15.51033,-1.99992) (-15.59919,-2.41450) (-15.67689,-2.83060) (-15.74344,-3.24792) (-15.79884,-3.66615) (-15.84312,-4.08499) (-15.87630,-4.50414) (-15.89840,-4.92330) (-15.90948,-5.34219) (-15.90958,-5.76051) (-15.89876,-6.17797) (-15.87708,-6.59429) (-15.84461,-7.00920) (-15.80143,-7.42241) (-15.74762,-7.83366) (-15.68327,-8.24267) (-15.60848,-8.64919) (-15.52336,-9.05295) (-15.42800,-9.45370) (-15.32252,-9.85119) (-15.20705,-10.24518) (-15.08170,-10.63542) (-14.94662,-11.02168) (-14.80192,-11.40373) (-14.64776,-11.78133) (-14.48427,-12.15428) (-14.31160,-12.52235) (-14.12991,-12.88533) (-13.93934,-13.24302) (-13.74007,-13.59521) (-13.53226,-13.94171) (-13.31607,-14.28233) (-13.09166,-14.61687) (-12.85923,-14.94517) (-12.61893,-15.26705) (-12.37096,-15.58232) (-12.11549,-15.89084) (-11.85271,-16.19243) (-11.58281,-16.48694) (-11.30596,-16.77423) (-11.02238,-17.05414) (-10.73224,-17.32654) (-10.43575,-17.59128) (-10.13310,-17.84824) (-9.82450,-18.09730) (-9.51013,-18.33833) (-9.19021,-18.57121) (-8.86494,-18.79583) (-8.53452,-19.01209) (-8.19916,-19.21989) (-7.85907,-19.41912) (-7.51445,-19.60970) (-7.16551,-19.79153) (-6.81248,-19.96453) (-6.45555,-20.12863) (-6.09494,-20.28374) (-5.73086,-20.42980) (-5.36352,-20.56675) (-4.99315,-20.69451) (-4.61995,-20.81304) (-4.24414,-20.92228) (-3.86594,-21.02218) (-3.48556,-21.11270) (-3.10322,-21.19379) (-2.71914,-21.26543) (-2.33352,-21.32758) (-1.94660,-21.38022) (-1.55859,-21.42331) (-1.16971,-21.45685) (-0.78017,-21.48082) (-0.39019,-21.49520) (-0.00000,-21.50000)};
\draw[aux line] (0.0000,-6.5943)--(-10.3041,6.5943);
\fill[courseink] (-10.3041,6.5943) circle[radius=.85];
\node[anchor=south west,inner sep=1mm] at (-10.3041,6.5943) {$P$};
\end{tikzpicture}
\par\figurelegend
{\usebeamerfont{caption}\color{coursemuted}
Земля показана в масштабе; $R_\oplus\approx6378$ км.\par}
\end{column}
\end{columns}
\end{frame}

% Слайд 31
\begin{frame}{Коррекция $a$ и $\omega$: скорости и импульсы}
При $p_{2A}=11\,136{,}6944$ км и $p_{2B}=26\,441{,}2291$ км:
\[
\mathbf r_1=\mathbf r_{2A}=\mathbf r_{2B}
=(-17\,160{,}0667,\;9822{,}8728,\;19\,644{,}9322)^{\mathsf T}\text{ км}.
\]
\centering
\renewcommand{\arraystretch}{1.04}
\begin{tabular}{@{}lrrr@{}}
\toprule
Скорость, м/с & $v_x$ & $v_y$ & $v_z$\\\midrule
$\mathbf v_1$ & $-489{,}8018$ & $-1622{,}4428$ & $-3244{,}7512$\\
$\mathbf v_{2A}$ & $-155{,}7771$ & $-1647{,}2598$ & $-3294{,}3832$\\
$\mathbf v_{2B}$ & $-2806{,}5448$ & $-1069{,}0535$ & $-2138{,}0184$\\\bottomrule
\end{tabular}\par
\begin{align*}
\Delta\mathbf v_A&=\mathbf v_{2A}-\mathbf v_1
=(334{,}0247,\;-24{,}8170,\;-49{,}6320)^{\mathsf T}\text{ м/с},\\
\Delta\mathbf v_B&=\mathbf v_{2B}-\mathbf v_1
=(-2316{,}7430,\;553{,}3893,\;1106{,}7327)^{\mathsf T}\text{ м/с}.
\end{align*}
\[
\boxed{\Delta v_A=338{,}6026\text{ м/с}},\qquad
\boxed{\Delta v_B=2626{,}4796\text{ м/с}}.
\]
\end{frame}

% Слайд 32
\begin{frame}{Выбор места коррекции $a$ и $\omega$}
\centering\begin{tikzpicture}
\begin{axis}[course plot,
width=128mm,height=48mm,ymin=0,ymax=3000,ytick={0,500,1000,1500,2000,2500,3000},
ylabel={$\Delta v$, м/с},
xlabel={$\nu_1$, град},
xmin=0,xmax=360,xtick={0,60,120,180,240,300,360},
unbounded coords=jump,filter discard warning=false,
legend style={at={(0.98,.98)},anchor=north east,legend columns=2},
]
\addplot[courseaccent,solid,line width=.85pt,no marks] table[x=nu,y=dvA] {images/opt-aw.dat};
\addplot[coursemuted,densely dashed,line width=.85pt,no marks] table[x=nu,y=dvB] {images/opt-aw.dat};
\legend{A,B}
\end{axis}
\end{tikzpicture}
\par
\smallskip\raggedright
A --- больший корень $e_2$, B --- меньший. Показан диапазон $\Delta v\le3000$ м/с.\newline
Для найденных минимумов проверяем высоту перицентра конечной орбиты.
\end{frame}

% Слайд 33
\begin{frame}{Когда коррекция невозможна}
\justifying
Конечная орбита должна проходить через текущую точку с радиусом $r$.
\begin{columns}[T,onlytextwidth]
\begin{column}{.48\textwidth}
\begin{block}{Заданы $a_2,e_2$ или радиусы апсид}
\centering
$r_{\text{п},2}\le r\le r_{\text{а},2}$.
\end{block}
\justifying
Нельзя поднять перицентр выше текущего радиуса или опустить апоцентр ниже него.
\par\medskip
В формуле для $\nu_2$ это проявляется как $|\cos\nu_2|>1$.
\end{column}
\begin{column}{.48\textwidth}
\begin{block}{Заданы $a_2,\omega_2$}
\justifying
У квадратного уравнения должен быть хотя бы один корень $0<e_2<1$.
\end{block}
\justifying
Нет действительных корней или все они вне этого диапазона --- требуемый эллипс не существует.
\end{column}
\end{columns}
\medskip
\justifying
Для найденного эллипса проверяем ограничения, в том числе высоту перицентра.
\end{frame}

% Изменение плоскости и пример a,e,Omega.
% Слайд 34
\begin{frame}{Плоскость орбиты после импульса}
\begin{columns}[T,onlytextwidth]
\begin{column}{.50\textwidth}
\justifying
Новая плоскость проходит через центр притяжения и точку манёвра.
\begin{block}{Геометрическое ограничение}
\centering
$\widehat{\mathbf h}_2\cdot\mathbf r=0$,\par\medskip
$\displaystyle\widehat{\mathbf h}_2=
\begin{pmatrix}
\sin i_2\sin\Omega_2\\
-\sin i_2\cos\Omega_2\\
\cos i_2
\end{pmatrix}$.
\end{block}
\justifying
В фиксированной точке $i_2$ и $\Omega_2$ связаны.
Произвольно задать оба угла нельзя.
\end{column}
\begin{column}{.48\textwidth}
\centering% Orthographic projection of two exact 3D planes; generated by plane_constraint_schematic().
\begin{tikzpicture}[course diagram,x=1mm,y=1mm]
\path[use as bounding box] (-34,-19) rectangle (34,35);
\fill[courseink!7!coursepaper] (-16.5466,5.2011)--(-7.8689,9.0751)--(25.2242,-1.3271)--(16.5466,-5.2011)--cycle;
\draw[courseink,line width=.65pt] (-16.5466,5.2011)--(-7.8689,9.0751)--(25.2242,-1.3271)--(16.5466,-5.2011);
\fill[courseaccent!13!coursepaper] (-16.5466,5.2011)--(-11.7124,15.6694)--(21.3807,5.2671)--(16.5466,-5.2011)--cycle;
\draw[courseaccent,line width=.65pt] (-16.5466,5.2011)--(-11.7124,15.6694)--(21.3807,5.2671)--(16.5466,-5.2011);
\fill[courseaccent!13!coursepaper] (-16.5466,5.2011)--(-21.3807,-5.2671)--(11.7124,-15.6694)--(16.5466,-5.2011)--cycle;
\draw[courseaccent,line width=.65pt] (-16.5466,5.2011)--(-21.3807,-5.2671)--(11.7124,-15.6694)--(16.5466,-5.2011);
\fill[courseink!7!coursepaper] (-16.5466,5.2011)--(-25.2242,1.3271)--(7.8689,-9.0751)--(16.5466,-5.2011)--cycle;
\draw[courseink,line width=.65pt] (-16.5466,5.2011)--(-25.2242,1.3271)--(7.8689,-9.0751)--(16.5466,-5.2011);
\draw[coursemuted,densely dashed,line width=.6pt] (-19.9938,6.2847)--(-16.5466,5.2011);
\draw[coursemuted,densely dashed,line width=.6pt] (16.5466,-5.2011)--(19.9938,-6.2847);
\draw[courseink,line width=.85pt] (-16.5466,5.2011)--(16.5466,-5.2011);
\draw[courseink,densely dashed,line width=.5pt] (0.0000,0.0000)--(4.0496,1.8079);
\draw[courseink,line width=.5pt] (0.0000,2.5034)--(1.7355,3.2782)--(1.7355,0.7748);
\draw[courseink,line width=.5pt] (0.0000,2.0862)--(1.7236,1.5444)--(1.7236,-0.5418);
\draw[courseink,line width=.95pt,->] (0.0000,0.0000)--(0.0000,29.2063) node[above,inner sep=1mm] {$\widehat{\mathbf h}_1$};
\draw[courseaccent,densely dashed,line width=.5pt] (-6.8944,2.1671)--(-4.2914,7.8039);
\draw[courseaccent,line width=.5pt] (-8.2239,3.1827)--(-7.1083,5.5985)--(-5.7788,4.5829);
\draw[courseaccent,line width=.5pt] (-8.0023,3.0135)--(-6.2787,2.4717)--(-5.1708,1.6253);
\draw[courseaccent,line width=.95pt,->] (-6.8944,2.1671)--(-22.4051,14.0160) node[above,inner sep=1mm] {$\widehat{\mathbf h}_2$};
\fill[courseaccent] (-6.8944,2.1671) circle[radius=.45];
\draw[vector line,line width=1pt] (0.0000,0.0000)--(13.0994,-4.1176) node[pos=.58,below=1mm] {$\mathbf r$};
\fill[courseink] (0.0000,0.0000) circle[radius=.7];
\node[below left,inner sep=1mm] at (0.0000,0.0000) {$O$};
\fill[courseink] (13.0994,-4.1176) circle[radius=.75];
\node[below right,inner sep=1mm] at (13.0994,-4.1176) {$P$};
\node[courseink,inner sep=.5mm] at (-0.0477,-5.2664) {1};
\node[courseaccent,inner sep=.5mm] at (12.6813,5.2352) {2};
\end{tikzpicture}
\par
{\usebeamerfont{caption}\color{coursemuted}
Общая линия плоскостей — $OP$.\par}
\end{column}
\end{columns}
\end{frame}

% Слайд 35
\begin{frame}{Коррекция наклонения или долготы узла}
Из условия $\widehat{\mathbf h}_2\cdot\mathbf r=0$ получаем
\[
r_x\sin i_2\sin\Omega_2-r_y\sin i_2\cos\Omega_2+r_z\cos i_2=0.
\]
\begin{columns}[T,onlytextwidth]
\begin{column}{.48\textwidth}
\begin{block}{Задана долгота узла $\Omega_2$}
\centering
$\displaystyle\tan i_2=\frac{-r_z}{r_x\sin\Omega_2-r_y\cos\Omega_2}$.
\end{block}
\raggedright
Находим $i_2$ в диапазоне $(0,180^\circ)$.
Затем рассчитываем коррекцию $a_2,e_2$ в найденной плоскости.
\end{column}
\begin{column}{.48\textwidth}
\begin{block}{Задано наклонение $i_2$}
\centering
$r_x\sin\Omega_2-r_y\cos\Omega_2=-r_z\cot i_2$.
\end{block}
\raggedright
Решаем тригонометрическое уравнение относительно $\Omega_2$.
Возможны две плоскости, одна или ни одной.
\end{column}
\end{columns}
\medskip
\raggedright
Точка вне экватора; для прямой орбиты $i_2\ge|\varphi|$, где $\varphi$ --- её широта.
\end{frame}

% Слайд 36
\begin{frame}{Коррекция $a$, $e$, $\Omega$: нахождение $i_2$}
Используем состояние из примера коррекции $a,e$:
\[
a_1=26\,000\text{ км},\quad e_1=0{,}6,\quad i_1=64^\circ,\quad
\Omega_1=0^\circ,\quad\omega_1=250^\circ,\quad\nu_1=60^\circ.
\]
Требуется $a_2=26\,554$ км, $e_2=0{,}7$, $\Omega_2=3^\circ$.
\[
\mathbf r=(8227{,}6814,\;-4298{,}3908,\;-8813{,}0072)^{\mathsf T}\text{ км}.
\]
Условие $\mathbf r\cdot\widehat{\mathbf h}_2=0$ даёт
\begin{align*}
A&=8227{,}6814\sin3^\circ-(-4298{,}3908)\cos3^\circ=4723{,}1036\text{ км},\\
\tan i_2&=\frac{-r_z}{A}=\frac{8813{,}0072}{4723{,}1036},
\qquad \boxed{i_2=61{,}812067^\circ.}
\end{align*}
\justifying
Наклонение изменилось вместе с долготой узла. Следующий шаг — определить
положение неизменного радиус-вектора в новой плоскости.
\end{frame}

% Слайд 37
\begin{frame}{Аргумент широты в новой плоскости}
\setlength{\abovedisplayskip}{2pt}
\setlength{\belowdisplayskip}{2pt}
\setlength{\jot}{0pt}
\justifying
Единичный вектор узла --- $\widehat{\mathbf n}_2$;
угол от него до $\mathbf r$ в плоскости орбиты --- $u_2$.
\[
\widehat{\mathbf n}_2=(\cos\Omega_2,\;\sin\Omega_2,\;0)^{\mathsf T},\qquad
\widehat{\mathbf m}_2=\widehat{\mathbf h}_2\times\widehat{\mathbf n}_2.
\]
$\widehat{\mathbf m}_2$ направлен на $90^\circ$ вперёд от $\widehat{\mathbf n}_2$ по движению:
\begin{gather*}
\widehat{\mathbf m}_2=(-\sin\Omega_2\cos i_2,\;\cos\Omega_2\cos i_2,\;\sin i_2)^{\mathsf T},\\
\mathbf r=r\bigl(\cos u_2\,\widehat{\mathbf n}_2+\sin u_2\,\widehat{\mathbf m}_2\bigr).
\end{gather*}
\begin{columns}[T,onlytextwidth]
\begin{column}{.49\textwidth}
{\color{courseaccent}Проекция на направление узла\par}
\smallskip\centering
$\widehat{\mathbf n}_2\cdot\mathbf r=r\cos u_2$\par\smallskip
$\widehat{\mathbf n}_2\cdot\mathbf r=r_x\cos\Omega_2+r_y\sin\Omega_2$\par\smallskip
$\displaystyle\boxed{\cos u_2=\frac{r_x\cos\Omega_2+r_y\sin\Omega_2}{r}.}$
\end{column}
\begin{column}{.49\textwidth}
{\color{courseaccent}Третья координата радиус-вектора\par}
\smallskip\centering
$r_z=r\bigl(0\cdot\cos u_2+\sin i_2\sin u_2\bigr)$\par\smallskip
$r_z=r\sin i_2\sin u_2$\par\smallskip
$\displaystyle\boxed{\sin u_2=\frac{r_z}{r\sin i_2}.}$
\end{column}
\end{columns}
\end{frame}

% Слайд 38
\begin{frame}{Коррекция $a$, $e$, $\Omega$: остальные углы}
\justifying
При $r=12\,800$ км, $\Omega_2=3^\circ$ и $i_2=61{,}8121^\circ$:
\begin{align*}
\cos u_2&=\frac{8227{,}6814\cos3^\circ-4298{,}3908\sin3^\circ}{12\,800}
\approx0{,}6243317,\\[2pt]
\sin u_2&=\frac{-8813{,}0072}{12\,800\sin61{,}8121^\circ}\approx-0{,}7811594.
\end{align*}
Косинус положителен, синус отрицателен: четвёртая четверть,
$\boxed{u_2=308{,}6332^\circ}$.
\par\smallskip
Целевые $a_2,e_2$ и радиус те же, что в плоском примере:
\[
\begin{array}{c|cc}
 & \text{A} & \text{B}\\\hline
\nu_2 & 85{,}2463^\circ & 274{,}7537^\circ\\
\omega_2=u_2-\nu_2 & 223{,}3869^\circ & 33{,}8794^\circ
\end{array}
\]
\justifying
Общие элементы: $a_2=26\,554$ км, $e_2=0{,}7$,
$i_2=61{,}8121^\circ$, $\Omega_2=3^\circ$.
\end{frame}

% Слайд 39
\begin{frame}{Коррекция $a$, $e$, $\Omega$: конечные орбиты}
\begin{columns}[T,onlytextwidth]
\begin{column}{.48\textwidth}
\centering
\renewcommand{\arraystretch}{1.14}
\begin{tabular}{@{}lrr@{}}
\toprule
Элемент & A & B\\\midrule
$a_2$, км & $26\,554$ & $26\,554$\\
$e_2$ & $0{,}7$ & $0{,}7$\\
$i_2$, град & $61{,}8121$ & $61{,}8121$\\
$\Omega_2$, град & $3$ & $3$\\
$\omega_2$, град & $223{,}3869$ & $33{,}8794$\\
$\nu_2$, град & $85{,}2463$ & $274{,}7537$\\\bottomrule
\end{tabular}\par\medskip
\justifying
Оба варианта лежат в одной конечной плоскости.
Угол между ней и исходной плоскостью --- $3{,}4523^\circ$.
\end{column}
\begin{column}{.50\textwidth}
\centering% Actual orbit coordinates, uniform scale, orthographic camera; generated by node_orbit_picture().
\begin{tikzpicture}[course diagram,x=1mm,y=1mm]
\path[use as bounding box] (-34.0,-23.5) rectangle (34.0,23.5);
\draw[courseink,line width=.75pt] plot coordinates {(-10.35273,0.19492) (-10.37099,0.16143) (-10.38887,0.12791) (-10.40638,0.09434) (-10.42351,0.06073) (-10.44027,0.02708) (-10.45665,-0.00661) (-10.47266,-0.04035) (-10.48829,-0.07413) (-10.50355,-0.10796) (-10.51843,-0.14183) (-10.53294,-0.17575) (-10.54708,-0.20971) (-10.56083,-0.24371) (-10.57422,-0.27776) (-10.58722,-0.31186) (-10.59985,-0.34601) (-10.61210,-0.38020) (-10.62398,-0.41444) (-10.63548,-0.44873) (-10.64660,-0.48306) (-10.65734,-0.51744) (-10.66770,-0.55188) (-10.67768,-0.58636) (-10.68728,-0.62089) (-10.69650,-0.65547) (-10.70533,-0.69010) (-10.71379,-0.72479) (-10.72186,-0.75952) (-10.72954,-0.79431) (-10.73684,-0.82914) (-10.74376,-0.86403) (-10.74615,-0.87684)};
\draw[courseink,line width=.75pt] plot coordinates {(28.43672,0.07558) (28.36879,0.15603) (28.25660,0.28219) (28.13971,0.40731) (28.01820,0.53133) (27.89218,0.65423) (27.76173,0.77598) (27.62697,0.89653) (27.48799,1.01586) (27.34489,1.13394) (27.19778,1.25074) (27.04676,1.36623) (26.89193,1.48039) (26.73341,1.59320) (26.57128,1.70462) (26.40567,1.81464) (26.23667,1.92323) (26.06440,2.03039) (25.88896,2.13608) (25.71045,2.24030) (25.52899,2.34303) (25.34467,2.44426) (25.15761,2.54397) (24.96790,2.64215) (24.77565,2.73879) (24.58097,2.83389) (24.38395,2.92743) (24.18471,3.01941) (23.98333,3.10982) (23.77991,3.19867) (23.57457,3.28594) (23.36739,3.37163) (23.15846,3.45575) (22.94789,3.53829) (22.73577,3.61925) (22.52218,3.69864) (22.30723,3.77645) (22.09099,3.85269) (21.87355,3.92737) (21.65501,4.00049) (21.43544,4.07206) (21.21492,4.14207) (20.99355,4.21055) (20.77139,4.27749) (20.54852,4.34290) (20.32503,4.40680) (20.10098,4.46920) (19.87644,4.53009) (19.65149,4.58951) (19.42619,4.64744) (19.20062,4.70392) (18.97483,4.75894) (18.74890,4.81253) (18.52288,4.86469) (18.29683,4.91544) (18.07081,4.96479) (17.84488,5.01275) (17.61909,5.05935) (17.39350,5.10459) (17.16816,5.14849) (16.94312,5.19106) (16.71842,5.23232) (16.49412,5.27228) (16.27026,5.31096) (16.04688,5.34838) (15.82402,5.38455) (15.60173,5.41949) (15.38005,5.45321) (15.15901,5.48573) (14.93865,5.51706) (14.71900,5.54722) (14.50010,5.57623) (14.28198,5.60410) (14.06468,5.63086) (13.84821,5.65651) (13.63262,5.68107) (13.41792,5.70455) (13.20414,5.72699) (12.99131,5.74838) (12.77945,5.76874) (12.56859,5.78810) (12.35874,5.80647) (12.14992,5.82386) (11.94216,5.84029) (11.73548,5.85578) (11.52988,5.87033) (11.32539,5.88397) (11.12202,5.89672) (10.91979,5.90857) (10.71871,5.91956) (10.51879,5.92970) (10.32004,5.93899) (10.12248,5.94746) (9.92612,5.95512) (9.73096,5.96198) (9.53701,5.96806) (9.34429,5.97337) (9.15280,5.97793) (8.96254,5.98174) (8.77353,5.98483) (8.58576,5.98721) (8.39925,5.98888) (8.21399,5.98986) (8.03000,5.99017) (7.84727,5.98982) (7.66580,5.98882) (7.48561,5.98717) (7.30668,5.98491) (7.12903,5.98202) (6.95265,5.97854) (6.77754,5.97446) (6.60371,5.96981) (6.43115,5.96459) (6.25986,5.95881) (6.08984,5.95248) (5.92108,5.94562) (5.75360,5.93823) (5.58738,5.93033) (5.42242,5.92193) (5.25872,5.91303) (5.09627,5.90364) (4.93508,5.89378) (4.77514,5.88346) (4.61644,5.87268) (4.45898,5.86145) (4.30276,5.84978) (4.14777,5.83769) (3.99401,5.82518) (3.84146,5.81225) (3.69014,5.79892) (3.54002,5.78520) (3.39112,5.77109) (3.24341,5.75660) (3.09689,5.74174) (2.95157,5.72652) (2.80742,5.71095) (2.66446,5.69502) (2.52266,5.67875) (2.38203,5.66215) (2.24255,5.64523) (2.10422,5.62798) (1.96704,5.61042) (1.83099,5.59256) (1.69607,5.57439) (1.56228,5.55593) (1.42960,5.53718) (1.29803,5.51815) (1.16756,5.49885) (1.03819,5.47928) (0.90990,5.45944) (0.78270,5.43934) (0.65656,5.41900) (0.53150,5.39840) (0.40749,5.37756) (0.28453,5.35649) (0.16262,5.33518) (0.04174,5.31365) (-0.07811,5.29190) (-0.19694,5.26993) (-0.31475,5.24774) (-0.43156,5.22535) (-0.54736,5.20276) (-0.66218,5.17997) (-0.77601,5.15698) (-0.88886,5.13380) (-1.00074,5.11044) (-1.11166,5.08689) (-1.22163,5.06316) (-1.33064,5.03927) (-1.43872,5.01520) (-1.54586,4.99096) (-1.65207,4.96656) (-1.75737,4.94200) (-1.86175,4.91728) (-1.96523,4.89242) (-2.06781,4.86740) (-2.16949,4.84223) (-2.27030,4.81692) (-2.37023,4.79148) (-2.46928,4.76589) (-2.56748,4.74017) (-2.66481,4.71432) (-2.76130,4.68834) (-2.85694,4.66224) (-2.95175,4.63601) (-3.04573,4.60966) (-3.13889,4.58320) (-3.23123,4.55662) (-3.32275,4.52993) (-3.41348,4.50312) (-3.50341,4.47621) (-3.59255,4.44919) (-3.68090,4.42207) (-3.76848,4.39485) (-3.85528,4.36753) (-3.94132,4.34011) (-4.02660,4.31260) (-4.11112,4.28499) (-4.19490,4.25730) (-4.27794,4.22951) (-4.36024,4.20164) (-4.44181,4.17368) (-4.52266,4.14564) (-4.60278,4.11751) (-4.68220,4.08931) (-4.76091,4.06102) (-4.83892,4.03266) (-4.91623,4.00423) (-4.99285,3.97572) (-5.06879,3.94713) (-5.14405,3.91848) (-5.21863,3.88976) (-5.29255,3.86096) (-5.36580,3.83211) (-5.43839,3.80318) (-5.51033,3.77419) (-5.58162,3.74514) (-5.65226,3.71602) (-5.72227,3.68685) (-5.79165,3.65761) (-5.86039,3.62832) (-5.92851,3.59896) (-5.99602,3.56956) (-6.06290,3.54009) (-6.12918,3.51057) (-6.19485,3.48100) (-6.25993,3.45137) (-6.32440,3.42169) (-6.38828,3.39197) (-6.45158,3.36219) (-6.51429,3.33236) (-6.57642,3.30248) (-6.63798,3.27256) (-6.69896,3.24259) (-6.75938,3.21257) (-6.81924,3.18250) (-6.87854,3.15240) (-6.93728,3.12224) (-6.99547,3.09205) (-7.05312,3.06181) (-7.11022,3.03152) (-7.16679,3.00120) (-7.22281,2.97083) (-7.27831,2.94043) (-7.33327,2.90998) (-7.38771,2.87949) (-7.44163,2.84897) (-7.49504,2.81840) (-7.54792,2.78780) (-7.60030,2.75716) (-7.65217,2.72648) (-7.70353,2.69576) (-7.75439,2.66500) (-7.80475,2.63421) (-7.85462,2.60339) (-7.90400,2.57252) (-7.95289,2.54163) (-8.00129,2.51069) (-8.04920,2.47972) (-8.09664,2.44872) (-8.14360,2.41768) (-8.19009,2.38661) (-8.23610,2.35550) (-8.28165,2.32436) (-8.32673,2.29319) (-8.37135,2.26198) (-8.41550,2.23074) (-8.45920,2.19946) (-8.50244,2.16815) (-8.54522,2.13681) (-8.58756,2.10544) (-8.62945,2.07403) (-8.67089,2.04259) (-8.71189,2.01112) (-8.75244,1.97962) (-8.79256,1.94808) (-8.83224,1.91651) (-8.87148,1.88491) (-8.91029,1.85327) (-8.94867,1.82161) (-8.98662,1.78991) (-9.02414,1.75817) (-9.06124,1.72641) (-9.09792,1.69461) (-9.13417,1.66279) (-9.17001,1.63092) (-9.20543,1.59903) (-9.24043,1.56710) (-9.27502,1.53515) (-9.30919,1.50315) (-9.34296,1.47113) (-9.37631,1.43907) (-9.40926,1.40698) (-9.44180,1.37486) (-9.47394,1.34270) (-9.50567,1.31051) (-9.53700,1.27829) (-9.56794,1.24603) (-9.59847,1.21374) (-9.62860,1.18142) (-9.65834,1.14906) (-9.68769,1.11666) (-9.71664,1.08424) (-9.74519,1.05177) (-9.77336,1.01928) (-9.80113,0.98675) (-9.82852,0.95418) (-9.85551,0.92158) (-9.88212,0.88894) (-9.90835,0.85627) (-9.93418,0.82356) (-9.95963,0.79081) (-9.98470,0.75803) (-10.00939,0.72521) (-10.03369,0.69236) (-10.05761,0.65946) (-10.08115,0.62653) (-10.10431,0.59357) (-10.12708,0.56056) (-10.14948,0.52752) (-10.17151,0.49443) (-10.19315,0.46131) (-10.21441,0.42815) (-10.23530,0.39495) (-10.25581,0.36172) (-10.27595,0.32844) (-10.29571,0.29512) (-10.31509,0.26176) (-10.33410,0.22836) (-10.35273,0.19492)};
\draw[courseaccent,line width=1pt] plot coordinates {(-7.24134,1.50522) (-7.27648,1.48424) (-7.31141,1.46315) (-7.34613,1.44197) (-7.38063,1.42069) (-7.41492,1.39931) (-7.44900,1.37783) (-7.48287,1.35625) (-7.51653,1.33457) (-7.54998,1.31278) (-7.58323,1.29090) (-7.61626,1.26892) (-7.64909,1.24683) (-7.68170,1.22464) (-7.71411,1.20234) (-7.74632,1.17994) (-7.77831,1.15744) (-7.81010,1.13483) (-7.84169,1.11212) (-7.87306,1.08930) (-7.90423,1.06637) (-7.93520,1.04334) (-7.96596,1.02020) (-7.99651,0.99695) (-8.02686,0.97359) (-8.05700,0.95012) (-8.08694,0.92654) (-8.11667,0.90284) (-8.14620,0.87904) (-8.17552,0.85512) (-8.20463,0.83110) (-8.23354,0.80695) (-8.26225,0.78270) (-8.29074,0.75832) (-8.31903,0.73383) (-8.34712,0.70923) (-8.37500,0.68450) (-8.40267,0.65966) (-8.43013,0.63470) (-8.45739,0.60962) (-8.48444,0.58442) (-8.51128,0.55910) (-8.53791,0.53365) (-8.56433,0.50808) (-8.59054,0.48239) (-8.61655,0.45657) (-8.64234,0.43063) (-8.66792,0.40456) (-8.69329,0.37836) (-8.71844,0.35203) (-8.74339,0.32558) (-8.76812,0.29899) (-8.79263,0.27227) (-8.81693,0.24542) (-8.84102,0.21844) (-8.86488,0.19132) (-8.88853,0.16407) (-8.91196,0.13668) (-8.93517,0.10915) (-8.95817,0.08148) (-8.98094,0.05368) (-9.00348,0.02573) (-9.02581,-0.00235) (-9.04791,-0.03058) (-9.06978,-0.05896) (-9.09142,-0.08748) (-9.11284,-0.11614) (-9.13403,-0.14496) (-9.15499,-0.17392) (-9.17572,-0.20303) (-9.19621,-0.23229) (-9.21647,-0.26170) (-9.23649,-0.29127) (-9.25628,-0.32099) (-9.27582,-0.35087) (-9.29513,-0.38090) (-9.31419,-0.41110) (-9.33301,-0.44145) (-9.35158,-0.47197) (-9.36990,-0.50264) (-9.37061,-0.50385)};
\draw[courseaccent,line width=1pt] plot coordinates {(20.62296,-1.25106) (20.43779,-1.14686) (20.24450,-1.04011) (20.05078,-0.93509) (19.85673,-0.83180) (19.66240,-0.73023) (19.46788,-0.63037) (19.27322,-0.53220) (19.07849,-0.43572) (18.88376,-0.34091) (18.68908,-0.24775) (18.49451,-0.15624) (18.30010,-0.06636) (18.10591,0.02190) (17.91199,0.10857) (17.71839,0.19365) (17.52516,0.27717) (17.33233,0.35913) (17.13996,0.43956) (16.94807,0.51847) (16.75673,0.59588) (16.56595,0.67180) (16.37577,0.74626) (16.18624,0.81927) (15.99737,0.89085) (15.80921,0.96102) (15.62178,1.02980) (15.43511,1.09719) (15.24922,1.16323) (15.06414,1.22793) (14.87989,1.29131) (14.69649,1.35339) (14.51397,1.41418) (14.33235,1.47370) (14.15164,1.53197) (13.97185,1.58901) (13.79302,1.64484) (13.61514,1.69948) (13.43823,1.75293) (13.26232,1.80523) (13.08740,1.85638) (12.91349,1.90641) (12.74060,1.95533) (12.56875,2.00316) (12.39792,2.04991) (12.22815,2.09561) (12.05942,2.14027) (11.89176,2.18391) (11.72516,2.22653) (11.55962,2.26817) (11.39516,2.30883) (11.23177,2.34853) (11.06946,2.38729) (10.90824,2.42512) (10.74809,2.46204) (10.58903,2.49805) (10.43106,2.53319) (10.27417,2.56745) (10.11836,2.60086) (9.96364,2.63343) (9.81000,2.66518) (9.65745,2.69611) (9.50597,2.72625) (9.35557,2.75560) (9.20625,2.78417) (9.05800,2.81199) (8.91082,2.83907) (8.76470,2.86541) (8.61965,2.89103) (8.47566,2.91594) (8.33272,2.94015) (8.19083,2.96368) (8.04999,2.98654) (7.91019,3.00874) (7.77143,3.03028) (7.63369,3.05119) (7.49698,3.07147) (7.36129,3.09113) (7.22662,3.11019) (7.09295,3.12865) (6.96029,3.14652) (6.82862,3.16382) (6.69794,3.18055) (6.56825,3.19673) (6.43953,3.21236) (6.31179,3.22745) (6.18501,3.24202) (6.05919,3.25606) (5.93432,3.26960) (5.81039,3.28263) (5.68740,3.29517) (5.56535,3.30723) (5.44421,3.31881) (5.32400,3.32992) (5.20469,3.34057) (5.08629,3.35077) (4.96879,3.36053) (4.85217,3.36984) (4.73644,3.37873) (4.62158,3.38720) (4.50759,3.39525) (4.39446,3.40289) (4.28219,3.41013) (4.17076,3.41698) (4.06017,3.42343) (3.95042,3.42951) (3.84149,3.43521) (3.73338,3.44054) (3.62608,3.44551) (3.51958,3.45012) (3.41389,3.45438) (3.30898,3.45830) (3.20486,3.46187) (3.10151,3.46511) (2.99894,3.46802) (2.89713,3.47061) (2.79607,3.47288) (2.69577,3.47484) (2.59621,3.47649) (2.49738,3.47784) (2.39929,3.47889) (2.30191,3.47964) (2.20526,3.48011) (2.10931,3.48029) (2.01407,3.48020) (1.91953,3.47982) (1.82567,3.47918) (1.73250,3.47827) (1.64001,3.47710) (1.54819,3.47567) (1.45704,3.47398) (1.36654,3.47205) (1.27670,3.46986) (1.18751,3.46744) (1.09895,3.46477) (1.01104,3.46187) (0.92375,3.45874) (0.83709,3.45537) (0.75104,3.45179) (0.66560,3.44798) (0.58078,3.44395) (0.49655,3.43970) (0.41292,3.43525) (0.32987,3.43058) (0.24741,3.42571) (0.16553,3.42063) (0.08423,3.41535) (0.00349,3.40988) (-0.07669,3.40420) (-0.15631,3.39834) (-0.23538,3.39229) (-0.31390,3.38605) (-0.39188,3.37963) (-0.46932,3.37302) (-0.54623,3.36623) (-0.62262,3.35927) (-0.69848,3.35213) (-0.77382,3.34482) (-0.84865,3.33734) (-0.92298,3.32969) (-0.99680,3.32187) (-1.07012,3.31389) (-1.14295,3.30575) (-1.21530,3.29745) (-1.28715,3.28899) (-1.35853,3.28038) (-1.42943,3.27161) (-1.49986,3.26269) (-1.56982,3.25361) (-1.63932,3.24439) (-1.70836,3.23503) (-1.77694,3.22551) (-1.84508,3.21586) (-1.91277,3.20606) (-1.98001,3.19612) (-2.04682,3.18604) (-2.11320,3.17583) (-2.17914,3.16548) (-2.24466,3.15499) (-2.30975,3.14438) (-2.37443,3.13363) (-2.43869,3.12275) (-2.50253,3.11174) (-2.56597,3.10060) (-2.62901,3.08934) (-2.69164,3.07796) (-2.75388,3.06644) (-2.81572,3.05481) (-2.87717,3.04306) (-2.93824,3.03118) (-2.99892,3.01918) (-3.05922,3.00707) (-3.11915,2.99484) (-3.17869,2.98249) (-3.23787,2.97003) (-3.29668,2.95746) (-3.35513,2.94477) (-3.41321,2.93196) (-3.47094,2.91905) (-3.52831,2.90602) (-3.58533,2.89289) (-3.64200,2.87965) (-3.69832,2.86629) (-3.75430,2.85284) (-3.80993,2.83927) (-3.86523,2.82560) (-3.92019,2.81182) (-3.97482,2.79794) (-4.02912,2.78395) (-4.08309,2.76986) (-4.13674,2.75567) (-4.19006,2.74138) (-4.24307,2.72698) (-4.29575,2.71248) (-4.34812,2.69789) (-4.40018,2.68319) (-4.45193,2.66839) (-4.50337,2.65350) (-4.55451,2.63851) (-4.60534,2.62341) (-4.65587,2.60823) (-4.70610,2.59294) (-4.75604,2.57756) (-4.80568,2.56208) (-4.85503,2.54650) (-4.90409,2.53083) (-4.95287,2.51507) (-5.00135,2.49921) (-5.04956,2.48325) (-5.09748,2.46720) (-5.14512,2.45106) (-5.19248,2.43482) (-5.23957,2.41849) (-5.28638,2.40207) (-5.33292,2.38555) (-5.37919,2.36894) (-5.42520,2.35224) (-5.47093,2.33545) (-5.51640,2.31856) (-5.56161,2.30158) (-5.60655,2.28451) (-5.65124,2.26735) (-5.69567,2.25009) (-5.73984,2.23274) (-5.78375,2.21531) (-5.82741,2.19778) (-5.87082,2.18015) (-5.91398,2.16244) (-5.95689,2.14464) (-5.99955,2.12674) (-6.04197,2.10875) (-6.08414,2.09067) (-6.12607,2.07250) (-6.16775,2.05424) (-6.20920,2.03589) (-6.25040,2.01744) (-6.29137,1.99891) (-6.33210,1.98028) (-6.37259,1.96156) (-6.41285,1.94275) (-6.45288,1.92384) (-6.49267,1.90484) (-6.53223,1.88575) (-6.57157,1.86657) (-6.61067,1.84730) (-6.64955,1.82793) (-6.68820,1.80847) (-6.72662,1.78892) (-6.76482,1.76927) (-6.80279,1.74953) (-6.84055,1.72969) (-6.87808,1.70976) (-6.91538,1.68974) (-6.95247,1.66962) (-6.98934,1.64941) (-7.02599,1.62910) (-7.06242,1.60869) (-7.09864,1.58819) (-7.13464,1.56759) (-7.17042,1.54690) (-7.20599,1.52611) (-7.24134,1.50522)};
\draw[coursemuted,line width=.8pt,dashed] plot coordinates {(4.80989,-0.85712) (4.83234,-0.82642) (4.85585,-0.79375) (4.87912,-0.76086) (4.90215,-0.72776) (4.92495,-0.69443) (4.94750,-0.66088) (4.96981,-0.62710) (4.99188,-0.59310) (5.01369,-0.55886) (5.03526,-0.52440) (5.05657,-0.48970) (5.07762,-0.45476) (5.09840,-0.41958) (5.11893,-0.38417) (5.13919,-0.34851) (5.15918,-0.31261) (5.17889,-0.27646) (5.19833,-0.24006) (5.21749,-0.20341) (5.23636,-0.16650) (5.25495,-0.12934) (5.27324,-0.09192) (5.29124,-0.05423) (5.30895,-0.01629) (5.32635,0.02193) (5.34344,0.06041) (5.36022,0.09916) (5.37669,0.13819) (5.39284,0.17749) (5.40867,0.21707) (5.42417,0.25694) (5.43934,0.29709) (5.45417,0.33752) (5.46866,0.37825) (5.48280,0.41927) (5.49659,0.46058) (5.51002,0.50219) (5.52310,0.54410) (5.53580,0.58632) (5.54814,0.62884) (5.56010,0.67168) (5.57167,0.71482) (5.58286,0.75829) (5.59365,0.80207) (5.60404,0.84617) (5.61402,0.89060) (5.62359,0.93535) (5.63274,0.98044) (5.64147,1.02586) (5.64976,1.07161) (5.65762,1.11771) (5.66502,1.16415) (5.67198,1.21094) (5.67848,1.25808) (5.68450,1.30558) (5.69006,1.35343) (5.69513,1.40164) (5.69971,1.45021) (5.70379,1.49915) (5.70737,1.54847) (5.71043,1.59815) (5.71297,1.64822) (5.71498,1.69866) (5.71645,1.74949) (5.71737,1.80071) (5.71773,1.85232) (5.71752,1.90433) (5.71674,1.95674) (5.71537,2.00955) (5.71341,2.06276) (5.71083,2.11639) (5.70764,2.17043) (5.70383,2.22489) (5.69937,2.27977) (5.69427,2.33508) (5.68850,2.39082) (5.68207,2.44698) (5.67495,2.50359) (5.66714,2.56064) (5.65862,2.61813) (5.64938,2.67607) (5.63941,2.73446) (5.62870,2.79331) (5.61723,2.85263) (5.60499,2.91240) (5.59197,2.97264) (5.57815,3.03336) (5.56351,3.09455) (5.54805,3.15622) (5.53175,3.21837) (5.51460,3.28102) (5.49658,3.34415) (5.47767,3.40778) (5.45786,3.47191) (5.43713,3.53654) (5.41546,3.60168) (5.39285,3.66733) (5.36927,3.73350) (5.34471,3.80018) (5.31915,3.86738) (5.29256,3.93511) (5.26494,4.00337) (5.23626,4.07217) (5.20651,4.14150) (5.17566,4.21137) (5.14369,4.28178) (5.11060,4.35274) (5.07635,4.42425) (5.04093,4.49631) (5.00431,4.56893) (4.96647,4.64211) (4.92740,4.71585) (4.88706,4.79016) (4.84544,4.86503) (4.80252,4.94048) (4.75827,5.01650) (4.71267,5.09309) (4.66568,5.17027) (4.61730,5.24802) (4.56750,5.32636) (4.51624,5.40528) (4.46351,5.48478) (4.40928,5.56488) (4.35352,5.64556) (4.29620,5.72683) (4.23730,5.80870) (4.17679,5.89116) (4.11465,5.97421) (4.05084,6.05785) (3.98534,6.14209) (3.91811,6.22691) (3.84913,6.31234) (3.77837,6.39835) (3.70580,6.48496) (3.63138,6.57215) (3.55509,6.65994) (3.47690,6.74831) (3.39677,6.83726) (3.31466,6.92680) (3.23056,7.01691) (3.14442,7.10760) (3.05622,7.19887) (2.96592,7.29069) (2.87348,7.38309) (2.77887,7.47604) (2.68206,7.56954) (2.58302,7.66359) (2.48170,7.75818) (2.37808,7.85330) (2.27211,7.94895) (2.16377,8.04511) (2.05302,8.14178) (1.93981,8.23895) (1.82413,8.33661) (1.70592,8.43475) (1.58516,8.53335) (1.46180,8.63241) (1.33582,8.73191) (1.20718,8.83184) (1.07583,8.93218) (0.94176,9.03292) (0.80491,9.13404) (0.66525,9.23553) (0.52276,9.33737) (0.37738,9.43953) (0.22910,9.54200) (0.07787,9.64477) (-0.07633,9.74780) (-0.23355,9.85107) (-0.39381,9.95457) (-0.55714,10.05826) (-0.72358,10.16212) (-0.89316,10.26613) (-1.06591,10.37025) (-1.24185,10.47446) (-1.42102,10.57872) (-1.60344,10.68301) (-1.78913,10.78729) (-1.97812,10.89152) (-2.17044,10.99568) (-2.36610,11.09972) (-2.56513,11.20361) (-2.76754,11.30731) (-2.97334,11.41077) (-3.18257,11.51395) (-3.39522,11.61682) (-3.61131,11.71932) (-3.83085,11.82140) (-4.05384,11.92303) (-4.28030,12.02415) (-4.51021,12.12471) (-4.74358,12.22465) (-4.98041,12.32393) (-5.22069,12.42249) (-5.46441,12.52026) (-5.71156,12.61721) (-5.96211,12.71325) (-6.21607,12.80834) (-6.47339,12.90240) (-6.73406,12.99539) (-6.99804,13.08722) (-7.26531,13.17784) (-7.53582,13.26717) (-7.80953,13.35515) (-8.08640,13.44171) (-8.36637,13.52678) (-8.64940,13.61028) (-8.93541,13.69214) (-9.22436,13.77229) (-9.51617,13.85065) (-9.81077,13.92714) (-10.10807,14.00169) (-10.40801,14.07423) (-10.71048,14.14466) (-11.01539,14.21293) (-11.32266,14.27894) (-11.63216,14.34261) (-11.94381,14.40388) (-12.25747,14.46266) (-12.57304,14.51887) (-12.89038,14.57244) (-13.20938,14.62328) (-13.52989,14.67132) (-13.85177,14.71648) (-14.17489,14.75870) (-14.49908,14.79788) (-14.82421,14.83396) (-15.15010,14.86688) (-15.47659,14.89655) (-15.80352,14.92291) (-16.13071,14.94589) (-16.45799,14.96543) (-16.78518,14.98147) (-17.11209,14.99394) (-17.43853,15.00280) (-17.76432,15.00797) (-18.08925,15.00942) (-18.41314,15.00709) (-18.73578,15.00094) (-19.05698,14.99093) (-19.37653,14.97701) (-19.69422,14.95916) (-20.00985,14.93734) (-20.32322,14.91153) (-20.63411,14.88169) (-20.94233,14.84782) (-21.24766,14.80990) (-21.54991,14.76791) (-21.84885,14.72186) (-22.14431,14.67174) (-22.43607,14.61755) (-22.72393,14.55931) (-23.00770,14.49702) (-23.28719,14.43070) (-23.56222,14.36038) (-23.83258,14.28607) (-24.09811,14.20782) (-24.35863,14.12566) (-24.61397,14.03963) (-24.86396,13.94977) (-25.10845,13.85614) (-25.34727,13.75878) (-25.58029,13.65776) (-25.80737,13.55314) (-26.02836,13.44499) (-26.24315,13.33337) (-26.45161,13.21836) (-26.65363,13.10003) (-26.84911,12.97848) (-27.03796,12.85378) (-27.22009,12.72602) (-27.39540,12.59530) (-27.56385,12.46170) (-27.72535,12.32532) (-27.87987,12.18627) (-28.02734,12.04463) (-28.16774,11.90052) (-28.30102,11.75404) (-28.42718,11.60529) (-28.54620,11.45438) (-28.65807,11.30141) (-28.76279,11.14651) (-28.86037,10.98976) (-28.95084,10.83130) (-29.03422,10.67122) (-29.11053,10.50964) (-29.17983,10.34666) (-29.24215,10.18240) (-29.29755,10.01696) (-29.34609,9.85046) (-29.38784,9.68299) (-29.42286,9.51467) (-29.45124,9.34561) (-29.47305,9.17589) (-29.48839,9.00563) (-29.49734,8.83493) (-29.50000,8.66389) (-29.49648,8.49260) (-29.48687,8.32115) (-29.47129,8.14964) (-29.44985,7.97817) (-29.42266,7.80681) (-29.38983,7.63565) (-29.35150,7.46478) (-29.30778,7.29428) (-29.25879,7.12422) (-29.20467,6.95469) (-29.14553,6.78575) (-29.08151,6.61748) (-29.01273,6.44993) (-28.93932,6.28319) (-28.86142,6.11731) (-28.77915,5.95235) (-28.69265,5.78837) (-28.60204,5.62542) (-28.50746,5.46356) (-28.40903,5.30283) (-28.30688,5.14328) (-28.20114,4.98497) (-28.09194,4.82791) (-27.97940,4.67217) (-27.86365,4.51777) (-27.74481,4.36475) (-27.62299,4.21315) (-27.49833,4.06299) (-27.37094,3.91429) (-27.24093,3.76710) (-27.10842,3.62142) (-26.97352,3.47729) (-26.83634,3.33472) (-26.69699,3.19373) (-26.55558,3.05434) (-26.41220,2.91655) (-26.26697,2.78039) (-26.11998,2.64587) (-25.97133,2.51299) (-25.82111,2.38175) (-25.66942,2.25218) (-25.51634,2.12427) (-25.36197,1.99802) (-25.20639,1.87344) (-25.04968,1.75053) (-24.89193,1.62929) (-24.73321,1.50972) (-24.57360,1.39181) (-24.41318,1.27557) (-24.25201,1.16097) (-24.09017,1.04804) (-23.92773,0.93674) (-23.76474,0.82708) (-23.60128,0.71906) (-23.43741,0.61265) (-23.27318,0.50786) (-23.10865,0.40467) (-22.94388,0.30307) (-22.77892,0.20305) (-22.61383,0.10460) (-22.44864,0.00771) (-22.28342,-0.08765) (-22.11820,-0.18147) (-22.10919,-0.18650)};
\filldraw[fill=coursetint,draw=courserule,line width=.45pt] (-2.4718,-0.6757) circle[radius=5.0051];
\draw[courseink,line width=.75pt] plot coordinates {(-10.74615,-0.87684) (-10.75028,-0.89898) (-10.75642,-0.93397) (-10.76217,-0.96902) (-10.76753,-1.00413) (-10.77250,-1.03928) (-10.77707,-1.07450) (-10.78125,-1.10977) (-10.78504,-1.14509) (-10.78843,-1.18047) (-10.79142,-1.21591) (-10.79401,-1.25140) (-10.79620,-1.28695) (-10.79799,-1.32256) (-10.79938,-1.35823) (-10.80036,-1.39395) (-10.80094,-1.42974) (-10.80111,-1.46558) (-10.80086,-1.50148) (-10.80021,-1.53745) (-10.79914,-1.57347) (-10.79766,-1.60956) (-10.79576,-1.64571) (-10.79344,-1.68192) (-10.79070,-1.71819) (-10.78754,-1.75452) (-10.78395,-1.79092) (-10.77994,-1.82738) (-10.77550,-1.86391) (-10.77062,-1.90049) (-10.76532,-1.93715) (-10.75957,-1.97387) (-10.75339,-2.01065) (-10.74677,-2.04750) (-10.73971,-2.08442) (-10.73220,-2.12141) (-10.72424,-2.15846) (-10.71583,-2.19558) (-10.70697,-2.23276) (-10.69765,-2.27002) (-10.68788,-2.30734) (-10.67764,-2.34474) (-10.66694,-2.38220) (-10.65577,-2.41973) (-10.64414,-2.45734) (-10.63203,-2.49501) (-10.61944,-2.53276) (-10.60637,-2.57057) (-10.59283,-2.60846) (-10.57879,-2.64642) (-10.56427,-2.68445) (-10.54926,-2.72256) (-10.53375,-2.76073) (-10.51774,-2.79899) (-10.50123,-2.83731) (-10.48421,-2.87571) (-10.46668,-2.91418) (-10.44864,-2.95273) (-10.43008,-2.99136) (-10.41101,-3.03005) (-10.39140,-3.06883) (-10.37127,-3.10768) (-10.35060,-3.14660) (-10.32940,-3.18561) (-10.30765,-3.22468) (-10.28536,-3.26384) (-10.26252,-3.30307) (-10.23913,-3.34238) (-10.21517,-3.38177) (-10.19066,-3.42123) (-10.16557,-3.46078) (-10.13992,-3.50040) (-10.11368,-3.54010) (-10.08687,-3.57987) (-10.05946,-3.61973) (-10.03147,-3.65966) (-10.00288,-3.69967) (-9.97368,-3.73976) (-9.94388,-3.77994) (-9.91347,-3.82018) (-9.88244,-3.86051) (-9.85079,-3.90092) (-9.81850,-3.94140) (-9.78559,-3.98197) (-9.75203,-4.02261) (-9.71783,-4.06334) (-9.68298,-4.10414) (-9.64747,-4.14502) (-9.61129,-4.18598) (-9.57445,-4.22702) (-9.53693,-4.26813) (-9.49873,-4.30933) (-9.45985,-4.35060) (-9.42026,-4.39195) (-9.37998,-4.43338) (-9.33899,-4.47489) (-9.29729,-4.51647) (-9.25486,-4.55813) (-9.21171,-4.59987) (-9.16782,-4.64169) (-9.12319,-4.68358) (-9.07781,-4.72554) (-9.03167,-4.76758) (-8.98477,-4.80970) (-8.93710,-4.85188) (-8.88865,-4.89415) (-8.83942,-4.93648) (-8.78939,-4.97889) (-8.73856,-5.02137) (-8.68692,-5.06392) (-8.63446,-5.10654) (-8.58117,-5.14923) (-8.52706,-5.19200) (-8.47210,-5.23482) (-8.41629,-5.27772) (-8.35961,-5.32068) (-8.30208,-5.36371) (-8.24366,-5.40680) (-8.18436,-5.44995) (-8.12417,-5.49317) (-8.06307,-5.53645) (-8.00105,-5.57979) (-7.93812,-5.62318) (-7.87426,-5.66664) (-7.80945,-5.71015) (-7.74369,-5.75371) (-7.67698,-5.79733) (-7.60929,-5.84100) (-7.54062,-5.88472) (-7.47097,-5.92849) (-7.40031,-5.97230) (-7.32865,-6.01616) (-7.25596,-6.06006) (-7.18224,-6.10401) (-7.10748,-6.14799) (-7.03167,-6.19201) (-6.95479,-6.23607) (-6.87684,-6.28016) (-6.79781,-6.32427) (-6.71768,-6.36842) (-6.63644,-6.41259) (-6.55409,-6.45679) (-6.47060,-6.50101) (-6.38598,-6.54524) (-6.30020,-6.58949) (-6.21325,-6.63376) (-6.12513,-6.67803) (-6.03583,-6.72231) (-5.94532,-6.76659) (-5.85360,-6.81087) (-5.76065,-6.85515) (-5.66647,-6.89942) (-5.57104,-6.94368) (-5.47435,-6.98792) (-5.37638,-7.03215) (-5.27713,-7.07635) (-5.17657,-7.12053) (-5.07471,-7.16468) (-4.97152,-7.20880) (-4.86700,-7.25287) (-4.76112,-7.29691) (-4.65388,-7.34089) (-4.54526,-7.38482) (-4.43526,-7.42869) (-4.32385,-7.47250) (-4.21102,-7.51624) (-4.09677,-7.55991) (-3.98107,-7.60350) (-3.86392,-7.64700) (-3.74530,-7.69041) (-3.62519,-7.73372) (-3.50359,-7.77693) (-3.38049,-7.82003) (-3.25585,-7.86302) (-3.12969,-7.90588) (-3.00197,-7.94861) (-2.87269,-7.99121) (-2.74184,-8.03366) (-2.60940,-8.07595) (-2.47535,-8.11809) (-2.33969,-8.16006) (-2.20241,-8.20186) (-2.06348,-8.24347) (-1.92290,-8.28489) (-1.78065,-8.32611) (-1.63672,-8.36712) (-1.49110,-8.40791) (-1.34377,-8.44847) (-1.19473,-8.48880) (-1.04396,-8.52887) (-0.89145,-8.56869) (-0.73718,-8.60824) (-0.58116,-8.64751) (-0.42335,-8.68650) (-0.26376,-8.72518) (-0.10238,-8.76355) (0.06082,-8.80159) (0.22583,-8.83930) (0.39267,-8.87667) (0.56135,-8.91367) (0.73187,-8.95030) (0.90425,-8.98655) (1.07849,-9.02239) (1.25461,-9.05783) (1.43261,-9.09284) (1.61250,-9.12741) (1.79428,-9.16153) (1.97796,-9.19518) (2.16355,-9.22834) (2.35105,-9.26101) (2.54046,-9.29316) (2.73180,-9.32478) (2.92507,-9.35586) (3.12026,-9.38637) (3.31738,-9.41630) (3.51643,-9.44564) (3.71741,-9.47437) (3.92032,-9.50246) (4.12517,-9.52991) (4.33195,-9.55669) (4.54065,-9.58278) (4.75127,-9.60817) (4.96382,-9.63283) (5.17828,-9.65676) (5.39465,-9.67992) (5.61291,-9.70230) (5.83307,-9.72387) (6.05512,-9.74462) (6.27903,-9.76453) (6.50481,-9.78358) (6.73243,-9.80174) (6.96188,-9.81899) (7.19316,-9.83531) (7.42623,-9.85068) (7.66109,-9.86507) (7.89771,-9.87847) (8.13607,-9.89085) (8.37615,-9.90219) (8.61793,-9.91246) (8.86138,-9.92165) (9.10646,-9.92972) (9.35317,-9.93666) (9.60145,-9.94243) (9.85129,-9.94703) (10.10264,-9.95042) (10.35548,-9.95257) (10.60975,-9.95347) (10.86544,-9.95309) (11.12248,-9.95141) (11.38084,-9.94839) (11.64048,-9.94403) (11.90134,-9.93829) (12.16338,-9.93114) (12.42655,-9.92257) (12.69079,-9.91256) (12.95605,-9.90107) (13.22227,-9.88808) (13.48938,-9.87357) (13.75734,-9.85752) (14.02607,-9.83990) (14.29550,-9.82069) (14.56558,-9.79987) (14.83622,-9.77741) (15.10736,-9.75330) (15.37892,-9.72751) (15.65083,-9.70002) (15.92299,-9.67081) (16.19534,-9.63986) (16.46779,-9.60715) (16.74025,-9.57267) (17.01263,-9.53638) (17.28484,-9.49829) (17.55680,-9.45836) (17.82840,-9.41659) (18.09955,-9.37296) (18.37016,-9.32745) (18.64011,-9.28005) (18.90931,-9.23074) (19.17766,-9.17953) (19.44504,-9.12638) (19.71135,-9.07130) (19.97649,-9.01428) (20.24033,-8.95530) (20.50277,-8.89437) (20.76370,-8.83147) (21.02299,-8.76661) (21.28054,-8.69977) (21.53622,-8.63096) (21.78992,-8.56018) (22.04151,-8.48743) (22.29088,-8.41271) (22.53791,-8.33602) (22.78248,-8.25737) (23.02445,-8.17677) (23.26372,-8.09422) (23.50016,-8.00974) (23.73364,-7.92334) (23.96405,-7.83502) (24.19127,-7.74481) (24.41516,-7.65271) (24.63562,-7.55875) (24.85253,-7.46295) (25.06575,-7.36533) (25.27519,-7.26590) (25.48071,-7.16470) (25.68221,-7.06174) (25.87956,-6.95706) (26.07267,-6.85069) (26.26142,-6.74265) (26.44570,-6.63297) (26.62541,-6.52170) (26.80044,-6.40886) (26.97069,-6.29450) (27.13606,-6.17864) (27.29647,-6.06133) (27.45181,-5.94261) (27.60201,-5.82253) (27.74696,-5.70111) (27.88659,-5.57842) (28.02083,-5.45450) (28.14959,-5.32938) (28.27281,-5.20313) (28.39041,-5.07579) (28.50234,-4.94741) (28.60853,-4.81804) (28.70894,-4.68774) (28.80350,-4.55655) (28.89217,-4.42454) (28.97491,-4.29175) (29.05168,-4.15824) (29.12246,-4.02407) (29.18720,-3.88929) (29.24589,-3.75397) (29.29850,-3.61815) (29.34503,-3.48189) (29.38546,-3.34526) (29.41979,-3.20830) (29.44801,-3.07109) (29.47013,-2.93367) (29.48616,-2.79610) (29.49611,-2.65845) (29.50000,-2.52076) (29.49785,-2.38309) (29.48969,-2.24551) (29.47554,-2.10807) (29.45545,-1.97082) (29.42945,-1.83382) (29.39758,-1.69713) (29.35990,-1.56079) (29.31645,-1.42487) (29.26729,-1.28940) (29.21248,-1.15446) (29.15208,-1.02008) (29.08616,-0.88631) (29.01477,-0.75321) (28.93801,-0.62082) (28.85594,-0.48919) (28.76863,-0.35837) (28.67618,-0.22839) (28.57867,-0.09930) (28.47617,0.02885) (28.43672,0.07558)};
\draw[courseaccent,line width=1pt] plot coordinates {(-9.37061,-0.50385) (-9.38798,-0.53348) (-9.40580,-0.56449) (-9.42338,-0.59566) (-9.44069,-0.62700) (-9.45775,-0.65851) (-9.47455,-0.69019) (-9.49109,-0.72204) (-9.50737,-0.75407) (-9.52338,-0.78627) (-9.53913,-0.81864) (-9.55460,-0.85120) (-9.56980,-0.88394) (-9.58473,-0.91685) (-9.59938,-0.94995) (-9.61375,-0.98324) (-9.62784,-1.01671) (-9.64164,-1.05036) (-9.65516,-1.08421) (-9.66839,-1.11825) (-9.68132,-1.15248) (-9.69396,-1.18690) (-9.70630,-1.22152) (-9.71834,-1.25634) (-9.73007,-1.29136) (-9.74149,-1.32658) (-9.75261,-1.36200) (-9.76341,-1.39763) (-9.77389,-1.43346) (-9.78405,-1.46950) (-9.79389,-1.50576) (-9.80340,-1.54222) (-9.81257,-1.57890) (-9.82141,-1.61579) (-9.82992,-1.65291) (-9.83808,-1.69024) (-9.84589,-1.72779) (-9.85335,-1.76557) (-9.86046,-1.80357) (-9.86720,-1.84180) (-9.87359,-1.88027) (-9.87960,-1.91896) (-9.88525,-1.95789) (-9.89051,-1.99705) (-9.89540,-2.03645) (-9.89990,-2.07609) (-9.90401,-2.11598) (-9.90772,-2.15611) (-9.91103,-2.19648) (-9.91394,-2.23711) (-9.91644,-2.27799) (-9.91852,-2.31912) (-9.92017,-2.36050) (-9.92141,-2.40215) (-9.92220,-2.44406) (-9.92256,-2.48622) (-9.92248,-2.52866) (-9.92194,-2.57136) (-9.92095,-2.61433) (-9.91950,-2.65757) (-9.91758,-2.70109) (-9.91518,-2.74489) (-9.91230,-2.78897) (-9.90893,-2.83332) (-9.90507,-2.87797) (-9.90070,-2.92290) (-9.89582,-2.96812) (-9.89043,-3.01363) (-9.88452,-3.05944) (-9.87807,-3.10554) (-9.87109,-3.15195) (-9.86355,-3.19866) (-9.85547,-3.24567) (-9.84682,-3.29300) (-9.83760,-3.34063) (-9.82780,-3.38858) (-9.81741,-3.43684) (-9.80643,-3.48542) (-9.79484,-3.53433) (-9.78263,-3.58355) (-9.76981,-3.63311) (-9.75635,-3.68299) (-9.74224,-3.73321) (-9.72749,-3.78377) (-9.71207,-3.83466) (-9.69598,-3.88589) (-9.67921,-3.93747) (-9.66175,-3.98939) (-9.64358,-4.04167) (-9.62470,-4.09429) (-9.60509,-4.14727) (-9.58475,-4.20061) (-9.56365,-4.25431) (-9.54180,-4.30838) (-9.51918,-4.36281) (-9.49577,-4.41761) (-9.47157,-4.47278) (-9.44655,-4.52832) (-9.42072,-4.58425) (-9.39405,-4.64056) (-9.36653,-4.69724) (-9.33816,-4.75432) (-9.30890,-4.81178) (-9.27876,-4.86964) (-9.24772,-4.92789) (-9.21575,-4.98654) (-9.18286,-5.04559) (-9.14902,-5.10504) (-9.11421,-5.16490) (-9.07843,-5.22516) (-9.04165,-5.28584) (-9.00387,-5.34692) (-8.96506,-5.40843) (-8.92520,-5.47035) (-8.88429,-5.53269) (-8.84229,-5.59546) (-8.79921,-5.65865) (-8.75501,-5.72227) (-8.70968,-5.78631) (-8.66321,-5.85079) (-8.61556,-5.91571) (-8.56673,-5.98106) (-8.51670,-6.04685) (-8.46544,-6.11307) (-8.41293,-6.17974) (-8.35916,-6.24686) (-8.30411,-6.31442) (-8.24774,-6.38242) (-8.19005,-6.45088) (-8.13101,-6.51978) (-8.07060,-6.58913) (-8.00880,-6.65894) (-7.94558,-6.72920) (-7.88092,-6.79991) (-7.81480,-6.87108) (-7.74719,-6.94270) (-7.67808,-7.01478) (-7.60743,-7.08731) (-7.53522,-7.16030) (-7.46143,-7.23375) (-7.38603,-7.30765) (-7.30900,-7.38200) (-7.23031,-7.45681) (-7.14993,-7.53207) (-7.06784,-7.60779) (-6.98401,-7.68395) (-6.89842,-7.76057) (-6.81103,-7.83763) (-6.72182,-7.91514) (-6.63076,-7.99309) (-6.53782,-8.07148) (-6.44298,-8.15031) (-6.34619,-8.22958) (-6.24745,-8.30927) (-6.14671,-8.38940) (-6.04394,-8.46994) (-5.93912,-8.55091) (-5.83221,-8.63229) (-5.72319,-8.71407) (-5.61202,-8.79626) (-5.49867,-8.87885) (-5.38312,-8.96182) (-5.26532,-9.04518) (-5.14525,-9.12892) (-5.02288,-9.21302) (-4.89818,-9.29748) (-4.77110,-9.38229) (-4.64163,-9.46744) (-4.50973,-9.55292) (-4.37536,-9.63873) (-4.23850,-9.72484) (-4.09911,-9.81125) (-3.95716,-9.89794) (-3.81262,-9.98491) (-3.66546,-10.07213) (-3.51564,-10.15960) (-3.36314,-10.24729) (-3.20792,-10.33520) (-3.04995,-10.42331) (-2.88921,-10.51159) (-2.72565,-10.60004) (-2.55926,-10.68862) (-2.39000,-10.77733) (-2.21785,-10.86614) (-2.04277,-10.95503) (-1.86473,-11.04398) (-1.68372,-11.13296) (-1.49971,-11.22195) (-1.31266,-11.31092) (-1.12256,-11.39986) (-0.92938,-11.48872) (-0.73310,-11.57749) (-0.53370,-11.66613) (-0.33115,-11.75462) (-0.12545,-11.84291) (0.08343,-11.93099) (0.29551,-12.01881) (0.51080,-12.10634) (0.72931,-12.19354) (0.95105,-12.28038) (1.17603,-12.36682) (1.40427,-12.45282) (1.63576,-12.53834) (1.87051,-12.62333) (2.10852,-12.70776) (2.34979,-12.79158) (2.59431,-12.87474) (2.84208,-12.95719) (3.09310,-13.03889) (3.34734,-13.11979) (3.60480,-13.19984) (3.86545,-13.27898) (4.12928,-13.35717) (4.39626,-13.43434) (4.66637,-13.51045) (4.93957,-13.58543) (5.21584,-13.65923) (5.49512,-13.73178) (5.77739,-13.80304) (6.06258,-13.87294) (6.35067,-13.94141) (6.64158,-14.00839) (6.93526,-14.07382) (7.23166,-14.13764) (7.53069,-14.19977) (7.83230,-14.26016) (8.13639,-14.31873) (8.44290,-14.37542) (8.75174,-14.43015) (9.06281,-14.48287) (9.37602,-14.53350) (9.69127,-14.58197) (10.00845,-14.62821) (10.32745,-14.67215) (10.64816,-14.71373) (10.97045,-14.75287) (11.29419,-14.78951) (11.61927,-14.82357) (11.94553,-14.85499) (12.27284,-14.88370) (12.60105,-14.90963) (12.93001,-14.93271) (13.25956,-14.95289) (13.58954,-14.97010) (13.91979,-14.98427) (14.25014,-14.99535) (14.58040,-15.00327) (14.91041,-15.00798) (15.23998,-15.00942) (15.56892,-15.00754) (15.89704,-15.00229) (16.22415,-14.99361) (16.55005,-14.98147) (16.87455,-14.96581) (17.19743,-14.94661) (17.51850,-14.92381) (17.83755,-14.89738) (18.15437,-14.86730) (18.46875,-14.83354) (18.78048,-14.79607) (19.08934,-14.75486) (19.39514,-14.70991) (19.69765,-14.66120) (19.99667,-14.60872) (20.29199,-14.55247) (20.58339,-14.49244) (20.87068,-14.42864) (21.15364,-14.36107) (21.43207,-14.28975) (21.70579,-14.21469) (21.97458,-14.13592) (22.23826,-14.05345) (22.49664,-13.96733) (22.74955,-13.87757) (22.99679,-13.78423) (23.23820,-13.68733) (23.47361,-13.58694) (23.70287,-13.48309) (23.92580,-13.37585) (24.14228,-13.26527) (24.35216,-13.15141) (24.55530,-13.03435) (24.75159,-12.91415) (24.94089,-12.79089) (25.12311,-12.66465) (25.29814,-12.53549) (25.46588,-12.40352) (25.62626,-12.26882) (25.77920,-12.13147) (25.92463,-11.99156) (26.06249,-11.84921) (26.19274,-11.70449) (26.31533,-11.55750) (26.43023,-11.40836) (26.53742,-11.25716) (26.63689,-11.10400) (26.72863,-10.94900) (26.81264,-10.79225) (26.88894,-10.63386) (26.95755,-10.47394) (27.01850,-10.31260) (27.07183,-10.14994) (27.11757,-9.98608) (27.15578,-9.82111) (27.18653,-9.65515) (27.20988,-9.48831) (27.22589,-9.32068) (27.23467,-9.15238) (27.23628,-8.98350) (27.23082,-8.81416) (27.21840,-8.64444) (27.19912,-8.47446) (27.17307,-8.30430) (27.14039,-8.13406) (27.10118,-7.96384) (27.05557,-7.79372) (27.00369,-7.62380) (26.94566,-7.45417) (26.88162,-7.28491) (26.81171,-7.11610) (26.73605,-6.94782) (26.65480,-6.78015) (26.56809,-6.61317) (26.47607,-6.44694) (26.37888,-6.28154) (26.27667,-6.11703) (26.16959,-5.95348) (26.05779,-5.79095) (25.94140,-5.62950) (25.82059,-5.46919) (25.69550,-5.31006) (25.56627,-5.15218) (25.43306,-4.99558) (25.29600,-4.84032) (25.15525,-4.68643) (25.01095,-4.53396) (24.86323,-4.38295) (24.71225,-4.23342) (24.55814,-4.08543) (24.40104,-3.93898) (24.24108,-3.79412) (24.07840,-3.65088) (23.91313,-3.50926) (23.74539,-3.36930) (23.57533,-3.23102) (23.40305,-3.09443) (23.22868,-2.95955) (23.05235,-2.82639) (22.87417,-2.69497) (22.69425,-2.56529) (22.51270,-2.43737) (22.32964,-2.31120) (22.14518,-2.18681) (21.95940,-2.06418) (21.77243,-1.94332) (21.58434,-1.82424) (21.39525,-1.70693) (21.20524,-1.59139) (21.01440,-1.47762) (20.82283,-1.36561) (20.63060,-1.25536) (20.62296,-1.25106)};
\draw[coursemuted,line width=.8pt,dashed] plot coordinates {(1.56499,-3.22298) (1.60654,-3.20542) (1.64791,-3.18774) (1.68910,-3.16995) (1.73011,-3.15204) (1.77094,-3.13402) (1.81159,-3.11588) (1.85207,-3.09763) (1.89237,-3.07925) (1.93249,-3.06077) (1.97244,-3.04216) (2.01222,-3.02344) (2.05182,-3.00460) (2.09125,-2.98563) (2.13051,-2.96655) (2.16959,-2.94736) (2.20850,-2.92804) (2.24724,-2.90859) (2.28581,-2.88903) (2.32421,-2.86935) (2.36244,-2.84954) (2.40050,-2.82961) (2.43839,-2.80956) (2.47610,-2.78938) (2.51365,-2.76908) (2.55103,-2.74866) (2.58825,-2.72810) (2.62529,-2.70742) (2.66217,-2.68662) (2.69887,-2.66568) (2.73541,-2.64462) (2.77178,-2.62343) (2.80799,-2.60211) (2.84402,-2.58065) (2.87989,-2.55907) (2.91559,-2.53735) (2.95113,-2.51550) (2.98649,-2.49352) (3.02169,-2.47140) (3.05672,-2.44914) (3.09159,-2.42675) (3.12628,-2.40423) (3.16081,-2.38156) (3.19517,-2.35876) (3.22936,-2.33582) (3.26338,-2.31273) (3.29723,-2.28951) (3.33092,-2.26614) (3.36443,-2.24263) (3.39777,-2.21897) (3.43095,-2.19517) (3.46395,-2.17123) (3.49679,-2.14713) (3.52945,-2.12289) (3.56194,-2.09850) (3.59425,-2.07396) (3.62640,-2.04926) (3.65837,-2.02442) (3.69017,-1.99942) (3.72179,-1.97427) (3.75324,-1.94896) (3.78451,-1.92349) (3.81560,-1.89787) (3.84652,-1.87208) (3.87726,-1.84614) (3.90782,-1.82003) (3.93820,-1.79376) (3.96840,-1.76733) (3.99842,-1.74073) (4.02826,-1.71397) (4.05791,-1.68703) (4.08738,-1.65993) (4.11666,-1.63266) (4.14576,-1.60521) (4.17467,-1.57760) (4.20339,-1.54980) (4.23192,-1.52183) (4.26026,-1.49369) (4.28841,-1.46536) (4.31636,-1.43686) (4.34412,-1.40817) (4.37168,-1.37930) (4.39904,-1.35024) (4.42621,-1.32100) (4.45317,-1.29157) (4.47993,-1.26195) (4.50649,-1.23214) (4.53284,-1.20213) (4.55898,-1.17193) (4.58492,-1.14154) (4.61064,-1.11095) (4.63615,-1.08015) (4.66145,-1.04916) (4.68653,-1.01796) (4.71139,-0.98656) (4.73603,-0.95495) (4.76045,-0.92314) (4.78464,-0.89111) (4.80860,-0.85887) (4.80989,-0.85712)};
\draw[coursemuted,line width=.8pt,dashed] plot coordinates {(-22.10919,-0.18650) (-21.95303,-0.27377) (-21.78795,-0.36457) (-21.62301,-0.45388) (-21.45824,-0.54171) (-21.29368,-0.62809) (-21.12936,-0.71302) (-20.96533,-0.79652) (-20.80161,-0.87861) (-20.63823,-0.95930) (-20.47523,-1.03860) (-20.31264,-1.11654) (-20.15048,-1.19313) (-19.98877,-1.26838) (-19.82755,-1.34231) (-19.66684,-1.41494) (-19.50665,-1.48628) (-19.34702,-1.55635) (-19.18796,-1.62517) (-19.02949,-1.69275) (-18.87163,-1.75911) (-18.71439,-1.82426) (-18.55780,-1.88822) (-18.40187,-1.95101) (-18.24661,-2.01264) (-18.09204,-2.07313) (-17.93817,-2.13249) (-17.78502,-2.19074) (-17.63258,-2.24790) (-17.48088,-2.30398) (-17.32992,-2.35899) (-17.17972,-2.41296) (-17.03028,-2.46590) (-16.88161,-2.51781) (-16.73371,-2.56873) (-16.58660,-2.61865) (-16.44027,-2.66760) (-16.29474,-2.71560) (-16.15001,-2.76265) (-16.00609,-2.80876) (-15.86297,-2.85396) (-15.72067,-2.89826) (-15.57917,-2.94167) (-15.43850,-2.98420) (-15.29865,-3.02588) (-15.15961,-3.06670) (-15.02140,-3.10668) (-14.88401,-3.14585) (-14.74745,-3.18420) (-14.61171,-3.22176) (-14.47680,-3.25853) (-14.34271,-3.29452) (-14.20944,-3.32976) (-14.07700,-3.36425) (-13.94538,-3.39800) (-13.81458,-3.43102) (-13.68460,-3.46333) (-13.55543,-3.49493) (-13.42708,-3.52585) (-13.29955,-3.55608) (-13.17282,-3.58564) (-13.04690,-3.61454) (-12.92179,-3.64279) (-12.79748,-3.67040) (-12.67397,-3.69738) (-12.55126,-3.72374) (-12.42933,-3.74949) (-12.30820,-3.77464) (-12.18785,-3.79920) (-12.06828,-3.82318) (-11.94949,-3.84658) (-11.83147,-3.86942) (-11.71423,-3.89171) (-11.59775,-3.91345) (-11.48202,-3.93465) (-11.36706,-3.95532) (-11.25285,-3.97548) (-11.13938,-3.99511) (-11.02666,-4.01425) (-10.91468,-4.03289) (-10.80343,-4.05104) (-10.69290,-4.06870) (-10.58311,-4.08589) (-10.47403,-4.10262) (-10.36567,-4.11889) (-10.25801,-4.13470) (-10.15106,-4.15007) (-10.04481,-4.16500) (-9.93926,-4.17949) (-9.83439,-4.19356) (-9.73021,-4.20722) (-9.62671,-4.22046) (-9.52389,-4.23329) (-9.42173,-4.24572) (-9.32024,-4.25776) (-9.21940,-4.26941) (-9.11923,-4.28068) (-9.01970,-4.29157) (-8.92081,-4.30209) (-8.82256,-4.31224) (-8.72495,-4.32204) (-8.62797,-4.33148) (-8.53161,-4.34057) (-8.43587,-4.34931) (-8.34074,-4.35772) (-8.24623,-4.36579) (-8.15231,-4.37353) (-8.05900,-4.38095) (-7.96628,-4.38804) (-7.87414,-4.39482) (-7.78260,-4.40129) (-7.69163,-4.40746) (-7.60124,-4.41332) (-7.51141,-4.41888) (-7.42215,-4.42415) (-7.33345,-4.42913) (-7.24531,-4.43382) (-7.15771,-4.43823) (-7.07067,-4.44237) (-6.98416,-4.44623) (-6.89819,-4.44982) (-6.81275,-4.45314) (-6.72784,-4.45620) (-6.64345,-4.45900) (-6.55958,-4.46155) (-6.47622,-4.46384) (-6.39337,-4.46588) (-6.31103,-4.46768) (-6.22918,-4.46923) (-6.14784,-4.47055) (-6.06698,-4.47163) (-5.98661,-4.47247) (-5.90672,-4.47309) (-5.82732,-4.47347) (-5.74838,-4.47364) (-5.66992,-4.47358) (-5.59192,-4.47330) (-5.51439,-4.47280) (-5.43731,-4.47210) (-5.36069,-4.47118) (-5.28452,-4.47005) (-5.20879,-4.46872) (-5.13351,-4.46718) (-5.05866,-4.46544) (-4.98425,-4.46350) (-4.91027,-4.46137) (-4.83671,-4.45904) (-4.76358,-4.45653) (-4.69087,-4.45382) (-4.61857,-4.45092) (-4.54669,-4.44784) (-4.47521,-4.44458) (-4.40414,-4.44113) (-4.33347,-4.43750) (-4.26319,-4.43370) (-4.19331,-4.42972) (-4.12383,-4.42557) (-4.05472,-4.42125) (-3.98601,-4.41675) (-3.91767,-4.41209) (-3.84971,-4.40726) (-3.78212,-4.40227) (-3.71490,-4.39711) (-3.64806,-4.39179) (-3.58157,-4.38631) (-3.51545,-4.38068) (-3.44968,-4.37488) (-3.38427,-4.36893) (-3.31921,-4.36283) (-3.25450,-4.35657) (-3.19014,-4.35016) (-3.12612,-4.34361) (-3.06244,-4.33690) (-2.99909,-4.33005) (-2.93608,-4.32305) (-2.87340,-4.31590) (-2.81105,-4.30861) (-2.74902,-4.30118) (-2.68732,-4.29361) (-2.62594,-4.28590) (-2.56487,-4.27805) (-2.50412,-4.27006) (-2.44369,-4.26193) (-2.38356,-4.25367) (-2.32374,-4.24528) (-2.26422,-4.23674) (-2.20500,-4.22808) (-2.14609,-4.21929) (-2.08747,-4.21036) (-2.02915,-4.20130) (-1.97111,-4.19212) (-1.91337,-4.18280) (-1.85592,-4.17336) (-1.79875,-4.16379) (-1.74186,-4.15409) (-1.68525,-4.14427) (-1.62893,-4.13432) (-1.57287,-4.12425) (-1.51709,-4.11406) (-1.46159,-4.10375) (-1.40635,-4.09331) (-1.35138,-4.08275) (-1.29668,-4.07207) (-1.24223,-4.06127) (-1.18805,-4.05035) (-1.13413,-4.03931) (-1.08047,-4.02815) (-1.02706,-4.01688) (-0.97391,-4.00548) (-0.92100,-3.99397) (-0.86835,-3.98235) (-0.81594,-3.97060) (-0.76378,-3.95874) (-0.71187,-3.94677) (-0.66019,-3.93468) (-0.60876,-3.92248) (-0.55756,-3.91016) (-0.50660,-3.89773) (-0.45588,-3.88518) (-0.40539,-3.87253) (-0.35513,-3.85975) (-0.30511,-3.84687) (-0.25531,-3.83387) (-0.20573,-3.82076) (-0.15639,-3.80754) (-0.10727,-3.79421) (-0.05837,-3.78077) (-0.00969,-3.76721) (0.03877,-3.75355) (0.08701,-3.73977) (0.13504,-3.72588) (0.18285,-3.71188) (0.23044,-3.69777) (0.27783,-3.68355) (0.32500,-3.66922) (0.37196,-3.65478) (0.41872,-3.64023) (0.46527,-3.62557) (0.51161,-3.61079) (0.55775,-3.59591) (0.60368,-3.58092) (0.64941,-3.56581) (0.69494,-3.55060) (0.74027,-3.53528) (0.78541,-3.51984) (0.83034,-3.50430) (0.87508,-3.48864) (0.91962,-3.47287) (0.96397,-3.45700) (1.00813,-3.44101) (1.05209,-3.42491) (1.09586,-3.40870) (1.13944,-3.39238) (1.18283,-3.37594) (1.22604,-3.35940) (1.26905,-3.34274) (1.31188,-3.32597) (1.35453,-3.30909) (1.39699,-3.29209) (1.43926,-3.27498) (1.48136,-3.25776) (1.52326,-3.24043) (1.56499,-3.22298)};
\draw[aux line] (-2.4718,-0.6757)--(-9.5745,-4.2270);
\node[inner sep=0pt,font=\fontsize{8.5}{10}\selectfont] at (-2.4718,-0.6757) {$\oplus$};
\fill[courseink] (-9.5745,-4.2270) circle[radius=.85];
\node[anchor=north east,inner sep=1mm] at (-9.5745,-4.2270) {$P$};
\draw[coursemuted,line width=.45pt,->] (-24.0000,-14.5000)--(-23.3916,-17.8459) node[anchor=east,inner sep=.6mm] {$x$};
\draw[coursemuted,line width=.45pt,->] (-24.0000,-14.5000)--(-20.8097,-19.8277) node[anchor=north west,inner sep=.6mm] {$y$};
\draw[coursemuted,line width=.45pt,->] (-24.0000,-14.5000)--(-17.7990,-11.4307) node[anchor=west,inner sep=.6mm] {$z$};
\end{tikzpicture}
\par\figurelegend
{\usebeamerfont{caption}\color{coursemuted}
Пространственная проекция; Земля в масштабе.\par
Общая точка манёвра --- $P$.\par}
\end{column}
\end{columns}
\end{frame}

% Слайд 40
\begin{frame}{Коррекция $a$, $e$, $\Omega$: векторы импульса}
Преобразование найденных элементов даёт $\mathbf r_{2A}=\mathbf r_{2B}=\mathbf r_1$.
Скорости в неподвижных осях:
\par\smallskip\centering
\renewcommand{\arraystretch}{1.05}
\begin{tabular}{@{}lrrr@{}}
\toprule
Скорость, м/с & $v_x$ & $v_y$ & $v_z$\\\midrule
$\mathbf v_1$ & $6508{,}7585$ & $938{,}8305$ & $1924{,}8879$\\
$\mathbf v_{2A}$ & $6821{,}7824$ & $654{,}2175$ & $552{,}8699$\\
$\mathbf v_{2B}$ & $1956{,}3901$ & $3196{,}0463$ & $5764{,}3913$\\\bottomrule
\end{tabular}\par
\begin{align*}
\Delta\mathbf v_A&=(313{,}0240,\;-284{,}6131,\;-1372{,}0179)^{\mathsf T}\text{ м/с},\\
\Delta\mathbf v_B&=(-4552{,}3684,\;2257{,}2157,\;3839{,}5035)^{\mathsf T}\text{ м/с}.
\end{align*}
\[
\boxed{\Delta v_A=1435{,}7652\text{ м/с}},\qquad
\boxed{\Delta v_B=6368{,}7414\text{ м/с}}.
\]
\justifying
При требовании $\Omega_2=3^\circ$ меньший импульс вырос: $1388{,}84\to1435{,}77$ м/с.
\end{frame}

% Слайд 41
\begin{frame}{Выбор места коррекции $a$, $e$, $\Omega$}
\centering\begin{tikzpicture}
\begin{axis}[course plot,
width=128mm,height=48mm,ymin=0,ymax=13000,ytick={0,2000,4000,6000,8000,10000,12000},
ylabel={$\Delta v$, м/с},
xlabel={$\nu_1$, град},
xmin=0,xmax=360,xtick={0,60,120,180,240,300,360},
unbounded coords=jump,filter discard warning=false,
legend style={at={(0.98,.98)},anchor=north east,legend columns=2},
]
\addplot[courseaccent,solid,line width=.85pt,no marks] table[x=nu,y=dvA] {images/opt-node.dat};
\addplot[coursemuted,densely dashed,line width=.85pt,no marks] table[x=nu,y=dvB] {images/opt-node.dat};
\legend{A,B}
\end{axis}
\end{tikzpicture}
\par
\smallskip\raggedright
Цель та же: $a_2=26\,554$ км, $e_2=0{,}7$, $\Omega_2=3^\circ$.\newline В каждой точке заново находим $i_2$, $u_2$ и оба варианта импульса.
\end{frame}

% Слайд 42
\begin{frame}{Источники}
\setbeamertemplate{bibliography item}{\insertbiblabel}
\begin{thebibliography}{9}
\bibitem{DiPrinzio}
DiPrinzio M. Methods of Orbital Maneuvering. – American Institute of Aeronautics and Astronautics, 2019.

\bibitem{Vallado}
Vallado D. A. Fundamentals of astrodynamics and applications. – Springer Science \& Business Media, 2001. – Т. 12.

\bibitem{Eliasberg}
Эльясберг П. Е. Введение в теорию полета искусственных спутников Земли. – 1965.
\end{thebibliography}
\end{frame}

\end{document}
